\documentclass[11pt,oneside,openany]{book}
\usepackage[T1]{fontenc}
\usepackage[utf8]{inputenc}
\usepackage[sc]{mathpazo}
\renewcommand{\sfdefault}{lmss}
\renewcommand{\ttdefault}{lmtt}
\usepackage{amsmath,amssymb,amsthm,mathtools}
\usepackage[paperwidth=7.5in,paperheight=10in,inner=.85in,outer=.85in,top=.85in,bottom=.85in,headheight=15pt,headsep=20pt]{geometry}
\usepackage{microtype}
\usepackage{xcolor,graphicx,booktabs,array,longtable}
\newcolumntype{L}[1]{>{\raggedright\arraybackslash}p{#1}}
\usepackage{tikz,pgfplots}
\usetikzlibrary{arrows.meta,calc,patterns,decorations.markings}
\pgfplotsset{compat=1.18}
\definecolor{ink}{HTML}{17334B}
\definecolor{teal}{HTML}{137C83}
\definecolor{copper}{HTML}{AE5932}
\definecolor{soft}{HTML}{EEF3F5}
\pgfplotsset{every axis/.append style={axis line style={black!65},tick style={black!60},tick label style={font=\small},label style={font=\small},legend style={font=\small,draw=none,fill=white},grid style={black!10},major grid style={black!10},line width=.6pt},every axis plot/.append style={line width=1.15pt}}
\usepackage{enumitem,needspace,etoolbox}
\setlist[enumerate]{topsep=5pt,itemsep=5pt,parsep=0pt}
\setlist[itemize]{topsep=4pt,itemsep=3pt,parsep=0pt}
\usepackage{titlesec,fancyhdr}
\titleformat{\chapter}[display]{\normalfont\color{ink}}{\large\sffamily\MakeUppercase{Chapter \thechapter}}{12pt}{\huge\bfseries\raggedright\hyphenpenalty=10000\exhyphenpenalty=10000}[\vspace{12pt}\titlerule]
\titlespacing*{\chapter}{0pt}{-10pt}{24pt}
\titleformat{\section}{\large\bfseries\color{ink}}{\thesection}{.65em}{}
\titleformat{\subsection}{\normalsize\bfseries\color{ink}}{\thesubsection}{.6em}{}
\titlespacing*{\section}{0pt}{18pt plus 3pt minus 2pt}{7pt}
\titlespacing*{\subsection}{0pt}{12pt plus 2pt minus 2pt}{5pt}
\pagestyle{fancy}
\fancyhf{}
\fancyhead[L]{\small\nouppercase{\leftmark}}
\fancyfoot[C]{\small\thepage}
\renewcommand{\headrulewidth}{.3pt}
\renewcommand{\chaptermark}[1]{\markboth{\thechapter\quad #1}{}}
\fancypagestyle{plain}{\fancyhf{}\fancyfoot[C]{\small\thepage}\renewcommand{\headrulewidth}{0pt}}
\usepackage[font=small,labelfont=bf,labelsep=period]{caption}
\usepackage{float,placeins}
\floatplacement{figure}{htbp}
\setcounter{topnumber}{3}
\setcounter{bottomnumber}{2}
\renewcommand{\topfraction}{.9}
\renewcommand{\textfraction}{.08}
\renewcommand{\floatpagefraction}{.8}
\usepackage{makeidx}
\makeindex
\usepackage[pdfencoding=auto,psdextra,colorlinks=true,linkcolor=ink,citecolor=teal,urlcolor=teal,bookmarksnumbered=true]{hyperref}
\usepackage{bookmark}
\hypersetup{pdftitle={Discrete Information Geometry},pdfsubject={Finite differences, probability, and information flows}}
\urlstyle{same}
\newtheorem{theorem}{Theorem}[chapter]
\newtheorem{proposition}[theorem]{Proposition}
\newtheorem{lemma}[theorem]{Lemma}
\newtheorem{corollary}[theorem]{Corollary}
\theoremstyle{definition}
\newtheorem{definition}[theorem]{Definition}
\newtheorem{example}{Example}[chapter]
\BeforeBeginEnvironment{example}{\par\addvspace{4pt}\noindent\begin{minipage}{\linewidth}}
\AfterEndEnvironment{example}{\end{minipage}\par\addvspace{4pt}}
\theoremstyle{remark}
\newtheorem*{remark}{Remark}
\newenvironment{solution}{\par\smallskip\noindent\textit{Solution.}\ }{\par\medskip}
\newenvironment{exercises}{\begin{enumerate}[label=\thechapter.\arabic*,ref=\thechapter.\arabic*,leftmargin=2.8em,itemsep=6pt]}{\end{enumerate}}
\newenvironment{answers}{\begin{description}[leftmargin=2.9em,labelwidth=2.5em,labelsep=.4em,itemsep=8pt,font=\normalfont\bfseries]}{\end{description}}
\newcommand{\R}{\mathbb{R}}
\newcommand{\E}{\mathbb{E}}
\newcommand{\Var}{\operatorname{Var}}
\newcommand{\Cov}{\operatorname{Cov}}
\newcommand{\KL}[2]{D\!\left(#1\,\middle\|\,#2\right)}
\newcommand{\grad}{\operatorname{grad}}
\newcommand{\arcosh}{\operatorname{arcosh}}
\newcommand{\sech}{\operatorname{sech}}
\newcommand{\softmax}{\operatorname{softmax}}
\newcommand{\diag}{\operatorname{diag}}
\newcommand{\dd}{\,\mathrm{d}}
\newcommand{\ip}[2]{\left\langle #1,#2\right\rangle}
\newcommand{\p}{\mathbf{p}}
\newcommand{\q}{\mathbf{q}}
\newcommand{\f}{\mathbf{f}}
\newcommand{\uvec}{\mathbf{u}}
\newcommand{\vvec}{\mathbf{v}}
\newcommand{\ones}{\mathbf{1}}
\newcommand{\simp}{\Delta^{n-1}}
\newcommand{\simpi}{\Delta^{n-1}_{\!\circ}}
\numberwithin{equation}{chapter}
\setlength{\parindent}{1.2em}
\setlength{\parskip}{0pt}
\linespread{1.07}
\setlength{\emergencystretch}{2em}
\clubpenalty=10000
\widowpenalty=10000
\displaywidowpenalty=10000
\raggedbottom
\setcounter{tocdepth}{1}
\makeatletter
\renewcommand*\l@section{\@dottedtocline{1}{1.5em}{3.2em}}
\makeatother
\setcounter{secnumdepth}{2}
\providecommand{\tightlist}{\setlength{\itemsep}{0pt}\setlength{\parskip}{0pt}}
\providecommand{\passthrough}[1]{#1}
\newcommand{\labheading}{\FloatBarrier\section*{Laboratory investigation}\addcontentsline{toc}{section}{Laboratory investigation}}

\begin{document}
\frontmatter
\begin{titlepage}\color{ink}\centering\vspace*{1.0in}{\small\sffamily PROBABILITY AND MATHEMATICAL MODELLING\par}\vspace{.45in}{\Huge\bfseries Discrete Information\\[.12in]Geometry\par}\vspace{.35in}{\Large Finite Differences, Probability,\\and Information Flows\par}\vfill\begin{tikzpicture}[scale=1.05]\draw[teal,line width=1pt](0,0)--(4,0)--(2,3.464)--cycle;\foreach \a in {0,...,6}{\foreach \b in {0,...,6}{\pgfmathtruncatemacro{\c}{6-\a-\b}\ifnum\c>-1\fill[teal] ({(\b+\c/2)*4/6},{\c*3.464/6}) circle(1.5pt);\fi}}\end{tikzpicture}\vfill{\large A course in ten chapters\par}\vspace{.3in}{\small\url{https://bobovski66.github.io/IntroDiscreteInfoGeom/}\par}\end{titlepage}
\chapter*{Preface}\addcontentsline{toc}{chapter}{Preface}
A record of daily arrivals, a population's changing composition, and a sequence of model predictions all describe quantities at separate steps. To understand them, we subtract neighboring values, accumulate changes, and ask what repeated updates will do. This book develops those elementary operations into a calculus of probability distributions and information.

The starting prerequisites are fractions, powers, basic algebra, and coordinate graphs. Logarithms, probability, matrices, and the trigonometry used for statistical distance are developed where they enter the argument. Continuous calculus is not assumed. A calculator with logarithm, exponential, square-root, and inverse-trigonometric functions is useful for checking the numerical examples.

The chapters follow one accumulating line of inquiry. Chapters 1 and 2 establish finite differences, summation, and difference equations. Growth and feedback provide settings in which the meaning of a step and the size of a correction can be examined exactly. Chapters 3 and 4 turn a probability vector into a geometric state, define directional changes on a token lattice, and connect currents to conservation. Finite sampling then explains the gap between a model's probabilities and an observed composition.

Chapters 5 and 6 ask what uncertainty and model discrepancy do during an update. Entropy grows from average surprise into a function with finite directional differences. Relative entropy compares a state with a reference and supplies a potential whose one-step change can establish contraction. The proofs use finite inequalities and explicit calculations.

Chapters 7 and 8 construct statistical distance through square-root coordinates on a sphere. Exact distances are distinguished from local quadratic estimates, and unrestricted shortest arcs are compared with permitted routes through a lattice. Chapter 9 uses these ideas to choose a learning update by balancing a score against the information cost of change. Chapter 10 brings the operators together on networks and examines which model distinctions remain visible after transport or observation.

Several distinctions recur because they change the mathematics. Discrete time does not require a finite state space. Expected proportions are different from sampled counts. A direction of movement does not determine its length. A quantity that is monotone under one rule need not be monotone under another. Each distinction is developed through examples before it is used in a larger argument.

The web readings and the textbook contain the same chapters. Each chapter begins with a setting and a question, develops its notation and calculations, and ends with exercises, selected solutions, and a laboratory investigation. The site supplies the accompanying demonstrations and check-ins. Read a chapter over several sittings: reproduce a worked calculation, explore its demonstration, and return to the argument with those numbers in view. The contents list on each web reading provides navigation through the longer developments.

Natural logarithms give information quantities in nats; base-two logarithms give bits. Probability vectors have nonnegative entries totaling one, and matrices act on them as columns. Zero-probability conventions are introduced before they are needed. Local Fisher formulas require positive starting probabilities, while exact spherical distance includes boundary states. Every recurrence states a finite update, and its conclusions concern that stated rule.

The companion course, \emph{Introduction to Information Geometry}, develops the continuous theory at \url{https://bobovski66.github.io/IntroInfoGeom/}. The present book supplies an algebraic entry to many of the same questions through differences, sums, finite probability, and exact geometric constructions.

\tableofcontents
\mainmatter

\chapter{Sequences, differences, and discrete integration}
Imagine recording the number of messages arriving at a help desk each day. The daily counts tell one story: when demand is high, when it is growing, and when it begins to slow. The running total tells another: how much work has accumulated. Passing between these two descriptions means subtracting consecutive totals or adding consecutive arrivals. These are the two basic operations of discrete calculus.

The word \textbf{discrete} tells us that we are working with separate steps: days, generations, observations, or rounds of an algorithm. We will begin with quantities measured at these steps. Later the quantities will be probabilities, and the questions will concern learning, uncertainty, and information. The same subtraction-and-addition relationship will run through the entire course.

\section{Reading a sequence as a record of change}\label{ch01-reading-a-sequence-as-a-record-of-change}

A \textbf{sequence} is an ordered list of values. We write it as \(a_0,a_1,a_2,\ldots\), where the subscript records the step. The symbol \(a_3\) means the value at step three; it does not mean \(a\) multiplied by three. Starting with subscript zero lets us distinguish an initial measurement from measurements made after one or more updates.

A formula such as \(a_n=n^2\) gives each value directly. Substituting \(n=0,1,2,3,4\) gives \(0,1,4,9,16\). Another way to specify a sequence is to give an initial value and an update rule. For example, \[a_0=10,\qquad a_{n+1}=a_n+3\] means: start at ten, then add three each step. The resulting list is \(10,13,16,19,\ldots\). The formula \(a_n=10+3n\) describes the same list.

Before interpreting a sequence, give its terms and its steps units. If \(a_n\) is the number of messages received by the end of day \(n\), then \(a_{n+1}-a_n\) counts messages arriving during the next day. If measurements are two hours apart, dividing the change by two gives a rate per hour. Throughout most of this course, one update is our time unit. A finite change and a rate per physical unit agree numerically only when that unit has been specified appropriately.

\section{The first difference measures one step}\label{ch01-the-first-difference-measures-one-step}

The \textbf{forward difference} of a sequence is \[\boxed{\Delta a_n=a_{n+1}-a_n.}\] Read the symbol \(\Delta\) as ``change in.'' The subscript \(n\) identifies the step that starts at \(n\) and ends at \(n+1\). A positive difference means an increase, a negative difference means a decrease, and a zero difference means no change during that step.

\begin{example}[Separating a total from its increments]

Suppose cumulative arrivals at a help desk are \(a_0=0\), \(a_1=2\), \(a_2=7\), \(a_3=15\), and \(a_4=26\). Subtract consecutive entries: \[\Delta a_0=2,\quad\Delta a_1=5,\quad\Delta a_2=8,\quad\Delta a_3=11.\] There are five totals and four intervals between them. The differences count the arrivals during those four intervals. The total rises throughout, and its increments are also rising.

\end{example}

This operation plays the role of a derivative. It asks how the value changes as the input advances one step. Everything needed to compute it is present in two neighboring entries. A finite table can therefore support a useful calculus even when no formula between the observations is available.

Differences obey simple algebraic rules. For sequences \(a_n,b_n\) and a fixed number \(c\), \[\Delta(a_n+b_n)=\Delta a_n+\Delta b_n,\qquad
\Delta(ca_n)=c\Delta a_n.\] For example, the first identity follows by expanding \((a_{n+1}+b_{n+1})-(a_n+b_n)\) and regrouping. Such identities let us analyze a complicated record by separating its parts.

\section{Differences themselves can change}\label{ch01-differences-themselves-can-change}

If \(a_n=n^2\), then \[\Delta a_n=(n+1)^2-n^2=(n^2+2n+1)-n^2=2n+1.\] The square sequence has first differences \(1,3,5,7,\ldots\). They increase by two each time. Applying the difference operation again gives the \textbf{second difference}: \[\Delta^2a_n=\Delta a_{n+1}-\Delta a_n
=a_{n+2}-2a_{n+1}+a_n.\] For the square sequence, every second difference is two.

\begin{longtable}[]{@{}llll@{}}
\toprule\noalign{}
Step \(n\) & \(a_n=n^2\) & \(\Delta a_n\) & \(\Delta^2a_n\) \\
\midrule\noalign{}
\endhead
\bottomrule\noalign{}
\endlastfoot
0 & 0 & 1 & 2 \\
1 & 1 & 3 & 2 \\
2 & 4 & 5 & 2 \\
3 & 9 & 7 & \\
4 & 16 & & \\
\end{longtable}

The blank entries cannot be obtained from this five-term table because they need values beyond its end. The formula supplies those additional values if we want them.

A positive second difference says the first differences are increasing. It does not, by itself, say the original values are increasing. A falling sequence can have negative increments that become less negative. For instance, the values \(9,4,1,0\) have differences \(-5,-3,-1\) and second differences \(2,2\). The quantity falls, but its fall slows.

This distinction will matter when we study entropy. The sign of its first difference will tell us whether uncertainty rises or falls; its second difference will tell us how that tendency changes across neighboring states.

\begin{example}[A triangular accumulation]

Suppose a workshop completes one item in its first interval, two in its second, three in its third, and so on. Let \(T_n\) be the number completed after \(n\) intervals. Then \[T_0=0,\qquad \Delta T_n=n+1.\] The initial values are \(0,1,3,6,10\). Pairing the terms of \(1+2+\cdots+n\) in forward and reverse order gives \(n\) pairs, each with sum \(n+1\). Therefore \[T_n=\frac{n(n+1)}2.\] Check the formula by subtracting neighboring values: \[\frac{(n+1)(n+2)}2-\frac{n(n+1)}2=n+1.\] The cumulative sequence is quadratic because its first differences increase linearly. Its second differences are one.

\end{example}

This gives a way to move in both directions between patterns. Constant first differences suggest an affine formula \(an+b\). Constant second differences suggest a quadratic formula. The suggestion becomes a proof after substitution verifies the formula and its starting value. A short table alone does not determine all future terms: many rules can agree on its entries and disagree later.

A table also makes clear what information a difference requires. Measurements at steps zero through four provide four changes. We cannot calculate the change starting at four until the fifth-step measurement arrives. A formula can supply information beyond the observed table, but that information comes from the formula's assumptions.

For measured data, subtraction can mix several effects. If yesterday's reading was too high and today's is too low, their difference exaggerates an apparent fall. The operation is exact arithmetic; its interpretation depends on measurement quality. A deterministic sequence lets us study the operation before adding measurement error.

\begin{figure}[htbp]\centering
\resizebox{.98\linewidth}{!}{\begin{tikzpicture}
\begin{axis}[width=6.5cm,height=5.3cm,xlabel={Step $n$},ylabel={$a_n=n^2$},xmin=-.2,xmax=5.2,ymin=0,ymax=27,xtick={0,...,5},grid=major]
\addplot[teal,only marks,mark=*,mark size=2.5pt] coordinates {(0,0)(1,1)(2,4)(3,9)(4,16)(5,25)};
\end{axis}
\begin{axis}[at={(6.5cm,0)},anchor=south west,width=6.5cm,height=5.3cm,xlabel={Starting step $n$},ylabel={$\Delta a_n=2n+1$},xmin=-.3,xmax=4.3,ymin=0,ymax=10,xtick={0,...,4},grid=major]
\addplot[copper,ycomb,mark=*,mark size=2.5pt] coordinates {(0,1)(1,3)(2,5)(3,7)(4,9)};
\end{axis}\end{tikzpicture}}
\caption{Six values of the square sequence determine five first differences. Each difference is indexed by the starting step of its interval.}\end{figure}

\section{Summation reconstructs accumulated change}\label{ch01-summation-reconstructs-accumulated-change}

Return to the arrivals \(2,5,8,11\). Starting from zero and adding them gives the cumulative values \(0,2,7,15,26\). Summation reverses the difference operation once an initial value is supplied.

The notation \[\sum_{k=0}^{n-1}b_k\] means ``add the terms \(b_k\) as \(k\) takes the integer values from zero through \(n-1\).'' For example, \(\sum_{k=0}^3 b_k=b_0+b_1+b_2+b_3\). The letter \(k\) is a running index; we use a different letter from the endpoint \(n\) to keep their roles distinct.

Define the running sum by \(S_0=0\) and \[S_n=\sum_{k=0}^{n-1}b_k.\] Then \(S_{n+1}\) contains all the terms in \(S_n\) plus one new term, \(b_n\). Thus \[\Delta S_n=S_{n+1}-S_n=b_n.\] A sequence whose difference is \(b_n\) is called an \textbf{antidifference} of \(b_n\). Adding a constant to an antidifference gives another one. Specifying the starting value removes that ambiguity.

\begin{example}[An antidifference found by inspection]

We already know that \(\Delta(n^2)=2n+1\). Therefore \(n^2\) is an antidifference of \(2n+1\). Starting at \(a_0=7\) instead of zero gives \(a_n=7+n^2\). The increments remain \(2n+1\) because the constant seven cancels in every subtraction.

\end{example}

\section{Why the discrete fundamental theorem works}\label{ch01-why-the-discrete-fundamental-theorem-works}

Write out three successive differences: \[\begin{aligned}
\Delta a_m+\Delta a_{m+1}+\Delta a_{m+2}
&=(a_{m+1}-a_m)+(a_{m+2}-a_{m+1})\\
&\quad +(a_{m+3}-a_{m+2})\\
&=a_{m+3}-a_m.
\end{aligned}\] Every intermediate value is added once and subtracted once. This cancellation is called \textbf{telescoping}. The same argument works for any number of steps and gives the discrete fundamental theorem of calculus: \[\boxed{\sum_{k=m}^{n-1}\Delta a_k=a_n-a_m.}\] There are \(n-m\) changes between steps \(m\) and \(n\). The last change starts at \(n-1\), which explains the upper summation limit.

This identity is useful even when the changes alternate in sign. A sequence \(5,8,6,10\) has differences \(3,-2,4\). Their sum is five, exactly the endpoint change \(10-5\). Adding the absolute changes would instead give nine, a measure of total movement. Net change and total movement answer different questions. We will encounter the same distinction for paths through probability space.

\begin{example}[A balance equation for a stockroom]

Let \(S_n\) be stock at the beginning of a day, \(R_n\) the deliveries received during that day, and \(D_n\) the items used. Assuming no other additions or losses, \[\Delta S_n=R_n-D_n.\] Take \(S_0=20\), deliveries \((5,0,8)\), and use \((3,6,4)\) over three days. Net changes are \((2,-6,4)\), giving stock levels \(20,22,16,20\). Summing the balance equations yields \[S_3-S_0=(5+0+8)-(3+6+4)=0.\] Final and initial stocks agree, although activity occurred throughout. Summing deliveries alone would describe throughput rather than the change in stored quantity.

\end{example}

This balance structure will reappear when probability moves through a network. A vertex's stored mass changes by arrivals minus departures. The form of that law is already visible in a stockroom calculation.

\section{Repeated multiplication gives a different kind of growth}\label{ch01-repeated-multiplication-gives-a-different-kind-of-growth}

A fixed addition gives constant first differences. A fixed multiplication gives a \textbf{geometric sequence}: \[a_{n+1}=r a_n,\qquad a_n=a_0r^n.\] To see the formula, compute \(a_1=ra_0\), then \(a_2=r(ra_0)=r^2a_0\), and continue. Here \(r\) is the multiplier per step. Growth of twenty percent corresponds to \(r=1.2\); a decrease of twenty percent corresponds to \(r=0.8\).

Its first difference is \[\Delta a_n=(r-1)a_n.\] The absolute increment depends on the current amount. Ten percent of a larger amount is a larger increment, even though the percentage rate is unchanged.

In a population model, a multiplier may produce noninteger counts. We then interpret counts as amounts or expected counts. A rule for a literal finite population would also need a sampling or rounding procedure. Chapter 4 develops the sampling interpretation.

There are two separate uses of the word discrete here. The time index may be discrete while the state takes arbitrary real values, as in \(a_{n+1}=1.2a_n\). Alternatively, the permitted states may form a finite list, as when a fixed number of tokens is allocated among categories. Both will occur in this course. A finite number of category labels does not force every probability to be an integer fraction with a fixed denominator.

\section{Composition changes differently from total size}\label{ch01-composition-changes-differently-from-total-size}

Suppose two types have counts \(A_n\) and \(B_n\). Their proportions are \[x_n=\frac{A_n}{A_n+B_n},\qquad 1-x_n=\frac{B_n}{A_n+B_n}.\] Dividing by the total is called \textbf{normalization}. It removes overall size and retains composition. Counts \((20,80)\) and \((200,800)\) both give \(x=1/5\).

Let each type multiply once per step: \[A_{n+1}=w_AA_n,\qquad B_{n+1}=w_BB_n,\qquad w_A,w_B>0.\] The new total is \(w_AA_n+w_BB_n\). Substituting the old proportions into the ratio gives \[\boxed{x_{n+1}=\frac{w_Ax_n}{w_Ax_n+w_B(1-x_n)}.}\] The denominator is the average multiplier in the old population. It is required because both the numerator and the total can change.

\begin{example}[Both counts rise while one share falls]

Start with counts \((50,50)\) and multipliers \((1.1,1.3)\). The new counts are \((55,65)\), so the \(A\) share becomes \(55/120=11/24\). Although \(A\) increased from fifty to fifty-five, its share fell from \(1/2\) to about \(0.4583\). The other type grew faster.

\end{example}

To measure composition change, subtract \(x_n\) from the update. Write \(Z=w_Ax+w_B(1-x)\) for the denominator at the current step. Then \[\begin{aligned}
x^+-x&=\frac{w_Ax-xZ}{Z}
=\frac{x[w_A-w_Ax-w_B(1-x)]}{Z}\\
&=\frac{x(1-x)(w_A-w_B)}{Z}.
\end{aligned}\] Here \(x^+\) is a shorter name for the next value. Consequently, \[\boxed{\Delta x_n=\frac{x_n(1-x_n)(w_A-w_B)}{w_Ax_n+w_B(1-x_n)}.}\] For \(0<x_n<1\), every factor other than \(w_A-w_B\) is positive. The difference therefore has the sign of the relative advantage. This is our first example of extracting a qualitative conclusion from an exact discrete derivative.

At \(x=0\) or \(x=1\), the numerator is zero. A missing type remains missing under multiplication. Positive multipliers also keep every initially present type positive at every finite step. These conclusions concern the stated update; adding migration would give another model.

\begin{example}[A finite calculus calculation on proportions]

With \(x_0=1/4\), \(w_A=2\), and \(w_B=1\), \[x_1=\frac{2(1/4)}{2(1/4)+3/4}=\frac25,\qquad
\Delta x_0=\frac25-\frac14=\frac3{20}.\] The next state is \(x_2=4/7\), so \(\Delta x_1=4/7-2/5=6/35\). Adding the two differences gives \(9/28\), which equals \(x_2-x_0\). Telescoping works for this nonlinear update just as it did for a table of message totals.

\end{example}

\begin{figure}[htbp]\centering
\resizebox{.98\linewidth}{!}{
\begin{tikzpicture}
\begin{axis}[width=6.7cm,height=5.3cm,xlabel={Step $t$},ylabel={Count},xmin=0,xmax=6,ymin=0,legend pos=north west]
\addplot[teal,mark=*] coordinates {(0,20) (1,40) (2,80) (3,160) (4,320) (5,640) (6,1280)};\addlegendentry{$A$}
\addplot[copper,mark=square*] coordinates {(0,80) (1,80) (2,80) (3,80) (4,80) (5,80) (6,80)};\addlegendentry{$B$}
\end{axis}
\begin{axis}[at={(6.5cm,0)},anchor=south west,width=6.7cm,height=5.3cm,xlabel={Step $t$},ylabel={$A$ proportion},xmin=0,xmax=6,ymin=0,ymax=1]
\addplot[teal,mark=*] coordinates {(0,0.2) (1,0.333333) (2,0.5) (3,0.666667) (4,0.8) (5,0.888889) (6,0.941176)};
\end{axis}\end{tikzpicture}}
\caption{Counts and proportions for initial counts (20,80) and multipliers (2,1). The total population changes at every step.}\end{figure}

\section{A change of coordinates can simplify a recurrence}\label{ch01-a-change-of-coordinates-can-simplify-a-recurrence}

For \(0<x<1\), define the \textbf{odds} \(O=x/(1-x)\). Odds of three mean three units of type \(A\) for each unit of type \(B\), so the proportion is \(3/4\). Solving the definition gives \(x=O/(1+O)\).

In the normalized growth model, the common denominator cancels when we divide the two new proportions: \[O_{n+1}=\frac{w_A}{w_B}O_n.\] The odds follow a geometric sequence. With \(r=w_A/w_B\), \[O_n=O_0r^n,\qquad
x_n=\frac{x_0r^n}{1-x_0+x_0r^n}.\] The final formula also makes sense for \(x_0=0\) and \(x_0=1\), even though the intermediate odds calculation assumed an interior starting value.

If \(x_0=1/5\) and \(r=2\), then \(O_n=2^n/4\). Reaching a proportion of at least \(0.9\) means reaching odds of at least nine, so \(2^n\geq36\). Since \(2^5=32\) and \(2^6=64\), the first qualifying step is six. A question about a nonlinear fraction has become a comparison of integer powers.

This use of coordinates anticipates information geometry. Different numerical descriptions of the same state can make different features easy to see. Proportions display composition; odds simplify relative multiplication.

We can also run a positive growth step backward. New odds satisfy \(O^+=rO\), so old odds are \(O=O^+/r\). Converting back gives \[x=\frac{x^+}{r(1-x^+)+x^+}.\] Substituting a forward update recovers the starting proportion. This checks a simulation and explains why exact positive multiplication cannot delete a present type in one step. Rounding can destroy that reversibility: a small positive proportion rounded to zero loses the information needed to reconstruct its earlier value.

\section{Turning a record into a model}\label{ch01-turning-a-record-into-a-model}

A sequence of observations and an update law are different starting points. With observations, compute differences and look for patterns. With a law, generate observations and prove identities. In applications we move between these tasks, and it helps to record which statements come from data and which come from assumptions.

A count table that doubles for three days supports the description ``doubling occurred during these intervals.'' Predicting doubling for a hundred days assumes a continuing multiplier of two. Resource limits or environmental changes might alter that rule. Similarly, a constant count ratio in a short record is evidence about composition, while the normalization calculation proves scale independence under the stated multiplier model.

A useful calculation record includes the initial state, the meaning of one step, the rule, and a check suited to the quantity. For cumulative totals, increments should add to the endpoint change. For probabilities, check nonnegativity and total one. For an explicit solution, check both its initial value and recurrence. These checks connect an answer to the question it is meant to resolve.

\section{From an interval to probability vectors}\label{ch01-from-an-interval-to-probability-vectors}

For three or more categories, collect the proportions into a vector \(p=(p_1,\ldots,p_k)\) with \[p_i\geq0,\qquad \sum_{i=1}^k p_i=1.\] The index \(i\) labels a category. Such a list is a \textbf{probability vector} because selecting one individual uniformly gives category \(i\) with probability \(p_i\). The same list can describe uncertain outcomes even without a population interpretation.

Positive multipliers give the general normalized update \[p_i^+=\frac{p_iw_i}{Z},\qquad Z=\sum_{j=1}^k p_jw_j.\] The new entries are nonnegative. Adding them gives \(Z/Z=1\), so they are still a probability vector. Multiplying every multiplier by the same positive constant changes both numerator and denominator by that factor and leaves the answer unchanged.

For \(p=(1/2,1/3,1/6)\) and \(w=(1,2,3)\), the unnormalized entries are \((1/2,2/3,1/2)\). Their sum is \(5/3\), giving \(p^+=(3/10,2/5,3/10)\). The changes add to zero because one unit of total probability has been redistributed. This conservation will become a vector calculus identity later in the course.

In the \href{https://bobovski66.github.io/IntroDiscreteInfoGeom/demos/operator-explorer.html}{operator explorer}, use the differences-and-sums view to compare a sequence with its difference table. Then use the \href{https://bobovski66.github.io/IntroDiscreteInfoGeom/demos/week01.html}{growth demonstration} to record a proportion trajectory and check a telescoping sum. Chapter 2 asks a new question about such trajectories: when repeated corrections move a state toward a target, what prevents them from overshooting and growing unstable?

\section{Exercises}\label{ch01-exercises}

\begin{exercises}

\item 

A sequence records cumulative arrivals \(0,3,8,15,24\). Find all first and second differences available from the table. Interpret both in the context of arrivals.

\item 

For \(a_n=3n^2+2n+5\), calculate \(\Delta a_n\) and \(\Delta^2a_n\) by expanding the expressions.

\item 

Starting from \(a_0=7\), reconstruct \(a_1,a_2,a_3,a_4\) when \(\Delta a_n=2n+3\). Check the endpoint change by summation.

\item 

Find an antidifference of \(2n+1\). Explain why specifying its value at zero determines it uniquely.

\item 

A sequence is \(10,7,5,4\). Calculate its differences and second differences. Explain how it can be decreasing while its second differences are positive.

\item 

Starting from \(a_0=12\), list five terms, including the initial value, of \(a_{n+1}=(3/2)a_n\). Give the first difference as a multiple of \(a_n\).

\item 

Normalize \((6,9,15)\) and \((14,21,35)\). Explain why the results agree.

\item 

Starting with counts \((30,70)\) and multipliers \((3,2)\), calculate counts and proportions after two steps. Verify the sum of the two proportion changes.

\item 

Derive the displayed formula for \(x^+-x\) without skipping the common-denominator step. Use it to explain why equal multipliers preserve composition.

\item 

Let \(x_0=1/10\) and \(w_A/w_B=3\). Find the first step at which \(x_n\geq1/2\) by using odds.

\item 

Compute one normalized update for \(p=(1/4,1/4,1/2)\) and \(w=(4,2,1)\). Prove that the answer is unchanged for \(w=(20,10,5)\).

\item 

For constant positive multipliers, prove by induction that \(p_{i,n}=p_{i,0}w_i^n/\sum_jp_{j,0}w_j^n\). State what changes if the multipliers depend on the current proportions.

\end{exercises}

\section{Selected solutions}\label{ch01-selected-solutions}

\textbf{1.} First differences are \(3,5,7,9\) and second differences are \(2,2,2\). Arrivals increase by two from one recorded interval to the next.

\textbf{2.} Expanding \(3(n+1)^2+2(n+1)+5-(3n^2+2n+5)\) gives \(6n+5\). Taking its difference gives six.

\textbf{3.} The successive increments are \(3,5,7,9\), giving \(7,10,15,22,31\). Their sum is \(24=31-7\).

\textbf{8.} Counts become \((90,140)\) and \((270,280)\). The proportions are \(x_0=3/10\), \(x_1=9/23\), \(x_2=27/55\). The sum of the changes is \(27/55-3/10=21/110\).

\textbf{10.} Initial odds are \(1/9\), so \(O_n=3^n/9\). A proportion of at least one half requires odds of at least one. The first such step is \(n=2\).

\textbf{11.} The unnormalized entries are \((1,1/2,1/2)\) and their total is two. The new table is \((1/2,1/4,1/4)\). Scaling all multipliers by five scales this total and each numerator by five, which cancels.
\labheading\begin{enumerate}
\def\labelenumi{\arabic{enumi}.}
\tightlist
\item
  Start with the sequence \(a_n=n^2\) for \(n=0,1,2,3,4,5\). (a) Calculate the six values. (b) List every first difference \(\Delta a_n\). (c) Calculate the second differences. (d) Explain what constant second differences tell you.
\item
  Suppose the daily number of new messages is \(b_n=3n+2\), beginning at day \(n=0\). (a) Calculate \(b_0,b_1,b_2,b_3\). (b) Compute the running totals \(S_0,S_1,S_2,S_3,S_4\), where \(S_0=0\). (c) Verify \(\Delta S_2=b_2\).
\item
  Let \(A_0=20\), \(B_0=60\), \(w_A=2\), and \(w_B=1\). (a) Compute counts at steps 0, 1, and 2. (b) Calculate the A proportions \(x_0,x_1,x_2\). (c) Find \(\Delta x_0\) and \(\Delta x_1\). (d) Verify the telescoping identity \(\Delta x_0+\Delta x_1=x_2-x_0\).
\item
  In the original growth demonstration, start at proportion \(x_0=0.2\) with relative multiplier \(2\). (a) Record five consecutive proportions. (b) Subtract to construct their difference sequence. (c) Add those differences and compare with the net change. (d) Explain why the differences get smaller near the limiting state.
\item
  Choose one sequence of your own whose first differences are constant and another whose second differences are constant. (a) Write five terms of each. (b) Show the appropriate difference tables. (c) Explain how the original sequence can be recovered by summing changes.
\end{enumerate}

\FloatBarrier

\chapter{Difference equations, feedback, and stability}
A thermostat, a stockroom reorder rule, and a learning algorithm share a basic structure. Each observes a current state, compares it with a desired state, and makes a correction. Repeating a sensible correction can produce smooth adjustment, repeated overshooting, or growing instability. The size of the correction matters as much as its direction.

Chapter 1 gave us a language for changes. In this chapter we use it to prescribe changes and study what repeated updates do. We will first solve a feedback model completely. Then we will return to population proportions, where the advantage of a type can depend on how common it is.

\section{An update rule generates a trajectory}\label{ch02-an-update-rule-generates-a-trajectory}

A \textbf{difference equation}, also called a \textbf{recurrence}, specifies how a sequence is generated. An equation of the form \[x_{n+1}=F(x_n)\] says: apply the function \(F\) to the current state to obtain the next one. Given \(x_0\), compute \(x_1=F(x_0)\), then \(x_2=F(x_1)\), and continue. The resulting sequence is a \textbf{trajectory} or \textbf{orbit}.

We can express the same rule using the first difference: \[\Delta x_n=F(x_n)-x_n.\] The first form emphasizes the next state; the second emphasizes the change. When designing a controller, the change form is often easier to interpret.

A model also needs a set of permitted states. A temperature can take many real values. A probability must stay between zero and one. A formula can make algebraic sense while sending a probability outside its permitted interval. Checking the state space is therefore part of understanding the recurrence.

\begin{example}[A small inventory correction]

Suppose \(x_n\) is the stock level at the start of a weekly restocking cycle, measured in crates. The target is ten crates. A simplified rule replaces half the gap: \[\Delta x_n=\frac12(10-x_n).\] Starting from two gives changes of four, two, and one, hence stock levels \(2,6,8,9\). These values can be interpreted as average stock levels if fractions occur. The rule describes a correction step; a full inventory model would also specify incoming orders and customer demand.

\end{example}

\section{Subtracting the target reveals the structure}\label{ch02-subtracting-the-target-reveals-the-structure}

Write the general correction rule as \[\Delta x_n=-k(x_n-c),\] where \(c\) is a target and \(k\) is the correction strength. If \(x_n<c\), the gap \(x_n-c\) is negative, so the prescribed change is positive when \(k>0\). If \(x_n>c\), the prescribed change is negative.

Expanding the difference gives \[x_{n+1}=x_n-k(x_n-c)=c+(1-k)(x_n-c).\] The state \(c\) is a \textbf{fixed point}: starting there gives zero change and leaves the state there forever. In general, a fixed point of \(F\) satisfies \(F(c)=c\).

Let \(e_n=x_n-c\) be the error. Subtract \(c\) from both sides of the update: \[e_{n+1}=(1-k)e_n.\] This is the geometric recurrence from Chapter 1. Therefore \[e_n=(1-k)^n e_0,\qquad
\boxed{x_n=c+(1-k)^n(x_0-c).}\] Introducing the error has turned a recurrence with a constant term into repeated multiplication. It also gives a direct answer to the question of long-term behavior.

\begin{example}[A replenishment rule with an implicit target]

Suppose a stockroom retains eighty percent of its stock after a day's use and then receives six replacement units. Its simplified recurrence is \[x_{n+1}=0.8x_n+6.\] The target is not stated in advance. Find a fixed point by asking what stock level reproduces itself: \[c=0.8c+6,\qquad 0.2c=6,\qquad c=30.\] Subtracting thirty gives \(x_{n+1}-30=0.8(x_n-30)\). Starting from ten, the first values are \(10,14,17.2,19.76\). The formula is \[x_n=30+0.8^n(10-30).\] A balance between removal and replenishment creates the equilibrium. At thirty units, the six units used match the six received.

\end{example}

More generally, for \(x^+=ax+b\) with \(a\ne1\), the fixed point is \(c=b/(1-a)\) and the error multiplier is \(a\). At \(a=1\), the rule is repeated addition: if \(b=0\), every state is fixed; if \(b\ne0\), there is no fixed state. Solving the fixed-point equation before dividing prevents this exceptional case from being lost.

\section{How corrections can settle or overshoot}\label{ch02-how-corrections-can-settle-or-overshoot}

The multiplier \(1-k\) controls both the sign and the size of the error. A negative multiplier crosses the target. A multiplier whose absolute value is less than one shrinks the error.

\begin{longtable}[]{@{}
  >{\raggedright\arraybackslash}p{(\linewidth - 4\tabcolsep) * \real{0.3333}}
  >{\raggedright\arraybackslash}p{(\linewidth - 4\tabcolsep) * \real{0.3333}}
  >{\raggedright\arraybackslash}p{(\linewidth - 4\tabcolsep) * \real{0.3333}}@{}}
\toprule\noalign{}
\begin{minipage}[b]{\linewidth}\raggedright
Correction strength
\end{minipage} & \begin{minipage}[b]{\linewidth}\raggedright
Error multiplier
\end{minipage} & \begin{minipage}[b]{\linewidth}\raggedright
Behavior from a nonzero error
\end{minipage} \\
\midrule\noalign{}
\endhead
\bottomrule\noalign{}
\endlastfoot
\(0<k<1\) & \(0<1-k<1\) & Approaches the target from one side \\
\(k=1\) & \(0\) & Reaches the target in one step \\
\(1<k<2\) & \(-1<1-k<0\) & Crosses the target with shrinking errors \\
\(k=2\) & \(-1\) & Alternates with unchanged error size \\
\(k>2\) & \(1-k<-1\) & Crosses the target with growing errors \\
\end{longtable}

At \(k=0\), nothing changes. For \(k<0\), a nonzero error grows because the correction points away from the target. For this model, attraction occurs exactly when \[\boxed{|1-k|<1,\quad\text{or equivalently }0<k<2.}\]

\begin{example}[Three correction strengths at the same target]

Take \(c=1/2\) and \(x_0=1/4\). With \(k=1/2\), the states begin \[\frac14,\quad\frac38,\quad\frac7{16},\quad\frac{15}{32}.\] With \(k=3/2\), they begin \[\frac14,\quad\frac58,\quad\frac7{16},\quad\frac{17}{32}.\] Both errors halve in magnitude each step; the second trajectory crosses the target. With \(k=5/2\), the states begin \[\frac14,\quad\frac78,\quad-\frac1{16}.\] The errors grow by a factor of \(3/2\). The last value also shows that this rule does not preserve the probability interval for these parameters.

\end{example}

A fixed point is \textbf{attracting} if trajectories starting sufficiently near it approach it. It is \textbf{stable} if sufficiently small initial errors remain small at all later steps. The case \(k=2\) illustrates the distinction: small errors remain small but never shrink to zero. For \(0<k<2\), the fixed point is both stable and attracting. We will normally describe this behavior as settling toward the target.

For \(0\leq k\leq1\) and \(c\in[0,1]\), the formula \(x^+=(1-k)x+kc\) is a weighted average of two points in the interval, so it stays in the interval. Attraction for \(1<k<2\) does not provide that same global guarantee. Stability and preservation of the permitted state space are separate checks.

The explicit formula also tells us how quickly a desired accuracy is reached. For the inventory rule, the starting error has magnitude twenty. To reduce it below one unit, require \[20(0.8)^n<1.\] Integer powers give \(20(0.8)^{13}\approx1.10\) and \(20(0.8)^{14}\approx0.88\). Thus fourteen steps suffice and thirteen do not. The trajectory never equals thirty at a finite step from this starting state, but it eventually lies within any positive tolerance.

A useful robustness question compares two initial states under the same rule. If \(x_n\) and \(y_n\) obey affine feedback with the same target and strength, then \[x_{n+1}-y_{n+1}=(1-k)(x_n-y_n).\] Initial disagreement shrinks at the same rate as the target error when \(0<k<2\). This compares whole trajectories, rather than comparing a single trajectory with a fixed point. It is an early example of contraction: a rule can systematically reduce the difference between states.

The word eventually concerns the tail of an infinite sequence. We can still make it concrete by giving a tolerance and a finite step after which the bound holds. In later chapters, information inequalities will provide similar stopping estimates for probability models.

\begin{figure}[htbp]\centering
\resizebox{.98\linewidth}{!}{\begin{tikzpicture}
\begin{axis}[width=6.5cm,height=5.3cm,xlabel={Step $n$},ylabel={Error $x_n-c$},xmin=0,xmax=5,ymin=-.3,ymax=.2,xtick={0,...,5},legend pos=north east,grid=major]
\addplot[teal,mark=*] coordinates {(0,-0.25) (1,-0.125) (2,-0.0625) (3,-0.03125) (4,-0.015625) (5,-0.0078125)};\addlegendentry{$k=1/2$}
\addplot[copper,mark=square*] coordinates {(0,-0.25) (1,0.125) (2,-0.0625) (3,0.03125) (4,-0.015625) (5,0.0078125)};\addlegendentry{$k=3/2$}
\end{axis}
\begin{axis}[at={(6.5cm,0)},anchor=south west,width=6.5cm,height=5.3cm,xlabel={Step $n$},ylabel={Error $x_n-c$},xmin=0,xmax=5,ymin=-1.5,ymax=2.1,xtick={0,...,5},legend pos=north west,grid=major]
\addplot[copper,mark=square*] coordinates {(0,-0.25) (1,0.375) (2,-0.5625) (3,0.84375) (4,-1.26562) (5,1.89844)};\addlegendentry{$k=5/2$}
\end{axis}\end{tikzpicture}}
\caption{Feedback errors from initial error -1/4. On the left, both magnitudes shrink and one alternates sign. On the right, an excessive correction amplifies the error; note the different vertical scale.}\end{figure}

\section{Fixed points are found before trajectories are solved}\label{ch02-fixed-points-are-found-before-trajectories-are-solved}

For a more complicated function \(F\), an explicit formula for every \(x_n\) may be unavailable. We can still look for fixed points by solving \(F(x)=x\), and then examine the increment \(F(x)-x\).

If the increment is positive below a candidate target and negative above it, single steps point toward the target. This is useful information, but the overshooting example shows why direction alone is incomplete. A large step can jump across the target and leave a larger error.

A practical investigation therefore records three items: the current state, the increment, and the new distance from the candidate fixed point. Numerical evidence can suggest attraction or instability. A proof needs an identity or inequality that covers every relevant step.

\section{Interactions make advantages depend on composition}\label{ch02-interactions-make-advantages-depend-on-composition}

In Chapter 1, one type always had the same relative multiplier. Many applications involve feedback: a strategy performs well when rare but poorly when common. A shared resource might support several methods of use, with the success of each depending on the mixture.

A \textbf{payoff matrix} records scores for pairwise encounters. For two types, write \[M=\begin{pmatrix}a&b\\c&d\end{pmatrix}.\] The row gives the focal type; the column gives the encountered type. Thus \(b\) is the score of an \(A\) individual encountering \(B\). The entries are scores, not probabilities, so they need not add to one.

Assume encounters are sampled according to the population proportions \(x\) and \(1-x\). The expected scores are weighted averages: \[f_A(x)=ax+b(1-x),\qquad f_B(x)=cx+d(1-x).\] For instance, if one quarter of encounters are with \(A\), the \(A\) score is one quarter of \(a\) plus three quarters of \(b\).

Subtract the two scores and call the result \(d(x)\): \[d(x)=f_A(x)-f_B(x)=(b-d)+(a-b-c+d)x.\] We use \(d(x)\) for a score gap; the symbol \(\Delta x\) remains reserved for a change between time steps.

\begin{example}[An advantage that reverses]

For \[M=\begin{pmatrix}2&4\\3&1\end{pmatrix},\] the scores are \(f_A=4-2x\) and \(f_B=1+2x\). Thus \(d(x)=3-4x\). At \(x=1/4\), the scores are \(7/2\) and \(3/2\), favoring \(A\). At \(x=9/10\), they are \(11/5\) and \(14/5\), favoring \(B\). They agree at \(q=3/4\), where both equal \(5/2\). This table creates feedback toward coexistence.

\end{example}

The encounter assumption matters. If individuals preferentially meet their own type, the weights need to reflect that preference. A matrix by itself does not specify how encounters are selected.

The table can be understood as a stylized model of two ways to use a shared resource. A method's score depends on which method it encounters. If a rare method receives a large score against the common one, rarity can be advantageous. If each method performs best among its own type, commonness can reinforce itself.

The numbers in a payoff table need an interpretation before a population update follows. They might represent a measured benefit, a relative reproductive score, or a reward assigned by an algorithm. The table specifies the result of an encounter; the averaging assumption specifies the mix of encounters; the conversion to multipliers specifies how scores affect the next population. Keeping these modeling steps visible makes it possible to change one assumption without confusing it with the others.

At a composition \(x=1/4\), for instance, the coexistence table predicts three times as many encounters with \(B\) as with \(A\). The average score \(f_A=(1/4)2+(3/4)4=7/2\) is between two and four. This interval check is a useful way to catch reversed weights or a mistaken row convention.

\section{Turning a score into a multiplier}\label{ch02-turning-a-score-into-a-multiplier}

A population multiplier must be nonnegative, and we will use strictly positive multipliers. Scores may be negative. We therefore state an explicit conversion.

The \textbf{linear multiplier rule} is \[w_i=1+\eta f_i,\] where \(\eta>0\) controls the intensity. It requires \(1+\eta f_i>0\) for every score that can occur. If every possible score is at least \(-2\), for example, choosing \(0<\eta<1/2\) ensures positivity.

The \textbf{exponential multiplier rule} is \[w_i=e^{\eta f_i}.\] Here \(e\approx2.71828\). An exponential satisfies \(e^{u+v}=e^u e^v\) and is positive for every real exponent. Negative scores give multipliers between zero and one. A calculator's exponential key evaluates them; the identities and positivity are the properties we need.

Normalize either set of multipliers as in Chapter 1. For the linear choice, \[x^+=\frac{x(1+\eta f_A)}{1+\eta\bar f},\qquad
\bar f=xf_A+(1-x)f_B.\] Subtracting \(x\) gives \[\Delta x=\frac{\eta x(1-x)d(x)}{1+\eta\bar f}.\] For the exponential choice, \[x^+=\frac{xe^{\eta f_A}}{xe^{\eta f_A}+(1-x)e^{\eta f_B}}.\] Its increment has the sign of \(e^{\eta f_A}-e^{\eta f_B}\), hence the sign of \(d(x)\). Both rules preserve the interval and have the same interior fixed points. Their step lengths can be very different.

\begin{example}[One score gap, two sizes of correction]

At \(x=1/4\) for the preceding table, take \(\eta=1\). Linear multipliers are \((9/2,5/2)\), giving \[x^+=\frac{(1/4)(9/2)}{(1/4)(9/2)+(3/4)(5/2)}=\frac38.\] For exponential multipliers, cancel the common factor \(e^{3/2}\): \[x^+=\frac{e^2}{e^2+3}\approx0.7112.\] Both changes are positive. The second is much larger.

\end{example}

\section{A convergence proof built from the error}\label{ch02-a-convergence-proof-built-from-the-error}

For the coexistence table, the linear rule approaches \(q=3/4\) without crossing it, for every \(\eta>0\). We can prove this by repeating the error calculation used for simple feedback.

First, \[\bar f=x(4-2x)+(1-x)(1+2x)=1+5x-4x^2.\] Since \(d(x)=-4(x-q)\), subtracting \(q\) yields \[x^+-q=(x-q)\left(1-\frac{4\eta x(1-x)}{1+\eta\bar f}\right).\] For \(0<x<1\), the subtracted fraction is positive. It is also less than one, because \[1+\eta\bar f-4\eta x(1-x)=1+\eta(1+x)>0.\] Thus the parenthesized factor lies strictly between zero and one. The error keeps its sign and shrinks.

A trajectory starting below \(q\) rises but cannot cross \(q\). It is increasing and bounded above, so its values approach a limit. This elementary fact means a quantity cannot keep making positive progress of a fixed minimum size while remaining below a fixed ceiling. The same reasoning applies to a decreasing sequence bounded below.

The rational update is continuous on the interval because its denominator is positive. Substituting the limit into the recurrence shows that it must be a fixed point. The only fixed point between the starting value and \(q\) is \(q\) itself. This completes the convergence argument.

The proof does more than a sign diagram: it controls crossing, establishes a limit, and identifies that limit. These are distinct steps in a stability argument.

The error identity suggests a general method. Suppose a recurrence has a fixed point \(q\) and can be written \[x^+-q=(x-q)R(x).\] If \(0<R(x)<1\) throughout a relevant interval, each nonfixed state moves toward \(q\) without crossing it. If a uniform bound \(R(x)\leq\rho<1\) is available, then repeated errors satisfy \[|x_n-q|\leq\rho^n|x_0-q|.\] This gives a rate as well as a direction. If \(R\) is negative, crossings occur; a bound \(|R(x)|<1\) controls whether those crossings shrink.

A statement about the sign of the increment answers a slightly different question. Below \(q\), \(\Delta x>0\) says the next point lies to the right of the current one. It does not specify whether it lies before or beyond \(q\). Drawing the current point, the target, and the computed next point on a line is often the clearest first check.

When investigating a nonlinear recurrence numerically, start on both sides of the proposed equilibrium and choose more than one distance from it. A very short trajectory near a slowly moving state may look stationary at the displayed precision. Record extra digits or record the error explicitly. Then ask what algebraic feature could explain the observed behavior.

\section{Thresholds, strong updates, and model choices}\label{ch02-thresholds-strong-updates-and-model-choices}

A different table can reverse the feedback. For \[M=\begin{pmatrix}4&1\\2&3\end{pmatrix},\] the score gap is \(4x-2\). Below one half, \(x\) decreases; above one half, it increases. The middle fixed point acts as a threshold. Under either positive multiplier rule, trajectories starting below it approach zero and those starting above it approach one. The direction of every step and boundedness give the convergence argument.

For the coexistence table, exponential updates can overshoot. At \(x=1/4\) with \(\eta=2\), \[x^+=\frac{e^4}{e^4+3}\approx0.9479>\frac34.\] Positivity ensures a valid probability, but does not ensure a small correction. This is the same distinction discovered in the affine feedback model.

\begin{figure}[htbp]\centering
\resizebox{.98\linewidth}{!}{\begin{tikzpicture}
\begin{axis}[width=6.7cm,height=5.4cm,xlabel={$A$ proportion $x$},ylabel={Score},xmin=0,xmax=1,ymin=0,ymax=4.5,legend pos=south east]
\addplot[teal,domain=0:1]{4-2*x};\addlegendentry{$f_A$}
\addplot[copper,domain=0:1]{1+2*x};\addlegendentry{$f_B$}
\draw[dashed] (axis cs:.75,0)--(axis cs:.75,4.3);
\end{axis}
\begin{axis}[at={(6.5cm,0)},anchor=south west,width=6.7cm,height=5.4cm,xlabel={Step $t$},ylabel={$A$ proportion},xmin=0,xmax=12,ymin=0,ymax=1,legend pos=south east]
\addplot[teal,mark=*] coordinates {(0,0.25) (1,0.4) (2,0.517483) (3,0.594476) (4,0.643506) (5,0.675429) (6,0.696853) (7,0.711616) (8,0.722007) (9,0.729437) (10,0.734814) (11,0.738741) (12,0.741628)};\addlegendentry{Linear}
\addplot[copper,mark=square*] coordinates {(0,0.25) (1,0.947915) (2,0.78886) (3,0.732469) (4,0.759042) (5,0.745565) (6,0.752236) (7,0.748887) (8,0.750558) (9,0.749721) (10,0.750139) (11,0.74993) (12,0.750035)};\addlegendentry{Exponential}
\addplot[ink,dashed,domain=0:12]{.75};
\end{axis}\end{tikzpicture}}
\caption{Coexistence scores and two finite trajectories at intensity 2. The linear rule approaches the equilibrium from one side; the exponential rule crosses it.}\end{figure}

Adding a common number \(C\) to every current score leaves exponential normalization unchanged, because each weight acquires the same factor \(e^{\eta C}\). Scaling all scores by a positive factor is equivalent to scaling \(\eta\).

For linear multipliers, a common shift changes the effective intensity. When \(1+\eta C>0\), \[1+\eta(f_i+C)=(1+\eta C)\left(1+\frac{\eta}{1+\eta C}f_i\right).\] The common first factor cancels. The resulting trajectory uses intensity \(\eta/(1+\eta C)\) for the original scores. Reporting an update therefore requires the score table, its conversion to multipliers, and its intensity.

Use the sequence view of the \href{https://bobovski66.github.io/IntroDiscreteInfoGeom/demos/operator-explorer.html}{operator explorer} to compare feedback trajectories, and the \href{https://bobovski66.github.io/IntroDiscreteInfoGeom/demos/week02.html}{payoff demonstration} to test the coexistence table. Keep the starting state fixed while changing one parameter. Then compare the signs and magnitudes of successive errors.

With three types, a single proportion no longer specifies the state. Chapter 3 replaces the interval by a triangle and gives us a way to measure changes in several directions.

\section{Exercises}\label{ch02-exercises}

\begin{exercises}

\item 

For \(\Delta x_n=\tfrac12(10-x_n)\) and \(x_0=2\), find four updates and a formula for \(x_n\).

\item 

With \(c=1/2\) and \(x_0=1/4\), calculate the first four states for \(k=3/2\). Give the error multiplier.

\item 

Repeat Exercise 2 for \(k=2\) and \(k=5/2\). Classify the behavior and check whether the values remain probabilities.

\item 

Prove that \(x^+=(1-k)x+kc\) preserves \([0,1]\) when \(c\in[0,1]\) and \(0\leq k\leq1\).

\item 

Explain the difference between solving \(F(c)=0\) and solving \(F(c)=c\). Give a recurrence for which the two equations have different solutions.

\item 

For \(M=\left(\begin{smallmatrix}5&-1\\2&3\end{smallmatrix}\right)\), find \(f_A,f_B,d(x)\) and the interior equal-score state.

\item 

Give a sufficient range of positive \(\eta\) for the linear multipliers in Exercise 6 to remain positive at every composition.

\item 

For the coexistence table, calculate one linear and one exponential update at \(x=1/2\), \(\eta=1\).

\item 

Verify the no-crossing proof for the linear coexistence rule by expanding its average score and error factor.

\item 

Starting at \(x=1/4\), compute four exponential coexistence updates with \(\eta=2\). Record the side of \(3/4\) occupied after each step.

\item 

Prove invariance of exponential normalization under a common score shift. Explain why the same conclusion need not hold for linear multipliers.

\item 

Design a numerical investigation of attraction to a fixed point. State the rule, the parameter, the initial states, the recorded errors, and what remains unproved by a finite run.

\end{exercises}

\section{Selected solutions}\label{ch02-selected-solutions}

\textbf{1.} The values are \(2,6,8,9,19/2\). The error is multiplied by one half, so \(x_n=10-8(1/2)^n\).

\textbf{2.} The first four states, including the initial state, are \(1/4,5/8,7/16,17/32\). The error multiplier is \(-1/2\).

\textbf{3.} For \(k=2\), the values alternate \(1/4,3/4\). For \(k=5/2\), they begin \(1/4,7/8,-1/16,43/32\). The first trajectory is stable without attraction; the second has growing errors and leaves the interval.

\textbf{6.} The scores are \(f_A=6x-1\) and \(f_B=3-x\). Their gap is \(7x-4\), which vanishes at \(4/7\).

\textbf{7.} All scores are at least \(-1\), so \(0<\eta<1\) suffices. At the endpoint \(x=0\), the \(A\) score is exactly \(-1\), explaining the strict upper bound.

\textbf{8.} The scores are \((3,2)\). Linear multipliers give \(x^+=4/7\). Exponential multipliers give \(x^+=e/(e+1)\approx0.7311\).
\labheading\begin{enumerate}
\def\labelenumi{\arabic{enumi}.}
\tightlist
\item
  Consider \(\Delta x_n=\tfrac12(10-x_n)\) with \(x_0=2\). (a) Rewrite as \(x_{n+1}=F(x_n)\). (b) Find the fixed point. (c) Calculate \(x_1,x_2,x_3\). (d) Find the errors relative to the fixed point and their ratio.
\item
  Study \(\Delta x_n=-k(x_n-1/2)\) from \(x_0=1/4\). (a) Compute three steps for \(k=1/2\). (b) Repeat for \(k=3/2\). (c) Describe which sequence crosses the target and why. (d) Check the prediction from \(|1-k|\).
\item
  Use the same equation with \(k=2\) and \(k=5/2\). (a) Compute four steps. (b) Explain how the signs and sizes of the errors change. (c) State why a large corrective step can destabilize a discrete system.
\item
  With payoff matrix \(M=((2,4),(3,1))\), write \(f_A(x)\) and \(f_B(x)\) and solve \(f_A=f_B\). (a) Evaluate the sign of \(f_A-f_B\) at \(x=1/4\) and \(x=9/10\). (b) Use the original demo to compare finite updates with two intensities. (c) Explain why the sign test does not give the step size.
\item
  Choose an update \(x_{n+1}=F(x_n)\) on \([0,1]\). (a) List four iterates for two nearby initial states. (b) Compute each first-difference sequence. (c) Give numerical evidence for attraction or repulsion and distinguish this from a proof.
\end{enumerate}

\FloatBarrier

\chapter{Discrete gradients on the probability simplex}
Suppose a campus survey records three ways of reaching class: walking, cycling, and taking a bus. A distribution such as \((1/2,1/3,1/6)\) says that half the responses are in the first category, a third in the second, and a sixth in the third. Changing one share forces at least one other share to change. The three numbers describe a state with two independent degrees of freedom.

Chapter 2 studied motion on a probability interval. With three categories, the state space becomes a triangle. We will use a finite grid of possible survey compositions to develop directional differences, gradients, circulation, and divergence. Each operation will have a concrete meaning in terms of changing one response or moving a quantity along an edge.

\section{Why probability states form a triangle}\label{ch03-why-probability-states-form-a-triangle}

Write \(p=(p_A,p_B,p_C)\), with nonnegative coordinates adding to one. Knowing \(p_A\) and \(p_B\) determines \(p_C=1-p_A-p_B\). The permitted pairs satisfy \(p_A\geq0\), \(p_B\geq0\), and \(p_A+p_B\leq1\).

It is often more symmetric to place pure states at the vertices of an equilateral triangle: \[A=(0,0),\qquad B=(1,0),\qquad C=(1/2,\sqrt3/2).\] The \textbf{barycentric representation} places the distribution at the weighted point \[X=p_AA+p_BB+p_CC
=\left(p_B+\frac{p_C}{2},\frac{\sqrt3}{2}p_C\right).\] The square root appears because an equilateral triangle of side one has height \(\sqrt3/2\). It only sets the picture's proportions.

At vertex \(A\), all the probability belongs to \(A\). On edge \(AB\), the \(C\) coordinate is zero. The center corresponds to \((1/3,1/3,1/3)\). Lines parallel to edge \(AB\) have constant height, hence constant \(p_C\).

\begin{example}[Locating and reading a mixture]

For \(p=(1/2,1/3,1/6)\), the plotted point is \((5/12,\sqrt3/12)\). Its height is one sixth of the triangle's total height, matching \(p_C=1/6\).

Conversely, if a plotted point has coordinates \((u,v)\), solve the plotting equations: \[p_C=\frac{2v}{\sqrt3},\qquad
p_B=u-\frac{v}{\sqrt3},\qquad
p_A=1-u-\frac{v}{\sqrt3}.\] At \((u,v)=(1/2,\sqrt3/6)\), these give three equal probabilities.

\end{example}

The set of probability vectors on \(k\) labels is called the \textbf{probability simplex}. Two labels give an interval, three give a triangle, and four give a tetrahedron. The definition is always the same: nonnegative entries with total one. We can calculate with more labels even when we cannot draw their simplex.

\section{Fixing a number of tokens makes a finite lattice}\label{ch03-fixing-a-number-of-tokens-makes-a-finite-lattice}

Imagine a survey of exactly \(N\) responses. Counts \((a,b,c)\) satisfy \[a,b,c\text{ are nonnegative integers},\qquad a+b+c=N.\] Their empirical proportions are \(p=(a/N,b/N,c/N)\). The possible points form a triangular \textbf{lattice}, a finite grid inside the simplex. We denote the set by \(\mathcal P_N\).

For \(N=4\), group the triples by their \(C\) count:

\begin{longtable}[]{@{}lll@{}}
\toprule\noalign{}
Count \(c\) & Possible pairs \((a,b)\) & Number \\
\midrule\noalign{}
\endhead
\bottomrule\noalign{}
\endlastfoot
0 & \((4,0),(3,1),(2,2),(1,3),(0,4)\) & 5 \\
1 & \((3,0),(2,1),(1,2),(0,3)\) & 4 \\
2 & \((2,0),(1,1),(0,2)\) & 3 \\
3 & \((1,0),(0,1)\) & 2 \\
4 & \((0,0)\) & 1 \\
\end{longtable}

The total is fifteen. In general, fixing \(c\) leaves \(N-c+1\) choices, so the number of states is \[1+2+\cdots+(N+1)=\frac{(N+1)(N+2)}2.\] Increasing \(N\) increases the resolution: one response is a smaller fraction of the whole.

A \textbf{neighboring state} is obtained by transferring one token from one category to another. From \((2,1,1)\), moving a token from \(C\) to \(A\) gives \((3,1,0)\). The probability displacement is \((1/4,0,-1/4)\). Every transfer preserves the total count and the total probability.

\begin{figure}[htbp]\centering
\resizebox{.98\linewidth}{!}{\begin{tikzpicture}\draw[ink,thick] (0,0)--(6,0)--(3,5.196)--cycle;\node[below left] at (0,0){$A$};\node[below right] at (6,0){$B$};\node[above] at (3,5.196){$C$};\fill[teal] (3,5.19615) circle (2pt);\fill[teal] (3.75,3.89711) circle (2pt);\fill[teal] (4.5,2.59808) circle (2pt);\fill[teal] (5.25,1.29904) circle (2pt);\fill[teal] (6,0) circle (2pt);\fill[teal] (2.25,3.89711) circle (2pt);\fill[teal] (3,2.59808) circle (2pt);\fill[teal] (3.75,1.29904) circle (2pt);\fill[teal] (4.5,0) circle (2pt);\fill[teal] (1.5,2.59808) circle (2pt);\fill[teal] (2.25,1.29904) circle (2pt);\fill[teal] (3,0) circle (2pt);\fill[teal] (0.75,1.29904) circle (2pt);\fill[teal] (1.5,0) circle (2pt);\fill[teal] (0,0) circle (2pt);\node[anchor=west,align=left] at (6.8,3.7){$N=4$: fifteen states\\One token changes two counts.};\draw[copper,thick] (2.25,1.29904) circle (5pt);\node[below=6pt] at (2.25,1.29904){$(2,1,1)$};\draw[copper,->] (2.25,1.29904)--(3.75,1.29904);\draw[copper,->] (2.25,1.29904)--(3,2.59808);\draw[copper,->] (2.25,1.29904)--(0.75,1.29904);\draw[copper,->] (2.25,1.29904)--(1.5,2.59808);\draw[copper,->] (2.25,1.29904)--(1.5,0);\draw[copper,->] (2.25,1.29904)--(3,0);\end{tikzpicture}}
\caption{The denominator-four token lattice. Arrows show the six one-token moves from counts (2,1,1).}\end{figure}

At \((2,1,1)\) all three categories can donate, so there are six neighbors. At \((4,0,0)\) only \(A\) can donate, giving two neighbors. Boundary restrictions are part of the geometry, not a failure of the formula.

These tokens represent units of composition. A token transfer through probability space need not describe a person physically traveling between locations. Chapter 4 will introduce another graph whose vertices are physical or categorical locations. Keeping these two meanings of a vertex separate will help.

The lattice is also a graph: its count states are vertices and its legal transfers are edges. A drawing of the triangle helps us see the graph, but the definition does not depend on the drawing. A computer can store each triple and list its legal neighbors directly.

To generate those neighbors, choose a donor category with positive count and a different receiver. Subtract one from the donor and add one to the receiver. There are six ordered donor-receiver pairs for three categories, but some are forbidden at the boundary. This procedure gives both an exact state space and an exact rule for adjacency.

The number of interior states can be counted separately. A strictly positive triple has at least one token in each category. Reserve those three tokens; the remaining \(N-3\) are distributed freely. For \(N\geq3\), there are \((N-1)(N-2)/2\) interior states. Subtracting from the full count gives \(3N\) boundary states. Vertices belong to two triangle edges in the picture but are counted only once in this total.

For \(N=4\), there are three interior states and twelve boundary states. Thus boundaries are a large part of a coarse lattice. Ignoring illegal moves would change the model substantially, rather than introduce only a small numerical error.

\section{A directional difference compares two permitted states}\label{ch03-a-directional-difference-compares-two-permitted-states}

A \textbf{scalar field} assigns one number to every state. For a first example, consider \[C(p)=p_A^2+p_B^2+p_C^2.\] This concentration score is one at a pure state and one third at the balanced state. Larger values indicate a more concentrated composition.

Let \(e_A=(1,0,0)\), \(e_B=(0,1,0)\), and \(e_C=(0,0,1)\). These are unit coordinate vectors. A transfer from \(C\) to \(A\) changes \(p\) by \((e_A-e_C)/N\). Write \[T_{A\leftarrow C}p=p+\frac{e_A-e_C}{N}.\] The arrow names the receiving and donating categories. The move is legal when \(p_C\geq1/N\).

For any field \(F\), define its \textbf{directional difference} \[\boxed{D_{A\leftarrow C}F(p)=F(T_{A\leftarrow C}p)-F(p).}\] First evaluate the field after the move, then subtract its value before the move. This is the same operation as a first difference, now taken in a specified direction through a state space.

\begin{example}[Two directions from the same composition]

At counts \((2,1,1)\) with \(N=4\), we have \(p=(1/2,1/4,1/4)\) and \(C(p)=3/8\).

Moving \(C\) to \(A\) gives \((3/4,1/4,0)\), with concentration \(5/8\). The directional difference is \(1/4\).

Moving \(A\) to \(B\) gives \((1/4,1/2,1/4)\), again with concentration \(3/8\). That directional difference is zero. The state changed, but this field assigns the same number to the two states.

\end{example}

Our \(D\) records the total change caused by one token. A difference quotient per unit of transferred probability would divide by \(1/N\) and equal \(ND_{A\leftarrow C}F\). A quotient per unit of ordinary Euclidean displacement would use the step length \(\sqrt2/N\). Stating the convention prevents factors from entering unnoticed. We use unscaled edge differences throughout the graph calculus.

We can derive concentration changes for every transfer, rather than calculate a new square sum each time. Let \(h=1/N\), and transfer from category \(j\) to category \(i\). Only those two terms change: \[\begin{aligned}
C(p+h(e_i-e_j))-C(p)
&=(p_i+h)^2-p_i^2+(p_j-h)^2-p_j^2\\
&=2h(p_i-p_j)+2h^2.
\end{aligned}\] The \(2h^2\) term is a finite-step correction. If the receiver and donor have equal probabilities, the transfer makes them unequal and raises concentration by \(2h^2\). If the donor has exactly one more token than the receiver, the move swaps their counts and leaves concentration unchanged.

This is a useful example of what a discrete derivative retains: the actual size of the move. Replacing it by only its term proportional to \(h\) would miss an exact effect at equal coordinates. Later, the local Fisher formula will also be presented alongside an exact finite expression, so that the effect of approximation remains visible.

We can compare fields on the same graph without changing its geometry of allowed moves. A cost field might assign greater weight to one survey response than another, while concentration and entropy treat labels symmetrically. Directional differences depend jointly on the chosen field and the chosen edge.

\section{An information field for the same lattice}\label{ch03-an-information-field-for-the-same-lattice}

A second scalar field is \textbf{entropy}: \[H(p)=-\sum_i p_i\ln p_i,\qquad 0\ln0:=0.\] Here \(\ln z\) is the natural logarithm: the exponent needed to express \(z\) as a power of \(e\). For example, \(\ln1=0\) and \(\ln(1/2)\approx-0.6931\). A calculator's \(\ln\) key is enough to evaluate this field. Chapter 5 will explain its meaning and derive its main properties.

For now, entropy measures uncertainty in the category of one randomly selected response. Pure states have entropy zero. Spreading probability more evenly generally gives greater uncertainty.

\begin{example}[Computing an entropy difference]

For \(p=(1/2,1/4,1/4)\), \[H(p)=-\frac12\ln\frac12-2\left(\frac14\ln\frac14\right)
=\frac32\ln2\approx1.0397.\] After the transfer \(C\) to \(A\), the state is \(q=(3/4,1/4,0)\) and \[H(q)=-\frac34\ln\frac34-\frac14\ln\frac14\approx0.5623.\] Thus \(D_{A\leftarrow C}H(p)\approx-0.4774\). Concentration increased while uncertainty decreased. The transfer \(A\) to \(B\) merely permutes the entries and leaves entropy unchanged.

\end{example}

A zero contribution for a zero probability means that an outcome which never occurs contributes nothing to the average. We do not ask the calculator to evaluate \(\ln0\).

\section{What a discrete gradient records}\label{ch03-what-a-discrete-gradient-records}

In an interval, there are two possible directions from an interior point. In a triangle, there are several. A \textbf{discrete gradient field} records changes of a scalar field along edges.

For an oriented edge \(p\to q\), define \[A_F(p,q)=F(q)-F(p).\] Reversing the edge reverses the sign: \(A_F(q,p)=-A_F(p,q)\).

Two independent lattice directions are \(u=e_A-e_C\) and \(v=e_B-e_C\). Where both forward moves are legal, one may record the pair \[\nabla_dF(p)=(D_uF(p),D_vF(p)).\] It lists changes in a chosen pair of directions. The directions meet at an angle in the triangle, so this pair is a coordinate record, rather than automatically an ordinary arrow perpendicular to contours. The full edge field is often easier to use because it also works at boundaries.

A simple ascent rule examines every legal neighbor, chooses one with the greatest positive increase in \(F\), and stays put if none improves it. Ties require a fixed choice rule. Because the lattice has finitely many states and every actual move strictly raises \(F\), this rule cannot revisit a state. It must stop after finitely many moves.

The stopping state is a local maximum with respect to the permitted edges. For a general field, a local maximum need not be the largest value anywhere on the graph. For entropy on this complete token-transfer lattice, balancing the counts identifies the maxima; Chapter 5 explains why.

A reproducible ascent rule requires more than the phrase ``move uphill.'' At a state, list its legal neighbors, evaluate the scalar field at each, and compare the increases. Select the largest strictly positive increase. If two increases tie, use a declared ordering of categories or another fixed rule. If none is positive, return the same state.

This algorithm needs no formula for a continuous gradient. On a graph with \(S\) states, it can make at most \(S-1\) strict-improvement moves, because each visited field value is greater than its predecessor and no vertex can recur. The argument does not guarantee that the selected path is the shortest route to a maximum. Ascent and shortest-path optimization are different tasks.

Allowing equal-valued moves changes the termination argument. Two adjacent states with the same value can be visited back and forth forever unless a tie rule prevents it. For entropy, permutations of a balanced count vector can provide such plateaus. The requirement of strict improvement is therefore part of the dynamics.

\section{Path sums reveal the cancellation behind gradients}\label{ch03-path-sums-reveal-the-cancellation-behind-gradients}

Suppose a path visits states \(p_0,p_1,p_2,p_3\). Adding gradient edge values gives \[\begin{aligned}
A_F(p_0,p_1)+A_F(p_1,p_2)+A_F(p_2,p_3)
&=[F(p_1)-F(p_0)]\\
&\quad+[F(p_2)-F(p_1)]+[F(p_3)-F(p_2)]\\
&=F(p_3)-F(p_0).
\end{aligned}\] The intermediate values cancel just as they did in Chapter 1. Any two paths with the same endpoints therefore have the same sum of gradient differences. Around a closed path, the sum is zero.

For an arbitrary signed edge field, the sum around a closed route is its \textbf{circulation}. A gradient has zero circulation around every closed route. A field that assigns one unit in a clockwise direction around a triangular face has circulation three around that face, so it cannot be a gradient of vertex values.

This test has a useful interpretation. If local changes are supposed to arise from one underlying score, following them around a loop must return that score to its starting value. A nonzero loop total reveals that the edge values contain additional directional information.

\section{Flux and divergence describe a different feature}\label{ch03-flux-and-divergence-describe-a-different-feature}

An edge field may instead describe \textbf{flux}, an amount transported in one step. Choose an orientation for each edge. Positive flux follows that arrow; negative flux goes against it. The \textbf{outward divergence} at a vertex is outgoing flux minus incoming flux.

For a clockwise unit flow around a triangle, one unit arrives and one unit leaves at every vertex. Every divergence is zero, but the circulation is three. Zero accumulation at a vertex can coexist with sustained circulation through it.

Conversely, a vertex may send mass to several neighbors and receive little. Its divergence is positive, so its stored mass decreases. This is the reason for the minus sign in the continuity equation developed in Chapter 4.

For any chosen collection of vertices, add their divergences. Every edge joining two vertices in the collection contributes once positively and once negatively, so it cancels. Only edges crossing the boundary remain: \[\boxed{\text{sum of outward divergences}
=\text{net flux leaving the collection}.}\] This is a finite version of a divergence theorem. Its content is exact accounting along edges.

\begin{example}[Divergence on a small directed flow diagram]

Take three vertices \(U,V,W\) with net flows two along \(U\to V\), one along \(V\to W\), and three along \(U\to W\). Outward divergence is five at \(U\), minus one at \(V\), and minus four at \(W\). The sum is zero.

Now consider only the collection \(\{U,V\}\). Its divergences sum to four. The internal flow of two from \(U\) to \(V\) cancels in that sum. The surviving outward flows are three from \(U\) to \(W\) and one from \(V\) to \(W\), also totaling four.

\end{example}

This example separates three levels of description. An edge stores a directional transport amount. A vertex divergence records net departure there. A collection of vertices exchanges mass only through its boundary edges. Moving from one level to the next requires addition and cancellation, rather than an approximation.

The graph in this flow example could represent actual locations. The token graph instead has whole probability distributions as its vertices. Both support the same divergence bookkeeping, but a field on probability states and a probability stored at physical locations should not be confused. The next chapter will explicitly choose locations as vertices and put probability mass on them.

\section{Application: cyclic motion need not be a gradient}\label{ch03-application-cyclic-motion-need-not-be-a-gradient}

Three competing strategies can favor one another cyclically. Consider the balanced score table \[M=\begin{pmatrix}0&1&-1\\-1&0&1\\1&-1&0\end{pmatrix}.\] At state \(p\), its scores are \[f_A=p_B-p_C,\quad f_B=p_C-p_A,\quad f_C=p_A-p_B.\] At the center all three are zero. Away from the center, the advantages encourage circulation rather than simple movement up one coordinate.

Exponential normalization gives \(p_i^+=p_i e^{\eta f_i}/Z\), with \(Z=\sum_i p_i e^{\eta f_i}\). At \(p=(1/2,1/3,1/6)\) the scores are \((1/6,-1/3,1/6)\). With \(\eta=1\), the next state is approximately \((0.5755,0.2327,0.1918)\). The scores must be recomputed before the next step because the composition has changed.

We can examine one exact diagnostic, the product \(P=p_Ap_Bp_C\). Direct cancellation gives \(\sum_i f_i=0\) and \(\sum_i p_if_i=0\). Multiplying the update equations therefore yields \[P(p^+)=\frac{P(p)}{Z^3}.\]

The weighted arithmetic-geometric mean inequality says that for positive numbers \(a_i\) and nonnegative weights \(p_i\) adding to one, \[\sum_i p_i a_i\geq\prod_i a_i^{p_i}.\] It generalizes the fact that the average of two positive numbers is at least their geometric mean. For rational weights it follows by repeating numbers in an ordinary arithmetic-geometric mean inequality. General weights follow by approximation.

Apply it to \(a_i=e^{\eta f_i}\). We obtain \[Z\geq e^{\eta\sum_i p_if_i}=1.\] For an interior state other than the center, the scores differ, giving strict inequality. Consequently \(P(p^+)<P(p)\). A finite cyclic update decreases this product; it does not conserve it.

The product is largest at the center, where it equals \(1/27\). Its decrease describes crossing toward lower product contours. It does not prove that every coordinate approaches zero or that the trajectory has a particular limit. A single scalar diagnostic gives one piece of a multidimensional story.

\begin{figure}[htbp]\centering
\resizebox{.98\linewidth}{!}{\begin{tikzpicture}\draw[ink,thick] (0,0)--(6,0)--(3,5.196)--cycle;\node[below left] at (0,0){$A$};\node[below right] at (6,0){$B$};\node[above] at (3,5.196){$C$};\draw[teal,line width=1pt,] plot coordinates {(2.5,0.866025) (2.30211,0.914991) (2.13433,0.989415) (2.00434,1.09113) (1.91704,1.22169) (1.87454,1.3817) (1.87635,1.57002) (1.91959,1.78302) (1.99934,2.014) (2.10912,2.25316) (2.2416,2.48824) (2.38957,2.70569) (2.54686,2.89222) (2.70911,3.03607) (2.87405,3.12795) (3.04129,3.16149) (3.21161,3.13339) (3.38597,3.04357) (3.56442,2.8955) (3.74508,2.69646) (3.92355,2.45754) (4.09295,2.19293) (4.24443,1.91836) (4.3683,1.64902) (4.45503,1.39764) (4.4961,1.17313) (4.48445,0.98029) (4.41464,0.820371) (4.28327,0.692054) (4.08965,0.592573) (3.83693,0.518653) (3.53324,0.467199) (3.19222,0.435747) (2.83231,0.422736) (2.4746,0.427675) (2.13975,0.451239) (1.84514,0.495337) (1.60317,0.563116) (1.42101,0.65886) (1.30151,0.787681) (1.24443,0.954869) (1.2474,1.16479) (1.30623,1.41932) (1.41474,1.71605) (1.56431,2.04684) (1.744,2.39749) (1.94142,2.74913) (2.14444,3.08112) (2.34321,3.37439) (2.5317,3.61395) (2.70818,3.78973) (2.87485,3.89591) (3.03679,3.92947) (3.20062,3.88899) (3.37309,3.7742) (3.55961,3.58665) (3.76275,3.33128) (3.98096,3.01809) (4.20804,2.66311) (4.4337,2.28732) (4.64535,1.91367) (4.83045,1.56301) (4.97842,1.25062) (5.08144,0.984702) (5.13404,0.766899) (5.13197,0.594087) (5.07111,0.460494) (4.94694,0.359436) (4.75496,0.28447) (4.49214,0.229994) (4.15941,0.191484) (3.76431,0.165501) (3.32271,0.149618) (2.85811,0.142316) (2.39816,0.142908) (1.96937,0.151513) (1.59225,0.169073) (1.27893,0.197432) (1.03359,0.239463) (0.85476,0.299217) (0.738192,0.382049) (0.679138,0.494619) (0.673725,0.644618) (0.719298,0.839994) (0.813842,1.08747) (0.954673,1.39033) (1.13681,1.74578) (1.35164,2.14306) (1.58679,2.56335) (1.82756,2.98255) (2.05966,3.37612) (2.27206,3.72397) (2.4587,4.01334) (2.61864,4.23892) (2.75509,4.40077) (2.87403,4.50162) (2.98296,4.54436) (3.0901,4.53041) (3.20382,4.45901) (3.33225,4.3274) (3.48251,4.13207) (3.65971,3.87086) (3.86549,3.5458) (4.09666,3.16584) (4.34467,2.74806) (4.5966,2.31618) (4.83773,1.89618) (5.05474,1.51085) (5.23811,1.17555) (5.38285,0.896891) (5.48765,0.673976) (5.55317,0.501037) (5.58036,0.370035) (5.56908,0.272596) (5.51723,0.201121) (5.42044,0.149249) (5.27216,0.111934) (5.0646,0.0853168) (4.79047,0.0665215) (4.44576,0.053457) (4.0331,0.044639) (3.56458,0.0390514) (3.06221,0.0360444) (2.555,0.0352672) (2.07288,0.0366313) (1.64035,0.0403013) (1.27234,0.0467149) (0.973804,0.0566307) (0.74195,0.0712098) (0.569614,0.0921336) (0.448243,0.121762) (0.369969,0.163332) (0.328758,0.221175) (0.320881,0.300924) (0.344943,0.409613) (0.401587,0.555512) (0.492862,0.74747) (0.621167,0.993495) (0.78772,1.29846) (0.990736,1.6612) (1.22388,2.07208) (1.47591,2.51243) (1.73205,2.95717) (1.9771,3.38005) (2.19865,3.75943) (2.38931,4.08194) (2.54698,4.34304) (2.67379,4.54499) (2.7745,4.6939) (2.855,4.79702) (2.92146,4.86085) (2.97988,4.89007) (3.03604,4.88695) (3.09566,4.85123) (3.16458,4.78006) (3.24888,4.66828) (3.35478,4.50906) (3.48814,4.29509) (3.65339,4.02072) (3.85203,3.68474) (4.08103,3.29339)};
\begin{axis}[at={(7cm,0.2cm)},anchor=south west,width=6.2cm,height=5.2cm,xlabel={Step $t$},ylabel={$p_Ap_Bp_C$},xmin=0,xmax=160,ymin=0]
\addplot[teal] coordinates {(0,0.0277778) (1,0.0275013) (2,0.0272369) (3,0.0269819) (4,0.0267289) (5,0.026468) (6,0.0261894) (7,0.0258869) (8,0.0255602) (9,0.025215) (10,0.0248622) (11,0.0245144) (12,0.024182) (13,0.02387) (14,0.0235766) (15,0.0232934) (16,0.0230078) (17,0.0227064) (18,0.0223788) (19,0.0220213) (20,0.0216382) (21,0.021241) (22,0.0208458) (23,0.0204676) (24,0.0201169) (25,0.019796) (26,0.0194993) (27,0.0192144) (28,0.0189254) (29,0.0186171) (30,0.0182787) (31,0.0179075) (32,0.01751) (33,0.0171006) (34,0.0166977) (35,0.0163182) (36,0.0159732) (37,0.0156654) (38,0.0153888) (39,0.0151309) (40,0.0148755) (41,0.014606) (42,0.0143097) (43,0.0139805) (44,0.0136212) (45,0.0132435) (46,0.0128648) (47,0.0125037) (48,0.0121745) (49,0.0118844) (50,0.0116325) (51,0.0114107) (52,0.0112064) (53,0.0110044) (54,0.0107901) (55,0.0105519) (56,0.0102839) (57,0.00998763) (58,0.00967252) (59,0.00935355) (60,0.00904742) (61,0.00876811) (62,0.00852376) (63,0.00831568) (64,0.00813921) (65,0.00798544) (66,0.00784306) (67,0.00770007) (68,0.00754534) (69,0.00737022) (70,0.00717027) (71,0.00694668) (72,0.00670663) (73,0.00646191) (74,0.00622596) (75,0.00601042) (76,0.00582259) (77,0.00566453) (78,0.00553376) (79,0.00542456) (80,0.00532941) (81,0.00524001) (82,0.00514812) (83,0.00504626) (84,0.0049286) (85,0.00479199) (86,0.00463703) (87,0.00446848) (88,0.00429461) (89,0.00412528) (90,0.00396945) (91,0.00383319) (92,0.00371888) (93,0.00362559) (94,0.00355002) (95,0.00348759) (96,0.00343313) (97,0.0033814) (98,0.00332729) (99,0.00326618) (100,0.00319423) (101,0.00310906) (102,0.00301042) (103,0.00290072) (104,0.00278492) (105,0.00266959) (106,0.00256129) (107,0.00246505) (108,0.00238349) (109,0.00231688) (110,0.0022637) (111,0.00222139) (112,0.00218692) (113,0.00215719) (114,0.00212915) (115,0.00209989) (116,0.00206665) (117,0.00202701) (118,0.00197908) (119,0.00192197) (120,0.00185626) (121,0.00178419) (122,0.00170947) (123,0.00163651) (124,0.00156934) (125,0.00151076) (126,0.00146195) (127,0.0014227) (128,0.00139184) (129,0.00136771) (130,0.0013485) (131,0.00133244) (132,0.00131786) (133,0.00130316) (134,0.00128681) (135,0.00126735) (136,0.00124346) (137,0.0012141) (138,0.00117883) (139,0.00113805) (140,0.0010932) (141,0.00104661) (142,0.00100108) (143,0.000959129) (144,0.000922544) (145,0.0008921) (146,0.000867697) (147,0.000848643) (148,0.000833956) (149,0.000822582) (150,0.000813507) (151,0.000805804) (152,0.00079862) (153,0.000791151) (154,0.000782608) (155,0.000772206) (156,0.000759189) (157,0.000742912) (158,0.000722983) (159,0.000699445) (160,0.000672936)};
\end{axis}\end{tikzpicture}}
\caption{Balanced cyclic exponential updates at intensity 0.35. The triangle trajectory circulates while the coordinate product decreases.}\end{figure}

For comparison, the permutation \(T(p_A,p_B,p_C)=(p_C,p_A,p_B)\) has \(T^3(p)=p\) and preserves the product exactly. Two rules can both look cyclic while having different invariants. An \textbf{invariant} is a quantity unchanged at every step of the particular recurrence.

Use the simplex view of the \href{https://bobovski66.github.io/IntroDiscreteInfoGeom/demos/operator-explorer.html}{operator explorer} to compare neighboring entropy values. The \href{https://bobovski66.github.io/IntroDiscreteInfoGeom/demos/week03.html}{weekly demonstration} also explores lattices and cyclic selection. Chapter 4 turns the flux bookkeeping into a model for transporting actual probability among labeled locations.

\section{Exercises}\label{ch03-exercises}

\begin{exercises}

\item 

Plot \((1,0,0)\), \((0,1/2,1/2)\), and \((1/4,1/4,1/2)\) in barycentric coordinates.

\item 

List all six states for \(N=2\), and verify the state-count formula.

\item 

List the neighbors of \((2,1,1)\) for \(N=4\). Explain why the corner \((4,0,0)\) has fewer neighbors.

\item 

For \(N=4\) and counts \((2,1,1)\), calculate the concentration differences for transfers \(A\) to \(B\) and \(B\) to \(C\).

\item 

Calculate the entropy difference for the transfer from counts \((2,1,1)\) to \((2,2,0)\). Compare its sign with the concentration change.

\item 

Explain why \(D_{A\leftarrow C}F\) and \(D_{C\leftarrow A}F\) at the same starting state need not be negatives of each other. Identify the edge reversal that does give the negative.

\item 

Choose two paths from \((4,0,0)\) to \((2,2,0)\). Add the concentration differences along each path and compare the totals.

\item 

Give a scalar potential on three vertices with values \(0,2,5\). Calculate its three gradient edge values and the circulation around the triangle.

\item 

Assign a clockwise unit flux around a triangle. Calculate divergence at each vertex and circulation around the triangle.

\item 

Count the strictly positive three-category states for \(N=10\) by first placing one token in each category.

\item 

For the balanced cyclic table, verify \(\sum_i f_i=0\) and \(\sum_i p_if_i=0\). Derive the product update identity.

\item 

Explain why strict ascent of a scalar field must stop on a finite graph. Describe what could change if equal-valued moves were allowed.

\end{exercises}

\section{Selected solutions}\label{ch03-selected-solutions}

\textbf{2.} The states are \((2,0,0),(1,1,0),(1,0,1),(0,2,0),(0,1,1),(0,0,2)\) in counts. The formula gives \((3)(4)/2=6\).

\textbf{3.} The neighbors are \((1,2,1),(1,1,2),(3,0,1),(2,0,2),(3,1,0),(2,2,0)\). At the corner only the occupied category can donate.

\textbf{4.} The initial concentration is \(3/8\). The first transfer gives \((1,2,1)\) and zero change. The second gives \((2,0,2)\) with concentration \(1/2\), hence change \(1/8\).

\textbf{5.} The new entropy is \(\ln2\approx0.6931\), giving change \(-\tfrac12\ln2\approx-0.3466\). Concentration rises from \(3/8\) to \(1/2\), so its change is positive.

\textbf{8.} Orient the triangle \(U\to V\to W\to U\). The edge values are \(2,3,-5\). Their circulation is zero.

\textbf{10.} After reserving three tokens, the remaining seven can be distributed in \((8)(9)/2=36\) ways. The total number of states, including the boundary, is sixty-six.
\labheading\begin{enumerate}
\def\labelenumi{\arabic{enumi}.}
\tightlist
\item
  Take \(N=4\). (a) List all 15 count triples \((a,b,c)\) adding to four. (b) Draw their triangular arrangement or use the lattice demo. (c) Count neighbors of \((2,1,1)\) and \((4,0,0)\). (d) Explain why the boundary has fewer neighbors.
\item
  Starting at \((2,1,1)\), move one token from C to A and one from A to B, in separate trials. (a) Record both new count triples. (b) Convert the triples to probabilities. (c) Calculate the entropy at each distribution, taking \(0\log0=0\). (d) Determine the signs of both directional differences.
\item
  At an interior state \(p\), define \(u=e_A-e_C\) and \(v=e_B-e_C\). (a) Write \(p+u/N\) and \(p+v/N\). (b) Explain when each transfer is legal. (c) Evaluate the two differences of \(H\) at \(p=(1/4,1/4,1/2)\) with \(N=4\).
\item
  Choose two different three-step routes on the token lattice with the same start and finish. (a) Label the entropy at each visited point. (b) Add the successive entropy differences along each route. (c) Explain why the route sums agree, even if the intermediate values differ.
\item
  Draw a triangle of three allowed states and assign a clockwise flow of one unit along each edge. (a) Compute divergence at each vertex. (b) Compute the circulation around the triangle. (c) Explain how zero divergence can coexist with nonzero circulation.
\end{enumerate}

\FloatBarrier

\chapter{Probability transport and conservation}
A shared bicycle can be at one of several stations at the end of a day. A mobile device can connect to one of several towers. A system can report one of several operating modes. In each case, we can describe uncertainty about its location or mode by a probability vector, then ask how that vector changes after one step.

The edge bookkeeping from Chapter 3 gives us a local description of these changes: what leaves one location enters another. A stochastic matrix packages all the transition rules into a single table. We will derive the matrix update from individual contributions, recover its currents, and then explain how a finite collection of observations fluctuates around the predicted probabilities.

\section{What one transition probability means}\label{ch04-what-one-transition-probability-means}

Let the possible locations be \(A,B,C\). The entry \(K_{ij}\) will mean the probability of arriving at location \(i\) when the starting location is \(j\). The first index gives the destination; the second gives the origin.

Each column therefore lists all possible destinations from one origin. Its entries must be nonnegative and add to one: \[K_{ij}\geq0,\qquad \sum_i K_{ij}=1.\] Such a table is called \textbf{column-stochastic}. The diagonal entry \(K_{jj}\) is the probability of remaining at \(j\).

Consider a path \(A\)--\(B\)--\(C\). During one step, send one quarter of the mass at each location to each adjacent location and retain the rest. From \(A\), the only neighbor is \(B\), so three quarters remain. From \(B\), there are two neighbors, so one half remains. The table is \[K=\begin{pmatrix}
3/4&1/4&0\\
1/4&1/2&1/4\\
0&1/4&3/4
\end{pmatrix}.\] Reading the middle column gives the destinations from \(B\): a quarter at \(A\), half at \(B\), and a quarter at \(C\).

This table can model a random transition of one object or a deterministic redistribution of a divisible mass. Both interpretations use the same arithmetic, but they describe different kinds of uncertainty.

\section{Matrix multiplication adds contributions from origins}\label{ch04-matrix-multiplication-adds-contributions-from-origins}

Suppose the current distribution is \(p=(p_A,p_B,p_C)\). To calculate the next probability at \(A\), add contributions from every origin: \[p_A^+=\frac34p_A+\frac14p_B+0p_C.\] The same calculation for every destination gives \[p_i^+=\sum_j K_{ij}p_j.\] This collection of weighted sums is \textbf{matrix multiplication}, written \(p^+=Kp\). We regard \(p\) as a column when multiplying, even when displaying it compactly as a tuple.

\begin{example}[Two steps starting at one location]

Start at \(p_0=(1,0,0)\). Multiplication selects the first column: \[p_1=\left(\frac34,\frac14,0\right).\] For the next step, \[\begin{aligned}
p_{A,2}&=\frac34\frac34+\frac14\frac14=\frac58,\\
p_{B,2}&=\frac14\frac34+\frac12\frac14=\frac5{16},\\
p_{C,2}&=\frac14\frac14=\frac1{16}.
\end{aligned}\] Thus \(p_2=(10/16,5/16,1/16)\). Probability reaches \(C\) only after passing through \(B\). The entries add to one, which checks the arithmetic.

\end{example}

A matrix of transition probabilities defines a \textbf{Markov update} when the next distribution depends on the present state through this rule, without additional dependence on earlier history. A model with delays may need a larger state that records the relevant history.

The identity matrix \(I\) leaves every vector unchanged. Subtracting \(p_n=Ip_n\) gives a vector difference equation: \[\boxed{p_{n+1}=Kp_n,\qquad \Delta p_n=(K-I)p_n.}\] A difference is now a vector of coordinate changes. The definition still means next minus current.

The calculation of \(p_{A,2}=5/8\) can also be read as a sum over two-step histories. Starting at \(A\), the histories that end at \(A\) are \(A\to A\to A\) and \(A\to B\to A\). Their probabilities are \((3/4)(3/4)\) and \((1/4)(1/4)\). Adding them gives \(10/16\). Matrix multiplication performs this path accounting simultaneously for all destinations.

This interpretation explains why transition matrices compose by multiplication. If one step uses \(K\) and the next uses \(R\), the final distribution is \(R(Kp)=(RK)p\). An entry of \(RK\) adds probabilities over all possible intermediate locations. Reversing the steps gives \(KR\), which need not be the same matrix. The chronological order is read from right to left in the product because the rightmost operation acts first.

A column convention should be tested on a concentrated input. Multiplying by \((1,0,0)\) must reproduce the destinations from \(A\). This quick check is especially useful when consulting a source that uses row distributions instead. Both conventions are valid when used consistently; mixing them changes the transition rule.

\section{Conservation follows from column sums}\label{ch04-conservation-follows-from-column-sums}

The sum of all new probabilities is \[\begin{aligned}
\sum_i p_i^+
&=\sum_i\sum_j K_{ij}p_j
=\sum_j p_j\left(\sum_i K_{ij}\right)\\
&=\sum_j p_j=1.
\end{aligned}\] We have exchanged the order of two finite sums: adding a rectangular collection by rows or by columns gives the same total. The column-sum condition then supplies the conservation law.

Nonnegativity is also preserved because each new entry is a sum of nonnegative products. The matrix therefore sends probability vectors to probability vectors.

Taking the difference of the two total masses gives \[\sum_i\Delta p_n(i)=0.\] Some coordinates can rise and others fall, but their changes balance. For the two-step example, \[\Delta p_0=(-1/4,1/4,0),\qquad
\Delta p_1=(-1/8,1/16,1/16).\] Both vectors have total zero.

\section{Currents recover the local picture}\label{ch04-currents-recover-the-local-picture}

The amount traveling from \(i\) to \(j\) is \(K_{ji}p_i\). The amount traveling the other way is \(K_{ij}p_j\). Define the \textbf{net current from \(i\) to \(j\)} by \[J_{ij}=K_{ji}p_i-K_{ij}p_j.\] Then \(J_{ji}=-J_{ij}\). A negative net current simply means that the opposite direction carries more mass.

For the symmetric path example, \[J_{AB}=\frac14(p_A-p_B),\qquad
J_{BC}=\frac14(p_B-p_C).\] Mass flows from the larger probability to the smaller one along each edge. With \(p=(1,0,0)\), the currents are \(J_{AB}=1/4\) and \(J_{BC}=0\).

Choose arrows \(A\to B\) and \(B\to C\). The outward divergences are \[\operatorname{div}J(A)=J_{AB},\quad
\operatorname{div}J(B)=J_{BC}-J_{AB},\quad
\operatorname{div}J(C)=-J_{BC}.\] At \(B\), the current on \(AB\) is incoming and therefore subtracted. The update is \[\boxed{\Delta p(i)=-\operatorname{div}J(i).}\] This \textbf{continuity equation} says that stored mass decreases by the net amount sent outward.

The equation holds for a general column-stochastic matrix as well. Subtracting outgoing transport from incoming transport gives \[\sum_{j\ne i}(K_{ij}p_j-K_{ji}p_i)
=\sum_jK_{ij}p_j-p_i,\] where the last step uses the sum of column \(i\). The diagonal stay-put term cancels correctly.

\begin{example}[Checking currents at the second step]

At \(p_1=(3/4,1/4,0)\), we obtain \(J_{AB}=1/8\) and \(J_{BC}=1/16\). Thus the divergences are \((1/8,-1/16,-1/16)\). Their negatives reproduce \(\Delta p_1\). Location \(B\) sends some mass onward to \(C\), yet gains overall because more arrives from \(A\).

\end{example}

For a group of locations, sum their continuity equations. Internal currents cancel and only currents crossing the group's boundary remain. This is the same cancellation proved on the lattice in Chapter 3, now describing changes of an actual probability distribution.

\begin{example}[Conservation across the boundary of a region]

Treat locations \(A\) and \(B\) as one region of the three-node path. Its stored probability is \(s_n=p_{A,n}+p_{B,n}\). The continuity equations give \[\Delta s_n=-J_{AB}+(J_{AB}-J_{BC})=-J_{BC}.\] Only the current from \(B\) to \(C\) changes the region's total. At \(p_1=(3/4,1/4,0)\), that current is \(1/16\), so the regional mass falls from one to \(15/16\) during the next step.

\end{example}

Over several steps, summation gives \[s_T-s_0=-\sum_{n=0}^{T-1}J_{BC,n}.\] Spatial cancellation removes the internal edge, and temporal cancellation removes intermediate stored values. The final identity combines the two fundamental operations of our calculus. It is a useful model for inventory crossing a facility boundary, traffic entering and leaving a district, or probability entering a selected set of states.

\begin{figure}[htbp]\centering
\resizebox{.98\linewidth}{!}{\begin{tikzpicture}[>=Stealth,every node/.style={font=\small}]
\node[draw=teal,circle,line width=1pt,minimum size=1cm] (A) at (0,0) {$A$};
\node[draw=teal,circle,line width=1pt,minimum size=1cm] (B) at (4.2,0) {$B$};
\node[draw=teal,circle,line width=1pt,minimum size=1cm] (C) at (8.4,0) {$C$};
\draw[->,copper,line width=1.4pt] (A)--node[above=5pt] {$J_{AB}=1/8$}(B);
\draw[->,copper,line width=1.4pt] (B)--node[above=5pt] {$J_{BC}=1/16$}(C);
\node[above=13pt] at (A) {$p_A=3/4$};
\node[above=13pt] at (B) {$p_B=1/4$};
\node[above=13pt] at (C) {$p_C=0$};
\node[below=15pt,align=center] at (A) {$\operatorname{div}J=1/8$\\$\Delta p_A=-1/8$\\$p_A^+=5/8$};
\node[below=15pt,align=center] at (B) {$\operatorname{div}J=-1/16$\\$\Delta p_B=1/16$\\$p_B^+=5/16$};
\node[below=15pt,align=center] at (C) {$\operatorname{div}J=-1/16$\\$\Delta p_C=1/16$\\$p_C^+=1/16$};
\end{tikzpicture}}
\caption{One transport step from (3/4,1/4,0). The middle vertex gains mass even while sending some onward, because its incoming current is larger.}\end{figure}

\section{Stationarity permits motion without accumulation}\label{ch04-stationarity-permits-motion-without-accumulation}

A \textbf{stationary distribution} \(\pi\) satisfies \(K\pi=\pi\). Every coordinate then has zero net change. For the symmetric path, \(\pi=(1/3,1/3,1/3)\) is stationary because every row also sums to one.

Stationarity does not mean that individual objects stay still. At uniform probabilities, equal amounts travel in both directions across each path edge. Net current is zero, although the separate one-way flows are positive.

Stationarity also does not guarantee convergence from every starting state. The cyclic rule \(T(p_A,p_B,p_C)=(p_C,p_A,p_B)\) fixes the uniform distribution but rotates every nonuniform state with period three.

A two-location example makes the possibilities explicit. Let the transition probability from \(A\) to \(B\) be \(a\) and the reverse probability be \(b\), with \(0\leq a,b\leq1\). Then \[K=\begin{pmatrix}1-a&b\\a&1-b\end{pmatrix}.\] Writing \(x=p_A\) gives \[x^+=(1-a)x+b(1-x)=b+(1-a-b)x.\] When \(a+b>0\), the stationary value is \(c=b/(a+b)\). Subtracting it gives \[x^+-c=(1-a-b)(x-c).\] This is precisely the feedback error equation from Chapter 2. If \(0<a+b<2\), errors shrink. If \(a=b=1\), the multiplier is \(-1\) and the two locations alternate. If \(a=b=0\), every state is stationary.

\section{A distribution and an observed count are different objects}\label{ch04-a-distribution-and-an-observed-count-are-different-objects}

Suppose a model predicts \(p=(1/2,1/3,1/6)\). Twelve independent observations have expected counts \((6,4,2)\), but the realized counts could be \((7,4,1)\) or another feasible triple. The empirical proportions are the observed counts divided by twelve.

An \textbf{outcome} is one possible result. An \textbf{event} is a collection of results. Event probabilities are added over their constituent outcomes. For \textbf{independent} observations, probabilities of specified successive outcomes multiply.

If two independent observations each have probabilities \(p_A,p_B,p_C\), the ordered pair \(AC\) has probability \(p_Ap_C\). A probability tree records this multiplication along a path. To find the probability of several paths, add their probabilities.

Independence is a model assumption. Drawing from a finite box with replacement can produce independent draws; drawing without replacement generally does not. From a box with two red tokens and one blue token, \[P(R\text{ then }B)=\frac23\frac13=\frac29\] with replacement, but it is \((2/3)(1/2)=1/3\) without replacement. The second procedure changes the probabilities available at the next draw.

\section{Counting ordered samples with the same composition}\label{ch04-counting-ordered-samples-with-the-same-composition}

For \(N\) independent binary observations with success probability \(p\), a particular ordered sample containing \(k\) successes has probability \(p^k(1-p)^{N-k}\). There are \[\binom Nk=\frac{N!}{k!(N-k)!}\] ways to choose the success positions. Here \(N!=1\cdot2\cdots N\), and \(0!=1\). Consequently \[P(\text{exactly }k\text{ successes})=\binom Nk p^k(1-p)^{N-k}.\] The coefficient counts orders; the powers give the probability of each order.

\begin{example}[One count, several orders]

With four trials and two successes, the six orders are \(SSFF,SFSF,SFFS,FSSF,FSFS,FFSS\). If \(p=1/3\), each has probability \((1/3)^2(2/3)^2=4/81\). The event of exactly two successes has probability \(6(4/81)=8/27\).

\end{example}

For \(k\) categories and counts \(n_1,\ldots,n_k\) adding to \(N\), the number of compatible ordered samples is the \textbf{multinomial coefficient} \[\Omega(n)=\frac{N!}{n_1!\cdots n_k!}.\] Choose the positions of category one, then category two among the remaining positions, and continue. The factorials cancel to this expression.

The count probability is \[\boxed{P(n_1,\ldots,n_k)=\Omega(n)\prod_i p_i^{n_i}.}\] The product sign \(\prod\) means multiply over the indices. A factor with exponent zero contributes one because no occurrence of that category is required.

For \(N=4\), counts \((2,1,1)\) have multiplicity twelve. Under \(p=(1/2,1/3,1/6)\), each compatible ordered sample has probability \(1/72\), so the count event has probability \(1/6\). The empirical distribution \((1/2,1/4,1/4)\) is a lattice point from Chapter 3.

All count probabilities add to one because every ordered sample belongs to exactly one count category. Grouping the samples changes the description, not the total probability.

\begin{example}[All possible count outcomes for two observations]

For two independent observations from \(p=(1/2,1/3,1/6)\), the six possible count states have probabilities:

\begin{longtable}[]{@{}lll@{}}
\toprule\noalign{}
Counts & Multiplicity & Count probability \\
\midrule\noalign{}
\endhead
\bottomrule\noalign{}
\endlastfoot
\((2,0,0)\) & 1 & \(1/4\) \\
\((0,2,0)\) & 1 & \(1/9\) \\
\((0,0,2)\) & 1 & \(1/36\) \\
\((1,1,0)\) & 2 & \(1/3\) \\
\((1,0,1)\) & 2 & \(1/6\) \\
\((0,1,1)\) & 2 & \(1/9\) \\
\end{longtable}

With denominator thirty-six, the numerators are \(9,4,1,12,6,4\), summing to thirty-six. This checks the whole count distribution.

\end{example}

The most probable single count state in this example is \((1,1,0)\), but many other states remain possible. The expected counts are \((1,2/3,1/3)\), which are not themselves a realizable integer count state. An average over possible samples need not be one of those samples.

The chance of observing no \(A\) can be read either from the table or directly. Adding the three no-\(A\) count probabilities gives \(1/4\); requiring two observations from the complement gives \((1-p_A)^2=(1/2)^2\). Agreement between these two calculations checks both the counting and the independence assumption.

\section{Expected values explain what the update predicts}\label{ch04-expected-values-explain-what-the-update-predicts}

A \textbf{random variable} assigns a numerical value to each outcome. Its \textbf{expectation} is the probability-weighted average: \[\mathbb E[X]=\sum_i p_ix_i.\] The \textbf{variance} measures average squared deviation from that mean: \[\operatorname{Var}(X)=\sum_i p_i(x_i-\mathbb E[X])^2
=\mathbb E[X^2]-(\mathbb E[X])^2.\] The second identity follows by expanding the square. Its square root is the standard deviation.

For a success indicator \(I\), the values are one with probability \(p\) and zero otherwise. Hence \(\mathbb E[I]=p\) and \(\operatorname{Var}(I)=p(1-p)\). A success count is the sum of \(N\) such indicators. Expectations add; for independent indicators the variances also add. Thus \[\mathbb E[K]=Np,\qquad \operatorname{Var}(K)=Np(1-p).\] The empirical proportion \(\widehat p=K/N\) satisfies \[\mathbb E[\widehat p]=p,\qquad
\operatorname{Var}(\widehat p)=\frac{p(1-p)}N.\]

For \(N=100\) and \(p=0.2\), the count has mean twenty and standard deviation four. The empirical proportion has mean \(0.2\) and standard deviation \(0.04\). Increasing the sample size fourfold halves this standard deviation. It reduces the typical scale of fluctuations across repeated samples; it does not force every larger sample to be closer.

\begin{figure}[htbp]\centering
\resizebox{.98\linewidth}{!}{\begin{tikzpicture}\begin{axis}[width=12.5cm,height=5.8cm,xlabel={Success count $k$},ylabel={Probability or frequency},xmin=-.5,xmax=10.5,ymin=0,ymax=.35,xtick={0,1,...,10},legend pos=north east,ybar,bar width=5pt]
\addplot[teal,fill=teal!55] coordinates {(0,0.0282475) (1,0.121061) (2,0.233474) (3,0.266828) (4,0.200121) (5,0.102919) (6,0.0367569) (7,0.00900169) (8,0.0014467) (9,0.000137781) (10,5.9049e-06)};\addlegendentry{Exact binomial}
\addplot[copper,fill=copper!55] coordinates {(0,0.03) (1,0.145) (2,0.22) (3,0.255) (4,0.18) (5,0.12) (6,0.04) (7,0.01) (8,0) (9,0) (10,0)};\addlegendentry{200 samples}
\end{axis}\end{tikzpicture}}
\caption{Exact probabilities for ten independent trials with success probability 0.3, compared with frequencies from 200 reproducible simulated samples.}\end{figure}

A finite bound also follows directly from the variance. On the event \(|\widehat p-p|\geq\varepsilon\), the squared deviation is at least \(\varepsilon^2\). Its average is therefore at least \(\varepsilon^2\) times the probability of that event. Rearranging gives \[P(|\widehat p-p|\geq\varepsilon)
\leq\frac{p(1-p)}{N\varepsilon^2}.\] This is an instance of Chebyshev's inequality. It can be a loose bound, but it requires no approximation to the sampling distribution.

Why do expectations add without independence? Expand an average over all joint outcomes and group the contribution of the first variable separately from that of the second. Each variable's marginal probabilities emerge when the other is summed out. Correlations do not affect this linear accounting.

Variances behave differently. Expanding \((X+Y)^2\) introduces the product term \(2XY\). Under independence its average factors as \(2\mathbb E[X]\mathbb E[Y]\), producing variance addition. Without independence, that cancellation need not occur. A batch of observations influenced by the same event can fluctuate together, making the independent-trial formula underestimate their collective variability.

For physical objects taking random transitions, a deterministic matrix update predicts mean counts. If there are twenty objects at \(A\), each independently following the first column of the path matrix, the expected numbers are fifteen at \(A\) and five at \(B\). The actual allocation is random. If there are three objects, the expected numbers are \(9/4\) and \(3/4\); these fractional expectations do not mean an object was divided.

For a fixed matrix, averaging sampled trajectories is compatible with the linear update: conditional expected next proportions are \(Kp_n\). For a nonlinear selection rule, averaging before and after applying the rule need not give the same result. This is another reason to specify which level of the model a probability vector describes.

\section{What can disappear in a finite observation}\label{ch04-what-can-disappear-in-a-finite-observation}

If a category has probability \(p>0\), the chance of missing it in \(N\) independent observations is \((1-p)^N\). With \(p=1/20\) and \(N=20\), this probability is about \(0.3585\). An expected count of one is compatible with observing zero.

This distinction matters for both geometry and information. Exact positive multiplication preserves an existing category. Finite sampling can land on a boundary lattice state where its observed count is zero. Treating an observed zero as an impossible event would be a further modeling choice.

One can combine transport and sampling: compute \(p^+=Kp\), then draw \(N\) observations from \(p^+\). The conditional mean of the sampled proportions is \(p^+\). The realized proportions fluctuate around it and need not obey the deterministic trajectory exactly. This two-stage model separates the transition mechanism from observation noise.

Use the transport view of the \href{https://bobovski66.github.io/IntroDiscreteInfoGeom/demos/operator-explorer.html}{operator explorer} to check that the negative divergence equals each coordinate change. Use the \href{https://bobovski66.github.io/IntroDiscreteInfoGeom/demos/week04.html}{sampling demonstration} to compare repeated empirical distributions.

So far we can describe where probability goes and how observations fluctuate. Chapter 5 asks what uncertainty does during those changes. It turns entropy from a field we can evaluate into a quantity whose differences we can interpret and analyze.

\section{Exercises}\label{ch04-exercises}

\begin{exercises}

\item 

Explain the meaning of every column in the three-location transition matrix. Check nonnegativity and column sums.

\item 

Starting from \(p_0=(0,1,0)\), compute \(p_1,p_2\) and both vector differences.

\item 

For \(p=(1/2,1/3,1/6)\), compute the two path currents, all divergences, and the next state.

\item 

Prove that the general net-current formula gives \(\Delta p=-\operatorname{div}J\).

\item 

For the two-location rule with \(a=1/4\), \(b=1/2\), find the stationary distribution and error multiplier. Calculate two updates from \((1,0)\).

\item 

Explain why the swap rule \(a=b=1\) has a stationary distribution but does not converge from \((1,0)\).

\item 

Calculate the probability of exactly three successes in five independent trials with success probability \(2/3\).

\item 

Find the multiplicity and probability of counts \((1,2,1)\) under \(p=(1/2,1/3,1/6)\).

\item 

For readings \(1,2,5\) with probabilities \(1/4,1/2,1/4\), calculate expectation and variance.

\item 

For \(N=100\) and \(p=0.2\), find the standard deviations of the count and empirical proportion. Repeat for \(N=400\).

\item 

A category has probability \(0.1\). Find the first integer sample size for which its probability of being missed is below \(0.05\).

\item 

Describe a simulation that first transports probability by \(K\) and then draws a finite sample. Explain which quantities are deterministic and which are random.

\end{exercises}

\section{Selected solutions}\label{ch04-selected-solutions}

\textbf{2.} The new states are \((1/4,1/2,1/4)\) and \((5/16,3/8,5/16)\). Their differences from their predecessors are \((1/4,-1/2,1/4)\) and \((1/16,-1/8,1/16)\).

\textbf{3.} Both currents are \(1/24\). Divergence is \((1/24,0,-1/24)\), so the next state is \((11/24,1/3,5/24)\).

\textbf{5.} The stationary state is \((2/3,1/3)\) and the multiplier is \(1-a-b=1/4\). Starting at \((1,0)\) gives \((3/4,1/4)\) and \((11/16,5/16)\).

\textbf{7.} The count probability is \(\binom53(2/3)^3(1/3)^2=80/243\). A single specified order with those counts has probability \(8/243\).

\textbf{8.} The multiplicity is \(4!/(1!2!1!)=12\). Multiplying by \((1/2)(1/3)^2(1/6)\) gives probability \(1/9\).

\textbf{11.} Since \(0.9^{28}\approx0.05233\) and \(0.9^{29}\approx0.04710\), the first qualifying size is twenty-nine.
\labheading\begin{enumerate}
\def\labelenumi{\arabic{enumi}.}
\tightlist
\item
  Use \(K=((3/4,1/4,0),(1/4,1/2,1/4),(0,1/4,3/4))\). (a) Check every column sum. (b) Starting with \(p_0=(1,0,0)\), calculate \(p_1\) and \(p_2\). (c) Write the two vector differences \(\Delta p_0,\Delta p_1\).
\item
  On the path A--B--C, suppose one quarter of the mass at each vertex travels to each adjacent vertex. (a) Find net current on each edge for \(p_0=(1,0,0)\). (b) Compute divergence at each vertex. (c) Verify that minus divergence agrees with \(\Delta p_0\).
\item
  Compare \(p=(1/2,1/3,1/6)\) with a direct count sample of 12 outcomes. (a) Give one integer count triple exactly matching those probabilities. (b) Give another feasible sample count triple. (c) Explain why two samples need not equal their expected counts.
\item
  Find a fixed point for the three-node matrix. (a) Verify the proposed vector sums to one. (b) Multiply by K to check stationarity. (c) Compare with the cyclic permutation of three labels, which can be stationary at uniform while other orbits rotate.
\item
  Design a two-vertex network with a transition probability \(a\) from A to B and \(b\) from B to A. (a) Construct the column-stochastic matrix. (b) Write \(\Delta p_A\) in terms of \(p_A,p_B\). (c) Find a stationary distribution when \(a+b>0\).
\end{enumerate}

\FloatBarrier

\chapter{Entropy and its discrete derivatives}
A sensor that always reports the same symbol is predictable. A sensor that reports four symbols equally often is harder to predict. Both devices have four labels available, but their probability distributions carry different amounts of uncertainty. We need a numerical function that depends on the probabilities themselves.

Entropy provides that function. We met its formula on the probability lattice in Chapter 3. Here we will build its interpretation from surprise, connect it to counting and communication, and then apply the difference calculus to its values. The central question is precise: for a stated change of probability, how much does uncertainty change?

\section{Why surprise is measured with a logarithm}\label{ch05-why-surprise-is-measured-with-a-logarithm}

An event of probability one is unsurprising. An event of probability \(1/2\) requires one fair binary choice to identify. An event of probability \(1/8\) is as selective as three specified fair binary choices. A useful surprise scale assigns these events values zero, one, and three.

For \(z>0\), the base-two logarithm \(\log_2 z\) is the exponent satisfying \(2^{\log_2z}=z\). Thus \[\log_2 8=3,\qquad \log_2(1/8)=-3.\] The \textbf{surprisal} of an event with probability \(p>0\) is \(-\log_2p\) bits. Lower probabilities give greater positive surprise.

Logarithms turn multiplication into addition: \[\log_2(ab)=\log_2a+\log_2b.\] If independent events have probabilities \(p\) and \(q\), the joint probability is \(pq\) and the joint surprisal is the sum of their surprisals. This additivity is the reason a logarithm is useful.

Natural logarithms use base \(e\). The same surprise scale in natural units, or \textbf{nats}, is \(-\ln p\). The conversion is \[\log_2p=\frac{\ln p}{\ln2}.\] Changing base changes units, as converting meters to feet changes the number attached to a length. It does not change which distribution has greater entropy.

\begin{example}[A message's surprise and its frequency]

Suppose a device reports \(A,B,C\) with probabilities \((1/2,1/4,1/4)\). An \(A\) report has surprisal one bit; a \(B\) or \(C\) report has surprisal two bits. An unlikely report is more surprising when it occurs, but it occurs less often. Average uncertainty must account for both facts.

\end{example}

\section{Entropy averages the surprise over observations}\label{ch05-entropy-averages-the-surprise-over-observations}

Using the expectation from Chapter 4, define \textbf{Shannon entropy} \[\boxed{H(p)=-\sum_i p_i\ln p_i.}\] The version in bits is \(H_2(p)=H(p)/\ln2\). In the three-symbol example, \[H_2(p)=\frac12(1)+\frac14(2)+\frac14(2)=\frac32\text{ bits}.\]

Entropy is an average under the distribution being described. It is not the surprise of every individual observation. Nor does it count how meaningful a message is to a person. It quantifies the uncertainty about which label will be observed within a specified probability model.

A zero-probability label contributes zero. It never occurs under the model, so it contributes nothing to the average. We write \(0\ln0:=0\) as a convention for the product and skip such entries in calculations.

Pure states such as \((1,0,0)\) have entropy zero. For a uniform distribution on \(k\) labels, \[H(1/k,\ldots,1/k)=-k\frac1k\ln\frac1k=\ln k.\] More equally plausible possibilities require more information to identify.

Entropy is unchanged by relabeling categories because the sum contains the same probability values. A change of labels is different from merging labels: merging may discard distinctions that were previously observable.

A helpful way to read the entropy sum is to imagine a long list of reports. A report of type \(i\) contributes surprise \(-\ln p_i\). The fraction of reports of that type is modeled by \(p_i\). Multiplying those two quantities gives its contribution to the mean. A rare outcome has high individual surprise but a small averaging weight.

For example, a binary source with probabilities \((0.99,0.01)\) has entropy about \(0.0560\) nats. Its rare outcome has surprisal \(-\ln0.01\approx4.6052\) nats. Those statements are compatible: most reports are easy to predict, while a small fraction are surprising. Confusing individual surprise with average entropy would reverse the interpretation of many information calculations.

One can express entropy as an effective number of equally likely outcomes: \[N_{\mathrm{eff}}=e^{H(p)}.\] A uniform distribution on \(k\) labels has \(N_{\mathrm{eff}}=k\). For \((1/2,1/4,1/4)\), the effective count is \(2^{3/2}\approx2.828\), between two and three. This need not be an integer because it matches an uncertainty level, rather than counting labels literally. The number is sometimes called perplexity in prediction settings.

\section{Why balanced distributions have the largest entropy}\label{ch05-why-balanced-distributions-have-the-largest-entropy}

Nonnegativity follows immediately: \(-\ln p_i\geq0\) when \(0<p_i\leq1\). An upper bound follows from the weighted arithmetic-geometric mean inequality introduced in Chapter 3.

First suppose all \(k\) probabilities are positive. Apply the inequality to numbers \(1/p_i\) with weights \(p_i\): \[\prod_i(1/p_i)^{p_i}\leq\sum_i p_i(1/p_i)=k.\] Taking logarithms gives \[\boxed{0\leq H(p)\leq\ln k.}\] Equality at the upper bound requires the numbers \(1/p_i\) to be equal, hence every probability to be \(1/k\).

If only \(s\) entries are positive, apply the same argument to those entries to obtain \(H(p)\leq\ln s\leq\ln k\). The set of positive-probability labels is the \textbf{support}. Removing available outcomes can lower the largest uncertainty attainable.

For four labels, compare \((1,0,0,0)\), \((1/2,1/4,1/4,0)\), and \((1/4,1/4,1/4,1/4)\). Their entropies in bits are zero, \(3/2\), and two. The bound is attained only by the last distribution.

\section{An entropy table turns the field into a sequence}\label{ch05-an-entropy-table-turns-the-field-into-a-sequence}

Restrict to a binary lattice with four tokens. Let \(k\) be the number assigned to \(A\), and define \[H_k=H(k/4,1-k/4),\qquad k=0,1,2,3,4.\] The result is a five-term sequence indexed by composition.

\begin{longtable}[]{@{}llll@{}}
\toprule\noalign{}
\(k\) & Probability of \(A\) & \(H_k\) in nats & Forward difference \\
\midrule\noalign{}
\endhead
\bottomrule\noalign{}
\endlastfoot
0 & 0 & 0 & \(0.562335\) \\
1 & \(1/4\) & \(0.562335\) & \(0.130812\) \\
2 & \(1/2\) & \(0.693147\) & \(-0.130812\) \\
3 & \(3/4\) & \(0.562335\) & \(-0.562335\) \\
4 & 1 & 0 & \\
\end{longtable}

For example, \[H_1=-\frac14\ln\frac14-\frac34\ln\frac34\approx0.562335.\] The first difference is \(\Delta H_k=H_{k+1}-H_k\). Its sign changes from positive to negative as the count passes through the balanced state. This identifies the maximum from a finite table.

Two different indices will appear in this chapter. The index \(k\) describes position on a lattice. The index \(n\) will describe time along an update trajectory. A spatial difference \(H_{k+1}-H_k\) and a time difference \(H(p_{n+1})-H(p_n)\) use the same subtraction operation, but compare different pairs of states.

\begin{figure}[htbp]\centering
\resizebox{.98\linewidth}{!}{\begin{tikzpicture}
\begin{axis}[width=6.5cm,height=5.3cm,xlabel={Count $k$ of four tokens},ylabel={Entropy $H_k$ (nats)},xmin=-.2,xmax=4.2,ymin=0,ymax=.78,xtick={0,...,4},grid=major]
\addplot[teal,mark=*] coordinates {(0,0)(1,.5623351446)(2,.6931471806)(3,.5623351446)(4,0)};
\end{axis}
\begin{axis}[at={(6.5cm,0)},anchor=south west,width=6.5cm,height=5.3cm,xlabel={Starting count $k$},ylabel={First difference $\Delta H_k$},xmin=-.5,xmax=3.5,ymin=-.7,ymax=.7,xtick={0,...,3},grid=major]
\addplot[copper,ybar,bar width=13pt] coordinates {(0,.5623351446)(1,.130812036)(2,-.130812036)(3,-.5623351446)};
\end{axis}\end{tikzpicture}}
\caption{Binary entropy on the four-token lattice and its forward differences. The differences decrease and change sign at the balanced composition.}\end{figure}

\section{Concavity means diminishing gains in one direction}\label{ch05-concavity-means-diminishing-gains-in-one-direction}

The first differences in the table decrease as \(k\) increases. Taking their differences gives \[\Delta^2H_k=H_{k+2}-2H_{k+1}+H_k.\] For the four-token table, these are approximately \[-0.431523,\quad-0.261624,\quad-0.431523.\] All are negative. Adding successive tokens to \(A\) brings smaller gains, then losses, as the distribution becomes more concentrated in \(A\).

A function on the simplex is \textbf{concave} when \[H((1-t)p+tr)\geq(1-t)H(p)+tH(r),
\qquad 0\leq t\leq1.\] The left side describes pooling the distributions and forgetting which source an observation came from. The right side averages the uncertainty when that source is retained.

For a binary sequence sampled at equal probability intervals, apply concavity to the midpoint: \[H_{k+1}\geq\frac12(H_k+H_{k+2}).\] Rearranging gives \(\Delta^2H_k\leq0\). This connects the geometric statement about mixtures with the discrete second difference.

Here is a finite proof of concavity. It is useful to follow the meaning of the quantities before the algebra. Fix one coordinate, let its two probabilities be \(a,b\), and let \(m=(1-t)a+tb>0\). Define weights \[\alpha=\frac{(1-t)a}{m},\qquad \beta=\frac{tb}{m}.\] They are nonnegative and add to one. Weighted arithmetic-geometric mean applied to \(m/a\) and \(m/b\) gives \[\alpha\ln(m/a)+\beta\ln(m/b)\leq0.\] Indeed, their weighted arithmetic mean is one. Multiplying by \(m\) and rearranging yields \[m\ln m\leq(1-t)a\ln a+tb\ln b.\] If one entry is zero, omit its zero-weight term; the corresponding arithmetic mean is at most one, giving the same inequality. If \(m=0\), every contribution is zero.

Now apply this coordinate inequality to every pair \(p_i,r_i\), add, and multiply by minus one. The result is the stated entropy inequality. No derivative has been used.

The averaging parameter in concavity deserves a concrete interpretation. Suppose source \(P\) is selected with probability \(1-t\) and source \(R\) with probability \(t\). Conditional on the selected source, draw a label using its table. If the source name is hidden, the label has distribution \(m=(1-t)p+tr\).

Knowing the source makes prediction easier on average. For two certain but different sources, \(p=(1,0)\) and \(r=(0,1)\), the conditional uncertainty is zero. Hiding an equally likely source name produces \(m=(1/2,1/2)\) with entropy \(\ln2\). The uncertainty has come entirely from not knowing which source was used.

Concavity does not say a mixture's entropy is at least the entropy of each input. Mixing a fair binary source with a certain source can lower the fair source's entropy while remaining above the weighted average of the two entropies. The two sides of the concavity inequality should be compared exactly as written.

For a denominator-\(N\) lattice, raw directional changes usually become smaller as the step \(1/N\) shrinks. Comparing slopes across resolutions therefore requires deciding whether to use \(\Delta H\) or the quotient \(N\Delta H\). A smaller one-token entropy change on a finer lattice does not by itself show that the underlying function has become flatter. It may simply mean the token represents less probability.

\section{A rule that balances tokens raises entropy}\label{ch05-a-rule-that-balances-tokens-raises-entropy}

Consider a lattice update that moves to a neighbor with strictly greater entropy whenever one exists, and stays put otherwise. On the binary four-token lattice, starting from \(k=0\) gives \[0\to1\to2\to2\to\cdots.\] The corresponding entropy sequence is \[0,\quad0.562335,\quad0.693147,\quad0.693147,\ldots.\] Its time differences are nonnegative. The rule was constructed to have that property.

For several categories, a transfer from a count \(a\) to a count \(b\) raises entropy whenever \(a\geq b+2\). To see why, write \[h(j)=-\frac jN\ln\frac jN.\] The entropy contribution of the two categories changes by \[[h(a-1)-h(a)]+[h(b+1)-h(b)].\] The increments \(h(j+1)-h(j)\) decrease with \(j\), by the concavity just proved. The gain at the smaller count exceeds the loss at the larger one. When \(a=b+1\), the transfer only swaps the counts and leaves entropy unchanged.

Thus entropy is maximized on the token lattice when counts differ by at most one. If \(N\) is divisible by the number of categories, exact uniformity is available. Otherwise the most balanced integer counts replace the continuous uniform point. A lattice constraint changes which maximum can actually be reached.

\begin{example}[The most balanced four-token state]

With three categories and four tokens, exact uniformity would require \(4/3\) tokens in each category. The maximizing count states are the three permutations of \((2,1,1)\). Moving toward \((3,1,0)\) concentrates the counts and lowers entropy. Moving between two permutations of \((2,1,1)\) leaves entropy unchanged.

\end{example}

\section{Entropy change is a property of the dynamics}\label{ch05-entropy-change-is-a-property-of-the-dynamics}

An entropy-increasing rule is one possible update. The growth rule from Chapter 1 can do something else. Starting at \((1/2,1/2)\) and giving \(A\) multiplier two and \(B\) multiplier one produces \((2/3,1/3)\). Its entropy is about \(0.6365\) nats, less than the starting value \(\ln2\approx0.6931\).

The normalized growth rule favors one type; the balancing rule favors uncertainty. A probability vector alone imposes neither direction of change. We need to state both the function being measured and the recurrence being used.

For any trajectory, telescoping gives the exact information balance \[\boxed{\sum_{n=0}^{T-1}[H(p_{n+1})-H(p_n)]
=H(p_T)-H(p_0).}\] Positive and negative steps can cancel. An endpoint entropy change does not reveal every intermediate change.

When a process always raises entropy, the summands measure the uncertainty gained during each step. In Chapter 6 we will use a related quantity that decreases toward a reference distribution, allowing stability to be proved by a one-step inequality.

\section{Prefix codes connect uncertainty to messages}\label{ch05-prefix-codes-connect-uncertainty-to-messages}

Suppose symbols are transmitted as bit strings. A \textbf{prefix code} has the property that no codeword begins another codeword. This allows the receiver to identify where each symbol ends.

For probabilities \((1/2,1/4,1/4)\), assign \[A\mapsto0,\qquad B\mapsto10,\qquad C\mapsto11.\] A binary tree describes the code: each branch contributes a bit, and each leaf identifies a symbol. The expected length is \(3/2\) bits per symbol, matching the entropy.

\begin{figure}[htbp]\centering
\resizebox{.98\linewidth}{!}{\begin{tikzpicture}[level distance=1.5cm,sibling distance=4cm,every node/.style={font=\small},edge from parent/.style={draw,ink,thick},label distance=2pt]
\node[circle,draw=ink,inner sep=3pt]{}
child{node[align=center]{$A:0$\\$p_A=1/2$} edge from parent node[above left]{$0$}}
child{node[circle,draw=ink,inner sep=3pt]{}
child[sibling distance=3cm]{node[align=center]{$B:10$\\$p_B=1/4$} edge from parent node[above left]{$0$}}
child[sibling distance=3cm]{node[align=center]{$C:11$\\$p_C=1/4$} edge from parent node[above right]{$1$}}
edge from parent node[above right]{$1$}};
\node[align=left,anchor=west] at (5,-1.5){Expected length: $3/2$ bits\\Entropy: $3/2$ bits\\Kraft sum: $1$};
\end{tikzpicture}}
\caption{A binary prefix tree with code lengths 1, 2, and 2, matched to probabilities (1/2,1/4,1/4).}\end{figure}

For three equally probable symbols, the same lengths \((1,2,2)\) give an expected length \(5/3\) bits, while the entropy is \(\log_2 3\approx1.5850\). Individual code lengths must be integers. Coding blocks makes finer average lengths possible: nine equally likely two-symbol blocks fit in a fixed four-bit code, or two bits per symbol; twenty-seven three-symbol blocks fit in five bits, or \(5/3\) bits per symbol. For uniform blocks of length \(m\), a fixed code needs \(\lceil m\log_2 3\rceil\) bits, whose per-symbol excess over \(\log_2 3\) is less than \(1/m\).

A tree also imposes a counting restriction. If the longest codeword has length \(L\), a leaf at depth \(\ell_i\) occupies \(2^{L-\ell_i}\) positions at depth \(L\). These sets are disjoint, so \[\sum_i2^{-\ell_i}\leq1.\] This is the Kraft inequality. Chapter 6 combines it with a probability inequality to show that entropy is a lower bound on the expected length of any prefix code.

\section{Joint observations and repeated information}\label{ch05-joint-observations-and-repeated-information}

For a pair of labels \((X,Y)\), let \(p_{ij}\) be the probability that \(X=i\) and \(Y=j\). Its joint entropy is \[H(X,Y)=-\sum_{i,j}p_{ij}\ln p_{ij}.\] The probabilities must add to one over all pairs. This is the same entropy definition applied to a larger label set.

If the observations are independent, \(p_{ij}=p_iq_j\). Expanding the logarithm gives \[\begin{aligned}
H(X,Y)
&=-\sum_{i,j}p_iq_j(\ln p_i+\ln q_j)\\
&=-\sum_i p_i\ln p_i\sum_jq_j
-\sum_jq_j\ln q_j\sum_i p_i\\
&=H(X)+H(Y).
\end{aligned}\] Independent uncertainties add because their probabilities multiply. Two independent fair binary observations have four equally likely ordered outcomes and entropy \(2\ln2\).

Now let \(Y\) be an exact copy of a fair binary \(X\). Only pairs \((0,0)\) and \((1,1)\) occur, each with probability one half. The joint entropy is just \(\ln2\). Each individual observation is fair, but the second adds no uncertainty once the first is known. Marginal probability tables alone therefore do not determine joint entropy.

This example also helps interpret coding blocks. A code for a sequence should reflect the sequence's joint probabilities. Repeatedly copying one symbol produces a very different information source from making independent draws with the same single-symbol frequencies. Counting and coding formulas must retain their independence assumptions.

\section{Counting gives another finite meaning to entropy}\label{ch05-counting-gives-another-finite-meaning-to-entropy}

For counts \(n_i\) adding to \(N\), let \(q_i=n_i/N\) be the empirical distribution. Chapter 4 showed that \[\Omega=\frac{N!}{\prod_i n_i!}\] ordered samples have this composition. Under the distribution \(q\) itself, each such sample has probability \[\prod_i q_i^{n_i}
=\exp\left(N\sum_i q_i\ln q_i\right)=e^{-NH(q)}.\] The probability of the whole count event is therefore \(\Omega e^{-NH(q)}\). An event has probability at most one, so \[\boxed{\frac{\ln\Omega}{N}\leq H(q).}\] The left side is a logarithmic count per observation. This is an exact finite bound, not an equality asserted at every sample size.

For balanced binary counts \((2,2)\), \(\Omega=6\). Hence \(\log_2\Omega/4\approx0.6462\) bits, while the entropy is one bit. The count event includes six of the sixteen equally likely ordered samples. That incomplete coverage accounts for the strict inequality.

The same entropy function thus connects three descriptions: average surprise, attainable coding rates, and the exponential scale of sample multiplicities. Our discrete calculus measures how this function changes between probability states.

Use the simplex view of the \href{https://bobovski66.github.io/IntroDiscreteInfoGeom/demos/operator-explorer.html}{operator explorer} to compare neighboring entropy values. In the \href{https://bobovski66.github.io/IntroDiscreteInfoGeom/demos/week05.html}{weekly demonstration}, explore surprise and code lengths. Two different distributions can still have the same entropy. Chapter 6 develops a comparison that keeps track of a specific reference model.

\section{Exercises}\label{ch05-exercises}

\begin{exercises}

\item 

Find the surprisal in bits of probabilities \(1\), \(1/2\), \(1/8\), and \(1/16\). Explain additivity for two independent events.

\item 

Compute entropy in bits and nats for \((1/2,1/4,1/4)\) and \((3/4,1/4)\).

\item 

Construct the full binary entropy table for \(N=4\), including every first and second difference.

\item 

Distinguish the spatial index \(k\) from the time index \(n\) for the trajectory \(k_n=0,1,2,2,\ldots\). Calculate both kinds of differences.

\item 

Find every maximum-entropy count state for three categories with \(N=5\). Explain why exact uniformity is unavailable.

\item 

Starting from counts \((5,0,0)\), give a sequence of entropy-increasing token transfers ending at one of the states in Exercise 5.

\item 

Starting at \((1/2,1/2)\), calculate three normalized growth states for multipliers \((2,1)\). Compare their entropies.

\item 

Use concavity to show that mixing any distribution with the uniform distribution cannot lower entropy.

\item 

Verify that \(A\mapsto0\), \(B\mapsto10\), \(C\mapsto11\) is a prefix code. Calculate its expected length under uniform probabilities.

\item 

Derive the Kraft inequality by counting available positions in a full binary tree.

\item 

For counts \((3,1)\), calculate \(\Omega\), \(\ln\Omega/N\), and \(H(3/4,1/4)\). Check the finite counting bound.

\item 

Prove that two independent observations from \(p\) have joint entropy \(2H(p)\) by expanding the finite sum over pairs.

\end{exercises}

\section{Selected solutions}\label{ch05-selected-solutions}

\textbf{2.} The first table has \(3/2\) bits or \((3/2)\ln2\approx1.0397\) nats. The second has about \(0.811278\) bits or \(0.562335\) nats.

\textbf{3.} To six decimal places, the entropy table and its differences are \[\begin{aligned}
(H_k)&=(0,\;0.562335,\;0.693147,\;0.562335,\;0),\\
(\Delta H_k)&=(0.562335,\;0.130812,\;-0.130812,\;-0.562335),\\
(\Delta^2H_k)&=(-0.431523,\;-0.261624,\;-0.431523).
\end{aligned}\]

\textbf{5.} The maximizing states are the three permutations of \((2,2,1)\). Their counts differ by at most one. Equal counts would require \(5/3\) tokens per category.

\textbf{7.} Including the initial state, the first three are \((1/2,1/2)\), \((2/3,1/3)\), and \((4/5,1/5)\). Their entropies are approximately \(0.6931,0.6365,0.5004\) nats.

\textbf{9.} No listed word begins another. Expected length is \((1+2+2)/3=5/3\) bits.

\textbf{11.} The multiplicity is four and \(\ln4/4\approx0.346574\). Entropy is about \(0.562335\), so the bound holds strictly.
\labheading\begin{enumerate}
\def\labelenumi{\arabic{enumi}.}
\tightlist
\item
  On a four-token binary lattice, calculate the entropy at \(k=0,1,2,3,4\) using natural logarithms. (a) Convert \(k\) to a probability \(k/4\). (b) Show the five entropy values to four decimals. (c) Explain the symmetry of the table.
\item
  Using your entropy table, (a) compute all four first differences, (b) compute all three second differences, (c) identify the sign change of the first differences, and (d) interpret that change as a maximum.
\item
  Define a rule that moves one token to an adjacent binary state with larger entropy, stopping when no neighbor is better. (a) Start at \(k=0\) and list all visited states. (b) Record entropy changes. (c) Verify that every listed change is nonnegative.
\item
  Choose a probability vector \(p\) with three categories and four tokens. (a) Calculate its entropy. (b) Find all legal neighboring states. (c) Calculate their entropy values. (d) Determine whether \(p\) is a local maximum on the lattice.
\item
  Take the trajectories of an entropy-increasing rule and an unequal-multiplier growth rule. (a) Calculate three consecutive entropy values for each. (b) Compare signs of \(\Delta H\). (c) Explain why entropy production is a statement about the update law, not probability in general.
\end{enumerate}

\FloatBarrier

\chapter{Relative entropy and discrete Lyapunov theory}
Two weather models can assign different chances of rain while having the same entropy. A binary table \((3/4,1/4)\) and its reversal \((1/4,3/4)\) have identical uncertainty, but they favor opposite outcomes. Entropy describes a distribution on its own. Comparing a changing model with a particular target requires another function.

Relative entropy supplies a directed comparison. We will derive its formula from the extra surprise incurred by using one table to describe observations from another. Then we will use its one-step difference to prove that a mixing update approaches its target. This joins information theory to the stability questions from Chapter 2.

\section{Comparing a source distribution with a model}\label{ch06-comparing-a-source-distribution-with-a-model}

Suppose observations follow \(p\), while a proposed model assigns probabilities \(q\). The model reports surprisal \(-\ln q_i\) when label \(i\) occurs. Averaging under the actual source gives the \textbf{cross-entropy} \[H(p,q)=-\sum_i p_i\ln q_i.\] The first argument determines how often outcomes are observed. The second determines the probabilities used to describe them.

Subtract the source's own entropy: \[\begin{aligned}
H(p,q)-H(p)
&=-\sum_i p_i\ln q_i+\sum_i p_i\ln p_i\\
&=\sum_i p_i\ln\frac{p_i}{q_i}.
\end{aligned}\] This is the \textbf{Kullback--Leibler divergence}, also called \textbf{relative entropy}: \[\boxed{D(p\|q)=\sum_i p_i\ln\frac{p_i}{q_i}.}\] The double bar separates the two arguments; it is not a conditional probability symbol. The logarithms are natural, giving units of nats.

The same formula can compare a current state with a reference even when neither is claimed to be the true source. What matters is that the roles of the arguments are explicit.

\begin{example}[The cost of using a fair model]

Let \(p=(3/4,1/4)\) and \(q=(1/2,1/2)\). The fair model assigns surprisal \(\ln2\) to either observation, so its cross-entropy is \(\ln2\). The source entropy is about \(0.562335\). Their difference is \[D(p\|q)=\frac34\ln\frac32+\frac14\ln\frac12
\approx0.130812.\] One summand is negative, but the total is positive. Individual outcomes can favor different models; the average is the relevant comparison.

\end{example}

\section{Direction and support belong to the definition}\label{ch06-direction-and-support-belong-to-the-definition}

Reversing the arguments changes both the averaging weights and the ratios. For the same two tables, \[D(q\|p)=\frac12\ln\frac23+\frac12\ln2
=\frac12\ln\frac43\approx0.143841.\] The two directed values differ. Relative entropy is not a symmetric distance.

If \(p_i=0\), its term contributes zero. If \(p_i>0\) but \(q_i=0\), the divergence is infinite: the reference model assigns zero probability to an outcome that the first table allows.

For example, \[D((1,0)\|(1/4,3/4))=\ln4,\] whereas \(D((1/4,3/4)\|(1,0))\) is infinite. This asymmetry at the boundary will motivate a different, symmetric geometry in Chapter 7.

In data analysis, a zero observed count need not mean a zero underlying probability. One way to avoid assigning impossible probabilities from limited data is to add a positive amount \(\alpha\) to every category count: \[\widetilde p_i=\frac{n_i+\alpha}{N+k\alpha}.\] The entries add to one and are positive. This is a deliberate smoothing rule. It changes the estimator and should be stated, rather than silently replacing zeros inside a calculation.

\section{Why the total divergence cannot be negative}\label{ch06-why-the-total-divergence-cannot-be-negative}

For positive tables, apply weighted arithmetic-geometric mean to \(q_i/p_i\) with weights \(p_i\): \[\prod_i\left(\frac{q_i}{p_i}\right)^{p_i}
\leq\sum_i p_i\frac{q_i}{p_i}=1.\] Take logarithms. The left side becomes \(\sum_i p_i\ln(q_i/p_i)=-D(p\|q)\), so \[\boxed{D(p\|q)\geq0.}\] This is \textbf{Gibbs' inequality}.

Equality requires all ratios \(q_i/p_i\) to be equal. Since both tables sum to one, the common ratio must be one, so \(p=q\).

If some \(p_i\) vanish, apply the inequality only on its support. The arithmetic mean on the right is then \(\sum_{i:p_i>0}q_i\leq1\). Equality in the final divergence bound requires that this sum be one and that the ratios on the support agree. Thus equality still forces \(p=q\), including agreement on zero coordinates. An infinite divergence also satisfies nonnegativity.

This argument explains why the source's own entropy is the smallest average logarithmic surprise available among probability models. It does not say every single observation will favor that model.

The proof of nonnegativity can be understood as an average comparison of likelihood ratios. Some labels have \(p_i/q_i>1\) and favor the first table; others favor the reference. Normalization restricts how these advantages can balance. Gibbs' inequality says their logarithmic average, weighted by \(p\), cannot be negative.

It does not make relative entropy an absolute measure of error in the world. If both tables are poorly chosen descriptions of a measurement process, their divergence still compares those tables. Interpreting it as the cost of using a model assumes that the first table adequately describes the source or empirical frequencies of interest.

For finite data, one may use the empirical table in the first argument and a proposed model in the second. For dynamics, the first argument may instead be a state evolving under an update law. The numerical formula remains unchanged, but the question being answered changes. Naming the roles before substituting numbers prevents that shift from happening silently.

\section{Uniform reference connects relative entropy to entropy}\label{ch06-uniform-reference-connects-relative-entropy-to-entropy}

Let \(u_i=1/k\). Expanding the definition, \[\begin{aligned}
D(p\|u)
&=\sum_i p_i\ln p_i-\sum_i p_i\ln(1/k)\\
&=-H(p)+\ln k.
\end{aligned}\] Hence \[\boxed{D(p\|u)=\ln k-H(p).}\] Relative entropy to uniform is exactly the gap from maximum entropy.

For a nonuniform reference \(q\), a small value of \(D(p\|q)\) means similarity to that particular table, rather than proximity to maximum uncertainty. Approaching a concentrated target may lower entropy while lowering relative entropy as well. The reference determines the question being asked.

\section{A mixing update makes the target explicit}\label{ch06-a-mixing-update-makes-the-target-explicit}

Fix a positive reference distribution \(q\) and a parameter \(0\leq\lambda\leq1\). Define \[p_{n+1}=(1-\lambda)p_n+\lambda q.\] This may describe replacing a fraction \(\lambda\) of a current allocation by a target allocation at every round. Each coordinate is a weighted average, so the update preserves nonnegativity and total probability.

Subtract the target: \[p_{n+1}-q=(1-\lambda)(p_n-q).\] Every coordinate error is multiplied by the same number. Repeating the identity gives \[\boxed{p_n-q=(1-\lambda)^n(p_0-q).}\] Equivalently, \[p_n=(1-\lambda)^n p_0+[1-(1-\lambda)^n]q.\] For \(0<\lambda<1\), the original state has geometrically shrinking weight. At \(\lambda=1\), the target is reached immediately. At \(\lambda=0\), nothing changes.

This is the vector version of stable affine feedback from Chapter 2. There the state was one number; here it is a probability vector. The algebra of the error remains the same.

\begin{example}[Following both the state and its discrepancy]

Take \(p_0=(3/4,1/4)\), \(q=(1/2,1/2)\), and \(\lambda=1/2\). The states and their forward divergences are:

\begin{longtable}[]{@{}lll@{}}
\toprule\noalign{}
Step \(n\) & First probability & \(V_n=D(p_n\|q)\) \\
\midrule\noalign{}
\endhead
\bottomrule\noalign{}
\endlastfoot
0 & \(3/4\) & \(0.130812\) \\
1 & \(5/8\) & \(0.031584\) \\
2 & \(9/16\) & \(0.007833\) \\
3 & \(17/32\) & \(0.001954\) \\
\end{longtable}

The coordinate error halves each step. The divergence falls even faster in this example. It is a nonlinear function of that error, so it need not halve exactly.

\end{example}

Mixing can move toward a nonuniform target while lowering entropy. Take \(p_0=(1/2,1/2)\), \(q=(0.9,0.1)\), and \(\lambda=1/2\). The new state is \((0.7,0.3)\). Entropy falls from about \(0.6931\) to \(0.6109\) nats, while \(D(p_n\|q)\) falls from about \(0.5108\) to \(0.1537\).

There is no conflict between these changes. Entropy asks how uncertain the current state is. Relative entropy to \(q\) asks how different the state is from the specified target. Moving toward a target that favors one outcome can reduce both uncertainty and target discrepancy.

The positivity assumption on \(q\) ensures the forward potential is finite for every starting distribution. If the target assigns zero to a category that the initial state occupies, the forward divergence is infinite until that category is removed entirely. The recurrence may still be well defined, but this particular finite-valued potential no longer describes every step.

\section{From a finite inequality to a stability estimate}\label{ch06-from-a-finite-inequality-to-a-stability-estimate}

The scalar argument in Chapter 5 established convexity of \(f(z)=z\ln z\): \[f((1-\lambda)a+\lambda b)
\leq(1-\lambda)f(a)+\lambda f(b).\] For fixed positive \(q\), rewrite divergence as \[D(p\|q)=\sum_i q_i f(p_i/q_i).\] Each coordinate of the ratio for \(p^+=(1-\lambda)p+\lambda q\) is \[p_i^+/q_i=(1-\lambda)(p_i/q_i)+\lambda.\] Apply convexity with \(a=p_i/q_i\) and \(b=1\). Since \(f(1)=0\), multiply by \(q_i\) and sum to obtain \[\boxed{D(p^+\|q)\leq(1-\lambda)D(p\|q).}\] Zeros in \(p\) use the continuous extension \(f(0)=0\), already reflected in the entropy convention.

Define a scalar sequence \(V_n=D(p_n\|q)\). The inequality becomes \[V_{n+1}\leq(1-\lambda)V_n,\qquad
\Delta V_n\leq-\lambda V_n.\] Iterating gives \(V_n\leq(1-\lambda)^nV_0\). A finite inequality about one step has become a quantitative statement about every later step.

For \(\lambda=1/2\), four steps suffice to reduce this bound by a factor of sixteen. The actual value may fall faster. A bound specifies a guaranteed ceiling, not an exact trajectory.

\section{What a Lyapunov function contributes}\label{ch06-what-a-lyapunov-function-contributes}

A \textbf{Lyapunov function} is a nonnegative potential used to monitor stability. Here \(V(p)=D(p\|q)\) vanishes only at the target and does not increase under the mixing update. The stronger geometric bound shows its value tends to zero.

The state formula already proved convergence for this simple rule. Why also study \(V\)? More complicated systems may lack an explicit formula for \(p_n\). A one-step inequality for a scalar potential can still control every trajectory covered by its assumptions. It replaces a difficult vector calculation by an informative scalar comparison.

A nonincreasing sequence alone need not tend to zero. It might approach a positive number or remain constant. In this example, the factor \(1-\lambda<1\) is what supplies the stronger conclusion. For another recurrence, the sign and size of \(\Delta V\) must be established anew.

The reverse divergence also contracts for mixing when initially finite: \[D(q\|p^+)\leq(1-\lambda)D(q\|p).\] A direct proof uses \(p_i^+\geq p_i^{1-\lambda}q_i^\lambda\), takes logarithms, multiplies by \(q_i\), and adds. Both directions are useful, but they are different sequences. The worked table above uses \(D(p_n\|q)\); the accompanying contraction figure uses \(D(q\|p_n)\).

\begin{figure}[htbp]\centering
\resizebox{.98\linewidth}{!}{\begin{tikzpicture}\begin{axis}[width=12.5cm,height=5.8cm,xlabel={Step $t$},ylabel={Relative entropy (nats)},xmin=0,xmax=12,ymin=0,legend pos=north east]
\addplot[teal,mark=*] coordinates {(0,0.510826) (1,0.223144) (2,0.113137) (3,0.0604663) (4,0.0331085) (5,0.0183531) (6,0.0102406) (7,0.0057345) (8,0.00321756) (9,0.00180733) (10,0.00101582) (11,0.000571144) (12,0.000321188)};\addlegendentry{$D(q\Vert p_t)$}
\addplot[copper,dashed,mark=square*] coordinates {(0,0.510826) (1,0.383119) (2,0.287339) (3,0.215505) (4,0.161628) (5,0.121221) (6,0.090916) (7,0.068187) (8,0.0511402) (9,0.0383552) (10,0.0287664) (11,0.0215748) (12,0.0161811)};\addlegendentry{$(1-\lambda)^tD(q\Vert p_0)$}
\end{axis}\end{tikzpicture}}
\caption{Mixture contraction toward binary target (0.5,0.5) from (0.9,0.1), with mixture fraction 0.25. The proven bound lies above the actual divergence.}\end{figure}

A simple counterexample shows why monotonicity needs care. Let \(q=(1/2,1/2)\) and let an update swap the coordinates of \(p\). Starting at \((3/4,1/4)\) gives the cycle \[(3/4,1/4)\to(1/4,3/4)\to(3/4,1/4).\] The divergence to \(q\) is constant and positive. Thus \(\Delta V=0\) at every step, yet the state does not approach the reference. A nonincreasing potential can coexist with persistent motion along one of its level sets.

For the mixture update, the bound rules out that behavior by forcing a definite proportional decrease whenever \(V>0\). Given a target tolerance \(\varepsilon<V_0\) and \(0<\lambda<1\), it suffices that \[(1-\lambda)^nV_0\leq\varepsilon.\] Taking logarithms gives \[n\geq\frac{\ln(\varepsilon/V_0)}{\ln(1-\lambda)}.\] The denominator is negative; dividing by it reverses the inequality. Rounding the right side up gives a sufficient integer step count. The special case \(\lambda=1\) is handled separately by immediate arrival at the target.

These estimates are useful when an algorithm needs a stopping rule. An observed small change from one step to the next might result from a tiny intensity rather than proximity to the target. A proven bound on the remaining discrepancy answers a more direct question.

\section{Summing dissipation recovers the total loss}\label{ch06-summing-dissipation-recovers-the-total-loss}

Call \[\sigma_n=V_n-V_{n+1}\] the discrepancy dissipated during step \(n\). For the mixing rule it is nonnegative. The discrete fundamental theorem gives \[\boxed{\sum_{n=0}^{T-1}\sigma_n=V_0-V_T.}\] The first three losses in the table are approximately \(0.099228\), \(0.023751\), and \(0.005879\). Their sum is about \(0.128858\), equal to \(V_0-V_3\) to rounding precision.

If \(V_T\to0\), these finite totals approach \(V_0\). Thus an infinite sum of nonnegative one-step losses has a finite total. An \textbf{infinite series} here means the limit of its finite partial sums; the telescoping identity supplies that limit.

The word dissipation describes loss of this particular discrepancy. It does not claim that total probability disappears or that all information of every kind decreases. Probability remains conserved while a scalar comparison to the target changes.

\section{Finite samples give an exact likelihood interpretation}\label{ch06-finite-samples-give-an-exact-likelihood-interpretation}

Suppose an independent sample of size \(N\) has counts \(n_i\) and empirical distribution \(\widehat p_i=n_i/N\). The \textbf{likelihood} of a proposed probability table \(q\) is the probability it assigns to the observed data: \[L(q)=\prod_i q_i^{n_i}\] for one ordered sample. Taking logarithms gives \[\ln L(q)=N\sum_i\widehat p_i\ln q_i.\] Consequently, \[\boxed{\ln L(\widehat p)-\ln L(q)=ND(\widehat p\|q).}\] If we use count probabilities instead, the multinomial coefficients cancel in the likelihood ratio.

Gibbs' inequality proves that the empirical table has at least as large a likelihood as any other table for those observations. This is a statement about fitting the current sample. It does not guarantee the best predictions for unseen observations, particularly when rare categories were missed.

\begin{example}[A four-observation comparison]

For counts \((3,1)\), the empirical model is \((3/4,1/4)\). Its likelihood relative to the fair model is \[\frac{(3/4)^3(1/4)}{(1/2)^4}=\frac{27}{16}.\] The logarithm of this ratio is \(4D((3/4,1/4)\|(1/2,1/2))\). The discrepancy per observation accumulates across the four observations.

\end{example}

For independent batches with specified models, log-likelihood ratios add. This is another instance of logarithms turning repeated multiplicative evidence into accumulated increments.

Likelihood compares how models score observations already recorded. In the four-observation example, both models assign small probabilities to the exact ordered sample; the empirical model's advantage is a ratio of those probabilities. A likelihood value is not a probability that the model itself is true. To assign probabilities to models, one needs an additional model-selection framework, such as the Bayesian reweighting introduced in Chapter 9.

The empirical table maximizes the likelihood of the current sample because its own observed frequencies fit that sample exactly. A category missed by chance receives zero empirical mass, even though it might appear next time. Smoothing keeps such possibilities available at the cost of reducing the current sample's maximum likelihood. The choice reflects a prediction goal extending beyond the fitting identity.

A related support issue distinguishes the two mixture potentials. Starting from \(p_0=(1,0)\) with positive target \(q=(1/2,1/2)\) gives finite \(D(p_0\|q)=\ln2\) but infinite \(D(q\|p_0)\). After any mixture step with \(0<\lambda\leq1\), all coordinates become positive and both values are finite. A contraction bound starting from an infinite reverse value is formally uninformative; one can instead begin its numerical bound at the first positive state.

\section{Mixtures and codes reveal two further consequences}\label{ch06-mixtures-and-codes-reveal-two-further-consequences}

Let \(m=(1-t)p+tr\). Expanding the definitions yields \[H(m)-(1-t)H(p)-tH(r)
=(1-t)D(p\|m)+tD(r\|m).\] The right side is nonnegative. This recovers entropy concavity and quantifies its gap: pooling distinct sources adds uncertainty when their source labels are forgotten. For \(0<t<1\), the gap vanishes exactly when the two distributions agree.

Now consider prefix-code lengths \(\ell_i\) with Kraft sum \(C=\sum_i2^{-\ell_i}\leq1\). Normalize these lengths into probabilities \(r_i=2^{-\ell_i}/C\). Then \[D(p\|r)=-H(p)+(\ln2)\sum_i p_i\ell_i+\ln C.\] Nonnegativity gives \[\sum_i p_i\ell_i\geq H_2(p)-\log_2C\geq H_2(p).\] This completes the coding connection from Chapter 5. A counting condition on code trees and a probability inequality together bound the expected message length.

Use the \href{https://bobovski66.github.io/IntroDiscreteInfoGeom/demos/week06.html}{mixing demonstration} to compare the recorded divergence with its geometric bound. Keep track of the order of its arguments. Then calculate successive losses and verify their telescoping sum.

Relative entropy gives a powerful potential, but its asymmetry and infinite boundary values limit its use as an ordinary distance. Chapter 7 changes coordinates and constructs a symmetric geometry that remains finite at the boundary.

\section{Exercises}\label{ch06-exercises}

\begin{exercises}

\item 

Compute both directed divergences between \((3/4,1/4)\) and \((1/2,1/2)\). Identify which table supplies the averaging weights in each.

\item 

Evaluate \(D((1,0)\|(1/4,3/4))\) and the reverse divergence. Explain the boundary difference.

\item 

Prove the equality condition in Gibbs' inequality when the first table has zero entries.

\item 

Verify \(D(p\|u)=\ln k-H(p)\) and interpret it as an entropy gap.

\item 

Starting at \((1,0,0)\), mix toward uniform with \(\lambda=1/4\). Find \(p_1,p_2\).

\item 

Prove the explicit mixture formula by induction. Explain what happens at \(\lambda=0\) and \(\lambda=1\).

\item 

For the binary worked example, calculate \(V_1/V_0\) and compare it with the guaranteed factor \(1/2\).

\item 

How many steps guarantee a factor-of-sixteen reduction in the contraction bound for \(\lambda=1/2\)? State why this need not be the first step at which the actual value meets that threshold.

\item 

For the four values in the worked table, calculate three dissipation terms and verify their sum.

\item 

Derive the likelihood identity for counts \((3,1)\) directly from the two sample probabilities.

\item 

Apply add-one smoothing to counts \((3,1,0)\) and compare its support with that of the empirical table.

\item 

Derive the entropy mixture-gap identity by expanding both sides. Apply it to an equal mixture of \((1,0)\) and \((0,1)\).

\end{exercises}

\section{Selected solutions}\label{ch06-selected-solutions}

\textbf{1.} The forward value is \((3/4)\ln(3/2)+(1/4)\ln(1/2)\approx0.130812\). The reverse value is \(\tfrac12\ln(4/3)\approx0.143841\).

\textbf{2.} The first divergence is \(\ln4\). The reverse is infinite because the reference assigns zero probability to the second label while the first table assigns it positive probability.

\textbf{5.} The states are \((5/6,1/12,1/12)\) and \((17/24,7/48,7/48)\). Each sums to one.

\textbf{7.} The ratio is about \(0.24145\), below one half. The contraction estimate is a valid upper bound with room to spare.

\textbf{8.} Four steps suffice for the bound because \((1/2)^4=1/16\). The actual discrepancy can decrease more rapidly than the bound.

\textbf{11.} Adding one gives counts \((4,2,1)\) with total seven, hence \((4/7,2/7,1/7)\). All categories now receive positive probability.
\labheading\begin{enumerate}
\def\labelenumi{\arabic{enumi}.}
\tightlist
\item
  Let \(p=(3/4,1/4)\) and \(q=(1/2,1/2)\). (a) Compute \(D(p\|q)\) exactly as a logarithmic expression. (b) Compute \(D(q\|p)\). (c) Explain why the two differ.
\item
  Use the mixture update \(p_{n+1}=\tfrac12p_n+\tfrac12q\), with \(p_0=(3/4,1/4)\). (a) Calculate \(p_1,p_2,p_3\). (b) Check each is a probability vector. (c) Calculate the errors \(p_n-q\).
\item
  Using the sequence from Investigation 2, (a) calculate \(V_n=D(p_n\|q)\) for \(n=0,1,2,3\), (b) compute \(\Delta V_0,\Delta V_1,\Delta V_2\), and (c) test \(V_{n+1}\leq\tfrac12 V_n\) numerically.
\item
  Prove by induction that \(p_n-q=(1-\lambda)^n(p_0-q)\) for a general mixture parameter \(\lambda\). (a) Verify at \(n=0\). (b) Show the step from \(n\) to \(n+1\). (c) Explain why the formula establishes convergence for \(0<\lambda\leq1\).
\item
  Treat \(V_n\) as a sequence of potential values. (a) Write its first three dissipation terms \(V_n-V_{n+1}\). (b) Sum them. (c) Explain why the cancellation identity measures total information discrepancy lost over three updates.
\end{enumerate}

\FloatBarrier

\chapter{Discrete Fisher geometry and statistical distance}
Raising a probability from \(0.49\) to \(0.50\) and raising one from \(0.01\) to \(0.02\) both add one percentage point. The second change doubles the chance of a rare outcome. If probabilities describe measurements or competing models, equal coordinate changes need not represent equally substantial statistical changes.

Chapter 6 compared distributions with a directed divergence. This chapter constructs a symmetric distance. The main idea is simple: take square roots of the probabilities, place the resulting vectors on a sphere, and measure their angular separation. The construction uses finite sums and elementary geometry, so we can calculate exact distances before developing a continuous theory of metrics.

\section{What an ordinary coordinate distance misses}\label{ch07-what-an-ordinary-coordinate-distance-misses}

A Euclidean displacement from \(p\) to \(q\) is \(\delta=q-p\), with squared length \(\sum_i\delta_i^2\). This treats a coordinate change independently of the probability already present in that coordinate.

Sampling provides a reason to look for a different scale. For a binary probability \(x\) and \(N\) independent observations, Chapter 4 gave the standard deviation of the empirical proportion: \[\sqrt{\frac{x(1-x)}N}.\] A small change \(\delta\) in \(x\) can be compared with that fluctuation scale. Its magnitude in these units is \[\frac{\sqrt N\,|\delta|}{\sqrt{x(1-x)}}.\] The denominator depends on the state. The same numerical change can be large relative to one sampling scale and small relative to another.

This motivates a statistical geometry. It does not supply an automatic hypothesis test or a complete measure of practical importance. We are building a distance between probability models, with their sampling behavior as one source of interpretation.

\section{Lengths, angles, and the formulas we need}\label{ch07-lengths-angles-and-the-formulas-we-need}

For a vector \(v=(v_1,\ldots,v_k)\), its Euclidean length is \[\|v\|=\sqrt{\sum_i v_i^2}.\] The dot product of \(v\) and \(w\) is \(v\cdot w=\sum_i v_iw_i\). For nonzero vectors, the angle \(\alpha\) between them satisfies \[\cos\alpha=\frac{v\cdot w}{\|v\|\|w\|}.\]

We measure angles in radians. A full turn is \(2\pi\), a right angle is \(\pi/2\), and an arc subtending angle \(\alpha\) on a circle of radius \(R\) has length \(R\alpha\). A calculator's \(\arccos\) function returns the angle having a specified cosine. Use radian mode.

The sphere of radius two is the set of vectors with squared length four. For three coordinates this is an ordinary sphere in space. With more coordinates, the equation still defines the object and its dot products still make sense, even when a picture is unavailable.

\section{Square roots turn normalization into a sphere equation}\label{ch07-square-roots-turn-normalization-into-a-sphere-equation}

Map a probability vector to \[s(p)=(2\sqrt{p_1},\ldots,2\sqrt{p_k}).\] Its squared length is \[\|s(p)\|^2=\sum_i(2\sqrt{p_i})^2=4\sum_i p_i=4.\] Thus every distribution lies on the radius-two sphere after this transformation. All coordinates are nonnegative, so only that portion of the sphere is used.

The inverse transformation is equally explicit: \[p_i=\frac{s_i^2}{4}.\] A nonnegative sphere vector gives nonnegative probabilities summing to one. We have a new coordinate system for the same probability states.

For two labels the image is a quarter circle. For three labels it is an octant of a sphere. A pure distribution \((1,0,0)\) maps to \((2,0,0)\); the uniform distribution maps to \((2/\sqrt3,2/\sqrt3,2/\sqrt3)\).

The factor two is a normalization choice. It makes the local squared length agree with the categorical Fisher quadratic form derived below. A sphere of radius one would give distances half as large, so the radius belongs in the definition.

\section{Overlap determines an exact statistical distance}\label{ch07-overlap-determines-an-exact-statistical-distance}

The dot product of the two transformed vectors is \[s(p)\cdot s(q)=4\sum_i\sqrt{p_iq_i}.\] Define the \textbf{square-root overlap}, also called the Bhattacharyya coefficient, by \[B(p,q)=\sum_i\sqrt{p_iq_i}.\] Since both vectors have length two, their angle is \(\alpha=\arccos B(p,q)\).

Every term in the overlap is nonnegative. The Cauchy--Schwarz inequality gives \[B(p,q)^2\leq\left(\sum_i p_i\right)\left(\sum_i q_i\right)=1.\] One finite proof of the inequality is the sum-of-squares identity \[\left(\sum_i a_i^2\right)\left(\sum_i b_i^2\right)
-\left(\sum_i a_ib_i\right)^2
=\sum_{i<j}(a_ib_j-a_jb_i)^2\geq0.\] Set \(a_i=\sqrt{p_i}\) and \(b_i=\sqrt{q_i}\). Thus \(0\leq B\leq1\) and \(0\leq\alpha\leq\pi/2\).

Define the categorical \textbf{Fisher--Rao distance} \[\boxed{d_{\mathrm{FR}}(p,q)=2\arccos B(p,q).}\] It is the length of the shorter great-circle arc between the transformed points. The distance lies between zero and \(\pi\).

\begin{example}[A distance that includes a boundary state]

Take \(p=(1,0)\) and \(q=(1/4,3/4)\). The overlap is \(\sqrt{1/4}=1/2\). Their angle is \(\pi/3\), so their distance is \(2\pi/3\).

Both distributions are allowed even though one coordinate of \(p\) is zero. The reverse KL comparison \(D(q\|p)\) is infinite, but the spherical distance is finite.

\end{example}

Identical tables have overlap one and distance zero. Tables with disjoint supports have overlap zero and distance \(\pi\). The construction treats the two arguments symmetrically.

\begin{example}[A three-category comparison carried out numerically]

Compare \(p=(1/2,1/3,1/6)\) with the uniform table \(u=(1/3,1/3,1/3)\). Compute the overlap one coordinate at a time: \[B(p,u)=\sqrt{\frac16}+\frac13+\sqrt{\frac1{18}}
\approx0.977284.\] Its inverse cosine is about \(0.21355\) radians, giving distance about \(0.4271\). The method is the same for three labels as for two: multiply corresponding probabilities, take square roots, add, and apply inverse cosine.

\end{example}

A frequent arithmetic mistake is to take the square root after summing the products. That would calculate \(\sqrt{\sum_i p_iq_i}\), a different expression. Another is to feed a degree-valued angle into a radian arc formula. Recording the overlap and the angle separately makes both mistakes easier to diagnose.

Computer arithmetic can produce an overlap such as \(1.0000000000000002\) for identical distributions. Before applying inverse cosine, clamp such a rounded value to the interval \([0,1]\). A substantial value outside that interval indicates a different issue, such as unnormalized inputs; clamping should not conceal it. For very close binary distributions, the arcsine-coordinate formula is often convenient because it directly compares their arc coordinates.

\section{Why this construction gives a distance}\label{ch07-why-this-construction-gives-a-distance}

On a sphere, the shorter great-circle arc is a shortest route between its endpoints. Our points are never antipodal because their coordinates are nonnegative. Moreover, the shorter arc remains in the nonnegative region: its points can be written as positive combinations of the endpoint sphere vectors, as Chapter 8 will show.

A route from \(p\) to \(q\) followed by one from \(q\) to \(r\) is a possible route from \(p\) to \(r\). It cannot be shorter than the shortest direct spherical route. Therefore \[d_{\mathrm{FR}}(p,r)\leq d_{\mathrm{FR}}(p,q)+d_{\mathrm{FR}}(q,r).\] This is the triangle inequality. Symmetry follows from symmetry of overlap. Zero distance forces identical square-root coordinates, hence identical probabilities.

Together, these properties justify calling the formula a distance. They also let us compare a route through several probability states with the distance between its endpoints.

\section{Binary probabilities have a particularly useful coordinate}\label{ch07-binary-probabilities-have-a-particularly-useful-coordinate}

For \(p=(x,1-x)\), introduce \[\phi=\arcsin\sqrt{x},\qquad 0\leq\phi\leq\pi/2.\] The function \(\arcsin\) returns the angle whose sine is the given value. Because \(\sin^2\phi+\cos^2\phi=1\), the sphere coordinates become \[s(p)=(2\sin\phi,2\cos\phi).\] Moving along the quarter circle is now measured by the angle \(\phi\). The arc coordinate from \(x=0\) is \(2\phi\), giving \[\boxed{d_{\mathrm{FR}}(x,y)
=2|\arcsin\sqrt y-\arcsin\sqrt x|.}\]

\begin{longtable}[]{@{}lll@{}}
\toprule\noalign{}
Probability \(x\) & Angle \(\phi\) & Arc coordinate \(2\phi\) \\
\midrule\noalign{}
\endhead
\bottomrule\noalign{}
\endlastfoot
0 & 0 & 0 \\
\(1/4\) & \(\pi/6\) & \(\pi/3\) \\
\(1/2\) & \(\pi/4\) & \(\pi/2\) \\
\(3/4\) & \(\pi/3\) & \(2\pi/3\) \\
1 & \(\pi/2\) & \(\pi\) \\
\end{longtable}

The equal probability steps \(0\to1/4\) and \(1/4\to1/2\) have distances \(\pi/3\) and \(\pi/6\). The first covers twice as much arc.

\begin{figure}[htbp]\centering
\resizebox{.98\linewidth}{!}{\begin{tikzpicture}
\draw[->,ink] (0,0)--(4.7,0) node[right]{$2\sqrt{x}$};\draw[->,ink] (0,0)--(0,4.7) node[above]{$2\sqrt{1-x}$};
\draw[teal,line width=1.5pt] (4,0) arc[start angle=0,end angle=90,radius=4];
\draw[ink!40] (0,0)--(2,3.464);\draw[ink!40] (0,0)--(2.828,2.828);
\fill[copper] (2,3.464) circle(2.5pt) node[above right]{$x=1/4$};
\fill[teal] (2.828,2.828) circle(2.5pt) node[above right]{$x=1/2$};
\node[below] at (4,0){$x=1$};\node[left] at (0,4){$x=0$};
\node[align=left] at (2,-1){Radius $2$\\Arc coordinate $2\arcsin\sqrt{x}$};
\begin{axis}[at={(6.7cm,0)},anchor=south west,width=6.1cm,height=5.8cm,xlabel={Probability $x$},ylabel={Arc coordinate},xmin=0,xmax=1,ymin=0,ymax=3.2,ytick={0,1.570796,3.141593},yticklabels={$0$,$\pi/2$,$\pi$}]
\addplot[teal,domain=0:1,samples=101]{2*asin(sqrt(x))*pi/180};
\end{axis}\end{tikzpicture}}
\caption{Binary distributions form a quarter circle on the radius-two sphere. Equal probability increments generally cover unequal arc lengths.}\end{figure}

This coordinate also makes ordered binary path lengths easy to calculate. If \(x_0\leq x_1\leq\cdots\leq x_m\), the angle differences all have the same sign. Adding the arc distances telescopes to the distance from \(x_0\) to \(x_m\).

\section{Local Fisher length emerges from a finite identity}\label{ch07-local-fisher-length-emerges-from-a-finite-identity}

We can see the local geometry without differentiating a square root. Let \(q=p+\delta\), where \(\sum_i\delta_i=0\). Rationalizing the difference gives the exact identity \[2\sqrt{p_i+\delta_i}-2\sqrt{p_i}
=\frac{2\delta_i}{\sqrt{p_i+\delta_i}+\sqrt{p_i}}.\] Consequently the squared chord length is \[\|s(q)-s(p)\|^2
=4\sum_i\frac{\delta_i^2}
{(\sqrt{p_i+\delta_i}+\sqrt{p_i})^2}.\]

If every \(p_i\) is positive and the changes are small relative to those probabilities, the denominator is close to \(4p_i\). This gives the approximation \[\boxed{g_p(\delta,\delta)=\sum_i\frac{\delta_i^2}{p_i}.}\] This state-dependent quadratic form is the \textbf{Fisher metric} for categorical distributions. Its square root is a local length estimate.

Why does a chord approximate an arc? If their angular separation is \(\alpha\), the chord has length \(4\sin(\alpha/2)\) and the arc has length \(2\alpha\). For a small angle, these are close. Elementary inequalities \(\sin t\leq t\leq\tan t\) bound their ratio between one and \(1/\cos t\), with \(t=\alpha/2\). The bound approaches one as the angle becomes small.

Thus the local form has two transparent approximations behind it: replacing the square-root denominators by their nearby starting values, and replacing a short arc by its chord. Neither replacement is needed in the exact distance formula.

For a binary step \((x,1-x)\to(x+\delta,1-x-\delta)\), \[g_x(\delta,\delta)=\frac{\delta^2}{x}+\frac{\delta^2}{1-x}
=\frac{\delta^2}{x(1-x)}.\] Its square root is \(|\delta|/\sqrt{x(1-x)}\), precisely the sampling scale considered at the start of the chapter, apart from the factor \(\sqrt N\).

The phrase small relative to the probabilities can be made quantitative. Suppose \[|\delta_i|\leq\varepsilon p_i\quad\text{for every }i,\qquad 0<\varepsilon<1.\] Then \(p_i+\delta_i\) lies between \((1-\varepsilon)p_i\) and \((1+\varepsilon)p_i\). Bounding each square-root denominator in the exact chord identity gives \[\frac{4}{(1+\sqrt{1+\varepsilon})^2}g_p(\delta,\delta)
\leq\|s(q)-s(p)\|^2
\leq\frac{4}{(1+\sqrt{1-\varepsilon})^2}g_p(\delta,\delta).\] Both factors approach one when \(\varepsilon\) becomes small. This is a finite error bound supporting the local approximation. It also shows why a tiny absolute change in a nearly absent category can be problematic: its relative change may still be large.

The form \(g_p\) is positive for every nonzero displacement at a positive state, because it adds nonnegative squares with positive denominators. It scales quadratically: \(g_p(c\delta,c\delta)=c^2g_p(\delta,\delta)\). Its square root therefore scales like a length. These elementary properties explain the term quadratic form without requiring a background in differential geometry.

The allowable displacements also satisfy \(\sum_i\delta_i=0\). A proposed direction that changes every probability upward cannot be a direction through normalized distributions. The local geometry describes this sum-zero space of changes at the current state.

\section{Comparing exact and approximate quantities}\label{ch07-comparing-exact-and-approximate-quantities}

Relative entropy has the same local quadratic scale: \[2D(p+\delta\|p)\approx g_p(\delta,\delta)\] for sufficiently small changes at a positive starting state. This is a local approximation, rather than a formula for arbitrary finite steps.

\begin{example}[A short step at the center]

Take \(p=(1/2,1/2)\) and \(q=(0.55,0.45)\). Then \(\delta=(0.05,-0.05)\) and \[g_p(\delta,\delta)=0.01,\qquad \sqrt{g_p}=0.1.\] The exact distance is about \(0.100167\), and \(2D(q\|p)\approx0.0100167\). The local and exact quantities are close, with the expected squared-length relationship.

\end{example}

A small absolute step need not be small relative to a rare probability. For \(x:0.01\to0.02\), the local estimate is about \(0.100504\) while the exact distance is about \(0.083459\). The probability doubled, so replacing the denominator by its starting value is relatively crude. For \(x:0.49\to0.50\), the local estimate is about \(0.020004\) and the exact distance about \(0.020001\).

At \(p_i=0\), an inward step makes \(\delta_i^2/p_i\) unavailable. This is a limitation of the local expression at the boundary. The exact square-root geometry still gives a finite answer. We therefore use exact distances for boundary edges.

\section{Chords and Hellinger distance provide another scale}\label{ch07-chords-and-hellinger-distance-provide-another-scale}

Expanding the squared chord directly gives \[\|s(p)-s(q)\|^2=8(1-B(p,q)).\] The \textbf{Hellinger distance}, with the convention used here, is \[h(p,q)=\sqrt{1-B(p,q)}
=\frac1{\sqrt2}\sqrt{\sum_i(\sqrt{p_i}-\sqrt{q_i})^2}.\] It is the Euclidean distance in square-root coordinates scaled by \(1/\sqrt2\). It therefore also satisfies the triangle inequality.

The radius-two chord length is \(2\sqrt2\,h\). Hellinger distance and Fisher--Rao distance increase together as overlap decreases, but their numerical scales differ. For two distinct pure states, \(h=1\) while \(d_{\mathrm{FR}}=\pi\).

Independent repeated observations supply a further interpretation of overlap. For product distributions \(p^{\otimes N}\) and \(q^{\otimes N}\), finite multiplication and summation give \[B(p^{\otimes N},q^{\otimes N})=B(p,q)^N.\] For two observations, expand \(\sum_{i,j}\sqrt{p_ip_jq_iq_j}\) and factor it into two identical sums. Repeating gives the formula. If \(p\ne q\), the overlap is less than one and its repeated-observation value shrinks with \(N\). Repeated data can separate models that overlap strongly in a single observation.

The sampling interpretation can now be read more carefully. At \(x=1/2\) and \(N=100\), the standard deviation of an empirical binary proportion is \(0.05\). A change of \(0.01\) is one fifth of that scale. At \(x=0.01\), the corresponding standard deviation is about \(0.00995\), so the same absolute change is about one full fluctuation scale.

This comparison concerns distributions of repeated samples. A single sample can vary for many reasons, and a formal statistical test must specify a null model and its sampling distribution. The geometry identifies a useful state-dependent scale; it does not replace those modeling decisions.

Distances and uncertainties should also be kept separate. Two pure distributions have equal entropy zero but maximal spherical separation \(\pi\). Two distributions near uniform may both have large entropy and still differ from one another. Entropy describes uncertainty within one table. A distance describes separation between two tables. The sphere construction adds a relational measurement that one scalar entropy cannot provide.

\section{Giving the lattice statistical edge lengths}\label{ch07-giving-the-lattice-statistical-edge-lengths}

Return to \(\mathcal P_N\), where an edge represents a one-token transfer. Give each edge the exact weight \[\ell(p,q)=d_{\mathrm{FR}}(p,q).\] The probability lattice is now a \textbf{weighted graph}. One token can represent different statistical lengths depending on the starting composition.

For an interior edge, a local estimate is \(\sqrt{\sum_i(q_i-p_i)^2/p_i}\). This estimate depends on which endpoint is treated as the start, so it is not an exact symmetric edge distance. The spherical formula is symmetric and works at boundaries.

For \(x:1/4\to1/2\), the exact distance is \(\pi/6\approx0.523599\). The local estimate at \(1/4\) is about \(0.577350\). Splitting through \(3/8\) gives local estimates about \(0.288675\) and \(0.258199\), totaling \(0.546874\), closer to the exact value. The exact arc lengths already add to \(\pi/6\) at every such ordered subdivision.

Use the Fisher-distance view of the \href{https://bobovski66.github.io/IntroDiscreteInfoGeom/demos/operator-explorer.html}{operator explorer} to compare these quantities. The \href{https://bobovski66.github.io/IntroDiscreteInfoGeom/demos/week07.html}{weekly demonstration} explores the spherical coordinates. Chapter 8 uses the exact edge lengths to ask which routes are shortest when only certain changes are permitted.

\section{Exercises}\label{ch07-exercises}

\begin{exercises}

\item 

Map \((1,0,0)\), \((1/2,1/2,0)\), and \((1/3,1/3,1/3)\) to radius-two sphere coordinates. Check their lengths.

\item 

Calculate overlap and Fisher--Rao distance between \((1,0)\) and \((1/4,3/4)\).

\item 

Calculate the two distances for \(x:0\to1/4\) and \(x:1/4\to1/2\). Explain the unequal results.

\item 

For \(p=(1/2,1/2)\) and \(q=(0.55,0.45)\), calculate the local squared length, local length, exact distance, and \(2D(q\|p)\).

\item 

Derive the binary Fisher form \(\delta^2/[x(1-x)]\).

\item 

Compare the local estimate for \(x:1/4\to1/2\) with the sum of estimates through \(3/8\). Use each step's starting state.

\item 

Explain why the local formula cannot be used for a move from \((0,1)\) to \((1/4,3/4)\). Calculate the exact distance.

\item 

Prove the chord identity \(\|s(p)-s(q)\|^2=8(1-B)\).

\item 

Find Hellinger and Fisher--Rao distances between two distinct pure states. Compare them with the directed KL values.

\item 

Prove the repeated-observation overlap identity for two observations and then for general \(N\).

\item 

Explain why local estimates from opposite endpoints of a finite edge may differ. Give numerical values for one binary edge.

\item 

On the binary \(N=4\) lattice, add all four exact edge lengths from zero to one. Explain why the answer agrees with the endpoint distance.

\end{exercises}

\section{Selected solutions}\label{ch07-selected-solutions}

\textbf{1.} The vectors are \((2,0,0)\), \((\sqrt2,\sqrt2,0)\), and \((2/\sqrt3,2/\sqrt3,2/\sqrt3)\). Each has squared length four.

\textbf{2.} The overlap is \(1/2\) and the distance is \(2\pi/3\).

\textbf{4.} The squared local length is \(0.01\), its square root is \(0.1\), the exact distance is about \(0.100167\), and twice the divergence is about \(0.0100167\).

\textbf{6.} The single estimate is \(0.25/\sqrt{(1/4)(3/4)}\approx0.577350\). The two estimates total \(0.125/\sqrt{(1/4)(3/4)}+0.125/\sqrt{(3/8)(5/8)}\approx0.546874\).

\textbf{7.} The term for the first coordinate divides by zero. The exact overlap is \(\sqrt3/2\), giving angle \(\pi/6\) and distance \(\pi/3\).

\textbf{12.} The arc coordinates are \(0,\pi/3,\pi/2,2\pi/3,\pi\). Their successive positive differences sum to \(\pi\).
\labheading\begin{enumerate}
\def\labelenumi{\arabic{enumi}.}
\tightlist
\item
  For binary distributions \(p=(1,0)\) and \(q=(1/4,3/4)\), (a) calculate their square-root vectors on the radius-two circle, (b) find the overlap \(\sum_i\sqrt{p_iq_i}\), and (c) calculate the exact Fisher--Rao distance.
\item
  Compare the moves \(x:0\to1/4\) and \(x:1/4\to1/2\). (a) Write both coordinate changes. (b) Compute the arc distances by \(2|\arcsin\sqrt y-\arcsin\sqrt x|\). (c) Explain why equal coordinate steps have unequal lengths.
\item
  Take \(p=(1/2,1/2)\) and \(q=(11/20,9/20)\). (a) Calculate \(\delta=q-p\). (b) Find \(g_p(\delta,\delta)=\sum_i\delta_i^2/p_i\). (c) Calculate \(\sqrt{g_p(\delta,\delta)}\) and the exact spherical distance. (d) Compare numerically.
\item
  Take \(p=(1/4,3/4)\) and move to \(q=(1/2,1/2)\). (a) Compute the local Fisher estimate at p.~(b) Compute the exact arc distance. (c) Break the move into two equal probability steps and add their local estimates. (d) Explain the effect of finer resolution.
\item
  Consider an edge starting at \(p=(0,1)\). (a) Explain why the expression \(\sum_i\delta_i^2/p_i\) is unavailable there for a move into the first category. (b) Compute the exact Fisher--Rao distance from \((0,1)\) to \((1/4,3/4)\). (c) Explain why the exact metric is useful.
\end{enumerate}

\FloatBarrier

\chapter{Path sums, circulation, and statistical geodesics}
A distribution can move from one state to another in many ways. A survey composition can be changed one response at a time. A model can be adjusted by mixing it with a target. A sequence of distributions can also follow equal increments of statistical distance. The endpoints alone do not specify the permitted journey.

Chapter 7 gave us an exact distance between distributions. We now distinguish that endpoint distance from quantities accumulated along a route. The discrete fundamental theorem will identify sums that depend only on endpoints. Positive edge lengths will lead to shortest-path problems, and the sphere construction will give an explicit shortest route when intermediate states are unrestricted.

\section{A path is a finite sequence of legal moves}\label{ch08-a-path-is-a-finite-sequence-of-legal-moves}

Let a graph have probability states as its vertices. Two states are connected when the model permits a direct move between them. A \textbf{path} is a sequence \[\gamma=(v_0,v_1,\ldots,v_m)\] in which every consecutive pair is connected. The number \(m\) counts moves, while \(m+1\) counts listed states.

Assign a signed number \(A(u,v)\) to a directed edge. Require \[A(v,u)=-A(u,v).\] This \textbf{antisymmetry} says reversing a step reverses its contribution. The discrete line integral is the finite sum \[\boxed{\int_\gamma^d A=\sum_{j=0}^{m-1}A(v_j,v_{j+1}).}\] The integral symbol is a name for this operation on edges. Its definition is ordinary finite addition.

Such an edge value can represent a change of potential, a signed work contribution, or another quantity whose sign reverses with direction. A general travel cost need not have this property and will be handled separately.

\section{Gradient integrals recover an endpoint change}\label{ch08-gradient-integrals-recover-an-endpoint-change}

If a scalar potential \(\phi\) is known at vertices, define \[A_\phi(u,v)=\phi(v)-\phi(u).\] Along any path, \[\begin{aligned}
\int_\gamma^d A_\phi
&=[\phi(v_1)-\phi(v_0)]+\cdots+[\phi(v_m)-\phi(v_{m-1})]\\
&=\phi(v_m)-\phi(v_0).
\end{aligned}\] This is the \textbf{discrete fundamental theorem for line integrals}. It is the telescoping identity from Chapter 1 applied to values encountered along a path.

\begin{example}[Different routes, identical potential change]

Suppose \(\phi(U)=0\), \(\phi(V)=2\), \(\phi(W)=5\), and \(\phi(Z)=3\). Along \(U,V,W\), the contributions are \(2\) and \(3\), totaling five. Along \(U,Z,W\), they are \(3\) and \(2\), again totaling five.

The route \(U,W,V,W\) contributes \(5,-3,3\), still totaling five. Backtracking adds movement, but its signed potential changes cancel.

\end{example}

If \(\phi\) is entropy, the integral gives the net entropy change. It does not report the largest entropy encountered or the total size of all fluctuations. Choosing the quantity being accumulated is essential.

Around a closed path, the endpoint is the starting point, so every gradient integral is zero. Chapter 3 used this fact as a test for a gradient field.

\section{Circulation tests whether a potential exists}\label{ch08-circulation-tests-whether-a-potential-exists}

Consider a triangle oriented \(U\to V\to W\to U\), with edge values all equal to one in that direction. Its circulation is three. If these values came from a potential, adding them would telescope to zero. Thus no potential can generate this field.

There is a converse on a connected graph: an antisymmetric edge field is a gradient exactly when every closed-path sum is zero.

To construct the potential, choose a root vertex \(r\) and assign \(\phi(r)=0\). For any other vertex \(v\), choose a path from \(r\) to \(v\) and define \(\phi(v)\) to be its edge sum. This definition must be independent of the chosen path. Two candidate paths followed in opposite directions form a closed route. Its sum is zero by assumption, so the two sums agree.

Finally, follow a path to \(u\) and then the edge \(u\to v\). The construction gives \[\phi(v)=\phi(u)+A(u,v),\] hence \(A(u,v)=\phi(v)-\phi(u)\). This proves the claim.

For a finite graph, one can choose a spanning tree: a connected collection of edges containing every vertex and no cycles. Every remaining edge creates one basic cycle with the tree. Checking those cycle sums suffices, because the tree paths already define a candidate potential. This is a finite way to test the global condition.

There is a useful local-to-boundary version of the circulation calculation. Suppose a patch of triangular lattice faces is assigned a consistent counterclockwise orientation. Add the circulation around every face. An edge shared by two faces is traversed in opposite directions, so its two contributions cancel. Only the outer boundary edges remain: \[\boxed{\text{sum of oriented face circulations}
=\text{circulation around the patch boundary}.}\] This is a discrete Stokes identity. It is another exact cancellation theorem, now relating local loops to a larger loop.

For two triangles sharing an edge, one can check it by writing out both three-edge sums. The shared edge appears once as \(A(u,v)\) and once as \(A(v,u)=-A(u,v)\). The remaining four terms form the boundary route. A larger patch works by the same edge-by-edge accounting.

On a triangle patch without holes, zero circulation on every small face forces zero circulation on its boundary loops. On a graph with a hole, a loop around the hole may not bound a collection of included faces. Local face checks then need not detect its circulation. The earlier criterion used every closed path precisely because it does not assume the graph has a particular filling.

\section{Path length accumulates positive costs}\label{ch08-path-length-accumulates-positive-costs}

Give each unoriented edge a positive weight \(\ell(u,v)=\ell(v,u)\). Define \[L(\gamma)=\sum_{j=0}^{m-1}\ell(v_j,v_{j+1}).\] These contributions stay positive when a step is reversed. A round trip therefore has positive length even when its gradient integral is zero.

\begin{longtable}[]{@{}
  >{\raggedright\arraybackslash}p{(\linewidth - 4\tabcolsep) * \real{0.3333}}
  >{\raggedright\arraybackslash}p{(\linewidth - 4\tabcolsep) * \real{0.3333}}
  >{\raggedright\arraybackslash}p{(\linewidth - 4\tabcolsep) * \real{0.3333}}@{}}
\toprule\noalign{}
\begin{minipage}[b]{\linewidth}\raggedright
Quantity along a path
\end{minipage} & \begin{minipage}[b]{\linewidth}\raggedright
Edge contribution
\end{minipage} & \begin{minipage}[b]{\linewidth}\raggedright
Effect of reversing a step
\end{minipage} \\
\midrule\noalign{}
\endhead
\bottomrule\noalign{}
\endlastfoot
Potential change & \(\phi(v)-\phi(u)\) & Sign reverses \\
Statistical length & \(d_{\mathrm{FR}}(u,v)\) & Length is unchanged \\
Transfer count & 1 per allowed transfer & Another move is counted \\
\end{longtable}

A shortest path minimizes the sum of the chosen costs among permitted paths with the same endpoints. Before solving such a problem, specify both the allowed moves and their weights.

For statistical edge weights, the triangle inequality gives \[L(\gamma)\geq d_{\mathrm{FR}}(v_0,v_m).\] Repeatedly apply the triangle inequality to insert the intermediate states. The endpoint distance is a lower bound on every route's statistical length.

\section{Token distance can be solved by counting surpluses}\label{ch08-token-distance-can-be-solved-by-counting-surpluses}

Let count vectors \(a\) and \(b\) each have total \(N\), and permit a transfer between any two categories. Each transfer decreases one count by one and increases another by one.

The total surplus relative to the target is \[S=\sum_{i:a_i>b_i}(a_i-b_i).\] Since the total counts agree, the total deficit is also \(S\). Thus \[\boxed{d_{\mathrm{token}}(a,b)=S=\frac12\sum_i|a_i-b_i|.}\] At least \(S\) transfers are needed because each transfer can correct at most one surplus token. A route attaining the bound repeatedly sends a token from any surplus category to any deficit category. The donor has a token available, so every move is legal.

For \((4,1,1)\to(1,2,3)\), the differences are \((3,-1,-2)\). Three transfers suffice: one from \(A\) to \(B\) and two from \(A\) to \(C\), in any order. Their statistical lengths depend on the intermediate compositions even though their transfer counts agree.

The formula assumes every category-to-category transfer is allowed. A network restriction can force additional intermediate moves and change the answer.

\section{A lattice restriction can produce a measurable detour}\label{ch08-a-lattice-restriction-can-produce-a-measurable-detour}

Consider the two-token, three-category lattice. Start at \(p=(1,0,0)\) and end at \(r=(0,1/2,1/2)\). Their supports are disjoint, so the direct spherical distance is \(\pi\).

Two legal transfers produce the route \[p\ \longrightarrow\ q=(1/2,1/2,0)\ \longrightarrow\ r.\] The first overlap is \(1/\sqrt2\), giving distance \(\pi/2\). The second overlap is \(1/2\), giving distance \(2\pi/3\). The route length is therefore \[\frac\pi2+\frac{2\pi}3=\frac{7\pi}6.\] Starting instead through \((1/2,0,1/2)\) gives the same length.

These are the only possible first neighbors of \(p\). After either first step, any remaining route has length at least \(2\pi/3\) by the triangle inequality. Hence \(7\pi/6\) is the shortest lattice-path length. Its excess over the unrestricted endpoint distance is \(\pi/6\).

This gap measures a concrete constraint. The unrestricted shortest arc passes through distributions unavailable to a two-token allocation.

\section{Finding weighted shortest paths systematically}\label{ch08-finding-weighted-shortest-paths-systematically}

For a finite graph with nonnegative edge weights, \textbf{Dijkstra's algorithm} computes shortest distances from a starting vertex.

Give the start tentative distance zero and every other vertex tentative distance infinity. At each round, select an unvisited vertex with the smallest tentative distance. For each neighbor, compare its current tentative distance with the selected vertex's distance plus the connecting edge weight. Replace the neighbor's value if this route is shorter, and record its predecessor. Mark the selected vertex visited.

A visited distance can be finalized because an alternative route that reaches it through an unvisited vertex must already have cost at least the smallest remaining tentative distance. Adding nonnegative edge costs cannot reduce that cost. Stop when the destination is finalized, then follow predecessor records backward to recover a shortest route.

\begin{example}[Following the tentative distances]

Suppose edges \(AB,AC,CB,BD,CD\) have weights \(4,1,2,1,6\). Starting at \(A\):

\begin{longtable}[]{@{}llll@{}}
\toprule\noalign{}
Vertex just finalized & Tentative distance to \(B\) & To \(C\) & To \(D\) \\
\midrule\noalign{}
\endhead
\bottomrule\noalign{}
\endlastfoot
\(A\) & 4 & 1 & \(\infty\) \\
\(C\) & 3 & 1 & 7 \\
\(B\) & 3 & 1 & 4 \\
\(D\) & 3 & 1 & 4 \\
\end{longtable}

The predecessor chain gives \(A,C,B,D\), of total length four. Taking the apparently direct edge \(AB\) would have cost four before even reaching \(D\).

\end{example}

On a token lattice, the vertices are count states and the weights can be exact Fisher--Rao distances. The same algorithm applies. A shortest weighted route answers a finite optimization question; it need not use the fewest transfers if a longer sequence of transfers has lower total weight.

The shortest-path algorithm distinguishes a discovered route from a finalized one. A tentative value is an upper bound because an actual path realizing it has been found. It may fall when a better route is discovered. Once the smallest tentative vertex is finalized, nonnegative weights establish that its upper bound is also a lower bound.

Ties do not invalidate the distance calculation. They can be resolved in any fixed order. A predecessor record may then choose one of several equally short paths. If the purpose is to list all shortest paths, more than one predecessor may need to be retained. The demonstrations need only produce one valid minimizing route.

A zero-weight edge can be handled by the same algorithm. Negative weights undermine the finalization argument because a later edge could reduce an accumulated cost. Fisher--Rao weights between distinct distributions are positive, so the probability lattice falls within the intended setting.

For the two-token detour example, the minimizing route has only one intermediate lattice state. Its unrestricted spherical midpoint is \((1/2,1/4,1/4)\), unavailable with two tokens. With four tokens this midpoint is available, but the direct connections to it are still constrained to one-token transfers. Admitting a desired intermediate state and admitting the needed moves are separate requirements.

\section{Two ways to interpolate between distributions}\label{ch08-two-ways-to-interpolate-between-distributions}

An \textbf{interpolation} specifies intermediate states between endpoints. The simplest is the ordinary mixture \[m(t)=(1-t)p+tr,\qquad 0\leq t\leq1.\] Every coordinate is nonnegative and the total remains one. At \(t=1/2\), it gives the coordinate average. Sampling \(t\) at \(0,1/M,\ldots,1\) produces equal coordinate increments.

To obtain equal statistical increments, work on the sphere. Let \[s=2\sqrt p,\qquad u=2\sqrt r,\qquad
\alpha=\arccos B(p,r).\] The square root is taken coordinate by coordinate. For different endpoints, define \[s(t)=\frac{\sin((1-t)\alpha)}{\sin\alpha}s
+\frac{\sin(t\alpha)}{\sin\alpha}u,\qquad
g_i(t)=\frac{s_i(t)^2}{4}.\] For identical endpoints, set \(g(t)=p\) to avoid dividing by zero. This is \textbf{spherical interpolation}.

The coefficients are nonnegative because their angles lie between zero and \(\pi/2\). The formula therefore stays in the nonnegative part of the sphere. To verify its length, choose orthogonal unit vectors \[e_1=s/2,\qquad
e_2=\frac{u/2-(\cos\alpha)e_1}{\sin\alpha}.\] The dot-product formula shows \(e_1\cdot e_2=0\) and \(\|e_2\|=1\). The sine subtraction identity then rewrites the interpolation as \[s(t)=2[\cos(t\alpha)e_1+\sin(t\alpha)e_2].\] Its squared length is four. Thus the recovered probabilities really add to one.

This form also shows that the angle between the states at \(t_1\) and \(t_2\) is \(|t_2-t_1|\alpha\). Consequently \[d_{\mathrm{FR}}(g(t_1),g(t_2))=2|t_2-t_1|\alpha.\]

\section{Equal statistical steps follow a geodesic}\label{ch08-equal-statistical-steps-follow-a-geodesic}

Choose an integer \(M>0\) and list \(g(0),g(1/M),\ldots,g(1)\). Consecutive distances are \(2\alpha/M\), so their sum is \(2\alpha=d_{\mathrm{FR}}(p,r)\). This route attains the endpoint lower bound. A shortest route of this kind is called a \textbf{geodesic}.

The midpoint has a simpler formula: \[\boxed{g_i(1/2)=
\frac{(\sqrt{p_i}+\sqrt{r_i})^2}{2(1+B(p,r))}.}\] To derive it, normalize the vector \(\sqrt p+\sqrt r\). Its squared length is \(2+2B\), precisely the denominator. This produces the point halfway along the shorter sphere arc.

\begin{example}[Two meanings of halfway]

For \(p=(1,0)\) and \(r=(1/4,3/4)\), the coordinate midpoint is \((5/8,3/8)\). Their overlap is \(1/2\), and the spherical midpoint formula gives \((3/4,1/4)\). The latter divides the statistical distance into equal halves.

For three-category endpoints \((0.8,0.1,0.1)\) and \((0.1,0.7,0.2)\), the coordinate midpoint is \((0.45,0.4,0.15)\), while the spherical midpoint is approximately \((0.4339,0.3935,0.1726)\).

\end{example}

\begin{figure}[htbp]\centering
\resizebox{.98\linewidth}{!}{\begin{tikzpicture}\draw[ink,thick] (0,0)--(6,0)--(3,5.196)--cycle;\node[below left] at (0,0){$A$};\node[below right] at (6,0){$B$};\node[above] at (3,5.196){$C$};\draw[copper,line width=1pt,dashed] plot coordinates {(0.9,0.519615) (4.8,1.03923)};\draw[teal,line width=1pt,] plot coordinates {(0.9,0.519615) (0.979372,0.540691) (1.06171,0.561733) (1.14689,0.582706) (1.23476,0.603576) (1.32519,0.624309) (1.41801,0.64487) (1.51308,0.665225) (1.61025,0.685342) (1.70935,0.705187) (1.81022,0.724727) (1.9127,0.743931) (2.01661,0.762767) (2.12178,0.781204) (2.22805,0.799211) (2.33524,0.816759) (2.44316,0.833819) (2.55166,0.850363) (2.66053,0.866365) (2.76962,0.881796) (2.87873,0.896633) (2.98769,0.910851) (3.09632,0.924427) (3.20444,0.937337) (3.31188,0.949561) (3.41844,0.96108) (3.52397,0.971873) (3.62829,0.981923) (3.73123,0.991214) (3.83261,0.99973) (3.93227,1.00746) (4.03005,1.01438) (4.12578,1.0205) (4.21931,1.02579) (4.31049,1.03025) (4.39916,1.03387) (4.48518,1.03664) (4.56841,1.03857) (4.64871,1.03965) (4.72595,1.03987) (4.8,1.03923)};
\node[anchor=west,align=left,text=teal] at (6.6,3.8){Spherical interpolation};
\node[anchor=west,align=left,text=copper] at (6.6,3.1){Ordinary mixture};
\node[anchor=west,align=left] at (6.6,1.9){$p=(0.8,0.1,0.1)$\\$r=(0.1,0.7,0.2)$};
\end{tikzpicture}}
\caption{A spherical geodesic mapped back to the probability triangle, compared with the ordinary mixture segment between the same endpoints.}\end{figure}

In the binary case, both constructions traverse the same quarter-circle interval after mapping to the sphere, but generally at different speeds. With three categories, even the geometric routes generally differ.

For binary probabilities, equal statistical steps can be generated with particularly little machinery. Convert the starting and ending probabilities to angles \[\phi_0=\arcsin\sqrt{x_0},\qquad \phi_1=\arcsin\sqrt{x_1}.\] Choose an integer number of steps \(M\), linearly interpolate the angles, \[\phi_j=\phi_0+\frac jM(\phi_1-\phi_0),\] and convert back with \(x_j=\sin^2\phi_j\). Consecutive distances are \(2|\phi_1-\phi_0|/M\).

For the route from \(x_0=0\) to \(x_1=1\) in four steps, the angles are \[0,\quad\pi/8,\quad\pi/4,\quad3\pi/8,\quad\pi/2.\] The first probabilities are approximately \(0,0.1464,0.5,0.8536,1\). Each statistical step has length \(\pi/4\), although the probability increments are unequal.

An ordinary mixture over the same four steps gives \(0,1/4,1/2,3/4,1\). Its coordinate increments are equal, but its statistical increments are \(\pi/3,\pi/6,\pi/6,\pi/3\). Both routes traverse the same ordered binary arc, so both total lengths are \(\pi\). What differs is where the finite observations are placed along the route.

This example is useful when planning a sequence of model adjustments. Equal changes of a parameter need not mean equally spaced statistical states. A chosen coordinate system determines what ``equal step'' means.

\section{Choosing which shortest route answers the question}\label{ch08-choosing-which-shortest-route-answers-the-question}

A mixture can represent pooling sources. A spherical geodesic gives equal statistical increments. A token path respects indivisible changes. Each construction is appropriate to a different constraint.

Refining a token lattice adds available states, but adjacent denominators need not give nested lattices. Even for nested lattices, a convergence claim about shortest paths needs an argument about the allowed directions and costs. A finer picture alone is not such a proof. The exact endpoint distance remains a valid lower bound at every resolution.

In the \href{https://bobovski66.github.io/IntroDiscreteInfoGeom/demos/week08.html}{path demonstration}, compare the mixture, spherical interpolation, and permitted lattice routes for the same endpoints. Record both the endpoint distance and the sum of edge lengths. For a scalar field such as entropy, also record the signed gradient sum and verify that it depends only on the endpoints.

Chapter 9 uses geometry to choose an update rather than merely measure one. We will balance improvement in a score against the information cost of changing the current distribution.

\section{Exercises}\label{ch08-exercises}

\begin{exercises}

\item 

For potential values \(\phi(U)=0\), \(\phi(V)=2\), \(\phi(W)=5\), compute the integral on \(U,V,W,U\).

\item 

Assign one unit clockwise around a triangle. Compute both routes from \(U\) to \(W\) and explain why no potential produces the field.

\item 

Explain the root-and-path construction of a potential when every closed-path sum vanishes.

\item 

Prove the token-distance formula and construct a shortest transfer sequence from \((4,1,1)\) to \((1,2,3)\).

\item 

Verify the exact shortest length \(7\pi/6\) for the two-token example. Identify its excess over the endpoint distance.

\item 

Run Dijkstra's algorithm on the five-edge example and record the predecessor of each finalized vertex.

\item 

Explain why positive edge lengths cannot be treated as gradient increments. Use a two-edge round trip.

\item 

Verify that the ordinary mixture preserves the probability simplex.

\item 

Derive the spherical midpoint formula by normalizing \(\sqrt p+\sqrt r\).

\item 

Compute coordinate and spherical midpoints between \((1/2,1/2,0)\) and \((1/2,0,1/2)\).

\item 

Divide the binary arc from \((1,0)\) to \((1/4,3/4)\) into four equal statistical steps. Give the three interior first-coordinate values.

\item 

Choose two lattice routes with the same endpoints. Calculate a signed entropy-gradient sum and a positive statistical length for each. Explain what each total measures.

\end{exercises}

\section{Selected solutions}\label{ch08-selected-solutions}

\textbf{1.} The contributions are \(2,3,-5\), totaling zero.

\textbf{2.} The route \(U,V,W\) has sum two. The direct reverse of the clockwise edge \(W\to U\) has value \(-1\), so \(U,W\) has sum minus one. Their difference is the circulation three.

\textbf{4.} There are three surplus tokens at \(A\) and three deficits elsewhere, giving distance three. Transfer once from \(A\) to \(B\) and twice from \(A\) to \(C\).

\textbf{5.} The two possible first steps each cost \(\pi/2\). From either neighbor the direct remaining edge costs \(2\pi/3\), and no route can improve that endpoint distance. The excess is \(\pi/6\).

\textbf{10.} The coordinate midpoint is \((1/2,1/4,1/4)\). Since overlap is \(1/2\), the spherical midpoint is \((2/3,1/6,1/6)\).

\textbf{11.} The angle \(\phi=\arcsin\sqrt{x}\) decreases from \(\pi/2\) to \(\pi/6\) in four increments of \(\pi/12\). The interior first probabilities are \(\sin^2(5\pi/12)\), \(\sin^2(\pi/3)\), and \(\sin^2(\pi/4)\), or \((2+\sqrt3)/4,3/4,1/2\).
\labheading\begin{enumerate}
\def\labelenumi{\arabic{enumi}.}
\tightlist
\item
  Take a path \(A\to B\to C\) with scalar values \(\phi(A)=1\), \(\phi(B)=4\), and \(\phi(C)=2\). (a) Calculate the two gradient edge values. (b) Add them. (c) Compare with \(\phi(C)-\phi(A)\).
\item
  On a directed triangle \(A\to B\to C\to A\), assign each oriented edge value \(1\). (a) Calculate the circulation. (b) Determine whether this edge field could be the gradient of a scalar potential. (c) Justify using telescoping.
\item
  Consider count states \(a=(4,1,1)\) and \(b=(1,2,3)\) with \(N=6\). (a) Compute \(\sum_i|a_i-b_i|\). (b) Find the minimum number of token transfers. (c) Construct an explicit route attaining that minimum.
\item
  Choose two different routes between the same endpoints in the six-token triangle. (a) Calculate a scalar potential such as entropy at every vertex of both routes. (b) Sum the gradient differences. (c) Compare the unsigned cumulative changes \(\sum|\Delta H|\). (d) Explain why the latter can depend on the route.
\item
  Use Fisher--Rao edge weights for two short legal routes between the same probability states. (a) Compute and add the edge distances with a calculator. (b) Find the exact Fisher--Rao distance between endpoints. (c) Check that neither route length is smaller than the endpoint distance.
\end{enumerate}

\FloatBarrier

\chapter{Discrete natural-gradient ideas and learning dynamics}
Suppose an algorithm divides attention among three prediction methods. Each method receives a score after a round of observations. The algorithm should give more attention to better-performing methods, while retaining a valid probability distribution and avoiding an unnecessarily abrupt change.

Earlier chapters introduced all the ingredients for such a rule. Scores specify a preferred direction. Relative entropy measures the cost of departure from the current table. Fisher geometry describes the scale of small changes. We will combine these ideas into an exact finite optimization problem, solve it algebraically, and then explain its connection to natural gradients.

\section{A score specifies what an update should improve}\label{ch09-a-score-specifies-what-an-update-should-improve}

Let \(p=(p_1,\ldots,p_k)\) be current probabilities, initially with every \(p_i>0\). Let \(s_i\) be the score of choice \(i\). If a proposed new table is \(q\), its average score is \[q\cdot s=\sum_i q_i s_i.\] This is a linear function of \(q\). Without any restriction on change, maximizing it would place all probability on a highest-scoring choice.

That may be too abrupt. Scores may be noisy, and the current table may encode accumulated evidence. Introduce a positive intensity \(\eta\) and maximize \[\boxed{J(q)=\eta\,q\cdot s-D(q\|p)}\] over probability vectors \(q\).

The score term rewards improvement; the relative-entropy term charges for changing the current table. Larger \(\eta\) gives scores greater importance relative to that charge. The product \(\eta s_i\) must be dimensionless, so score units and intensity should be chosen consistently.

The direction of KL is part of this objective: the candidate \(q\) supplies the averaging weights, and the old state \(p\) is the reference. Reversing those arguments defines a different problem.

A score can represent a reward or the negative of a loss. If \(s_i\) is an error count, directly favoring larger \(s_i\) would reward worse predictions. To discourage loss \(\ell_i\), use \(s_i=-\ell_i\), giving multipliers \(e^{-\eta\ell_i}\). All the algebra remains valid; the interpretation is determined by which values are desirable.

The current probabilities matter even when scores are fixed. They can represent past evidence, an initial allocation, or the frequency with which choices are used. The update is therefore a response to both the prior state and the new score, rather than a fresh ranking that discards history.

A tie in scores illustrates this clearly. If \(s_i=s_j\), then \(r_i/r_j=p_i/p_j\): their relative weight is unchanged. If every score is equal, the whole table remains unchanged. This provides a quick consistency check on both the formula and a numerical implementation.

The objective also explains the role of the intensity. At very small intensity, the departure penalty dominates and the optimizer remains close to \(p\). At a large intensity, distinctions among scores matter more. Exact probabilities remain nonnegative at every finite intensity, but numerical underflow can make a very small weight appear zero. That computational artifact should be distinguished from the support behavior of the exact recurrence.

\section{Reweight, then normalize}\label{ch09-reweight-then-normalize}

The optimizer will be \[\boxed{r_i=\frac{p_i e^{\eta s_i}}{Z},\qquad
Z=\sum_j p_j e^{\eta s_j}.}\] First multiply each current probability by a positive exponential score factor. Then normalize the resulting weights. This is called an \textbf{exponential update} or a \textbf{multiplicative-weights update}.

Every new entry is positive and the total is one. The rule is the same normalized multiplication used in population selection, now interpreted as a learning step.

\begin{example}[Redistributing attention among three methods]

Take \(p=(1/2,1/3,1/6)\), scores \(s=(0,1,2)\), and \(\eta=\ln2\). The multipliers are \((1,2,4)\). Thus \[Z=\frac12+\frac23+\frac23=\frac{11}6,\qquad
r=\left(\frac3{11},\frac4{11},\frac4{11}\right).\] The third method has the highest score but starts with the smallest probability. Its updated probability ties the second method because its larger multiplier compensates for its smaller starting weight.

\end{example}

\begin{figure}[htbp]\centering
\resizebox{.98\linewidth}{!}{\begin{tikzpicture}\begin{axis}[width=12.5cm,height=5.6cm,ybar,bar width=14pt,xlabel={Choice},ylabel={Probability},xmin=.5,xmax=3.5,ymin=0,ymax=.65,xtick={1,2,3},legend pos=north east]
\addplot[ink,fill=ink!40] coordinates {(1,0.5) (2,0.333333) (3,0.166667)};\addlegendentry{Current table}
\addplot[teal,fill=teal!60] coordinates {(1,0.272727) (2,0.363636) (3,0.363636)};\addlegendentry{Exponential update}
\end{axis}\end{tikzpicture}}
\caption{An exponential update of (1/2,1/3,1/6) with scores (0,1,2) and intensity ln 2 gives (3/11,4/11,4/11).}\end{figure}

Subtracting the old table gives the discrete derivative \[\Delta p=r-p.\] For the example, its coordinates are \((-5/22,1/33,13/66)\), whose sum is zero. The negative first coordinate records attention released to the other choices.

\section{Why this table solves the optimization problem}\label{ch09-why-this-table-solves-the-optimization-problem}

The proof follows by expanding one relative entropy. From the formula for \(r_i\), \[\ln r_i=\ln p_i+\eta s_i-\ln Z.\] Substitute into \(D(q\|r)\): \[\begin{aligned}
D(q\|r)
&=\sum_i q_i(\ln q_i-\ln r_i)\\
&=D(q\|p)-\eta\sum_iq_i s_i+\ln Z.
\end{aligned}\] The last term is \(\ln Z\) because \(\sum_iq_i=1\). Rearranging gives \[\boxed{J(q)=\ln Z-D(q\|r).}\] Gibbs' inequality makes the right side at most \(\ln Z\), with equality exactly when \(q=r\). Therefore \(r\) is the unique optimizer.

The identity does more than identify a best table. It measures the loss from any candidate: \[J(r)-J(q)=D(q\|r).\] A candidate's directed discrepancy from the optimizer is exactly its objective shortfall.

\begin{example}[Comparing an unchanged state with the optimizer]

For \(p=(1/2,1/2)\), scores \((0,1)\), and \(\eta=\ln3\), the update is \(r=(1/4,3/4)\) and \(Z=2\). The maximum objective is \(\ln2\). At the unchanged table, \[J(p)=\frac12\ln3.\] Their difference is \(\ln2-\tfrac12\ln3=\tfrac12\ln(4/3)\), equal to \(D(p\|r)\).

\end{example}

The optimizer ranges over all probability vectors, including boundary candidates. The problem concerns a finite-step update on a finite label set; the set of possible real-valued probability vectors is still continuous. If candidates must belong to a token lattice, the exact optimizer may be unavailable, and the identity can instead compare the permitted candidates.

\section{What the update guarantees about scores}\label{ch09-what-the-update-guarantees-about-scores}

Because \(p\) itself is an available candidate and \(D(p\|p)=0\), \[\eta\,r\cdot s-D(r\|p)\geq\eta\,p\cdot s.\] Hence \[\boxed{\eta(r-p)\cdot s\geq D(r\|p)\geq0.}\] The new table has at least as large an average under the scores used to choose the update. When those scores differ on the positive support, the table changes, the divergence is positive, and the average strictly improves.

This is a statement about a \textbf{fixed score vector during the step}. If scores are recomputed from the new state, the next vector of scores may differ. Improving the old average does not automatically improve an average evaluated under new scores. Cyclic interactions from Chapter 3 are a useful reminder of this distinction.

For the three-method example, the average rises from \(2/3\) to \(12/11\). The one-step improvement is \(14/33\). Repeatedly using the same scores keeps favoring the same methods. A learning problem with new observations at each round generally supplies changing scores.

\section{Odds reveal an exact additive update}\label{ch09-odds-reveal-an-exact-additive-update}

For two positive coordinates, the common normalization cancels: \[\frac{r_i}{r_j}=\frac{p_i}{p_j}e^{\eta(s_i-s_j)}.\] Taking logarithms yields \[\boxed{\ln\frac{r_i}{r_j}-\ln\frac{p_i}{p_j}
=\eta(s_i-s_j).}\] Relative log weights have an exact finite difference equal to the score gap times the intensity.

For binary choices with \(p=(x,1-x)\) and \(d=s_A-s_B\), \[x^+=\frac{xe^{\eta d}}{xe^{\eta d}+1-x}.\] At \(x=1/2\), \(d=\ln2\), and \(\eta=1\), the next value is \(2/3\). The proportion changes by \(1/6\), while the log odds change by \(\ln2\).

This is the odds simplification from Chapter 1 expressed on a logarithmic scale. Repeated multiplication has become accumulated addition.

\section{The geometry determines a local direction}\label{ch09-the-geometry-determines-a-local-direction}

Why is the update connected with a \textbf{natural gradient}? An ordinary gradient chooses steepness using ordinary Euclidean length. A natural gradient uses the statistical length supplied by the Fisher metric. We can derive the direction by completing a square, without differentiating a function.

Consider a small displacement \(\delta\) with \(\sum_i\delta_i=0\). Approximate the change penalty using the Fisher form from Chapter 7: \[D(p+\delta\|p)\approx\frac12\sum_i\frac{\delta_i^2}{p_i}.\] The part of the objective that depends on the displacement is approximately \[Q(\delta)=\eta\sum_i\delta_i s_i-\frac12\sum_i\frac{\delta_i^2}{p_i}.\]

Let \(\bar s=\sum_i p_i s_i\). Since \(\sum_i\delta_i=0\), replacing \(s_i\) by \(s_i-\bar s\) does not change the linear term. Completing the square gives \[\begin{aligned}
Q(\delta)
&=\frac{\eta^2}{2}\sum_i p_i(s_i-\bar s)^2\\
&\quad-\frac12\sum_i
\frac{[\delta_i-\eta p_i(s_i-\bar s)]^2}{p_i}.
\end{aligned}\] The second sum is nonnegative. It is zero exactly at \[\boxed{\delta_i=\eta p_i(s_i-\bar s).}\] These coordinates add to zero, so the proposed displacement satisfies the mass constraint.

Thus \[v_i=p_i(s_i-\bar s)\] is the Fisher natural-gradient direction for a linear score function. The factor \(p_i\) reflects the local statistical geometry. The subtraction of the mean score enforces conservation of total probability.

A large step \(p+\eta v\) may acquire negative entries. This local quadratic calculation therefore does not replace the exact finite update. It explains its small-step direction.

The same direction can be characterized by a fixed statistical movement budget. Suppose \(\sum_i\delta_i=0\) and \(g_p(\delta,\delta)\leq\varepsilon^2\). The score improvement is \[\sum_i\delta_i s_i
=\sum_i\frac{\delta_i}{\sqrt{p_i}}\sqrt{p_i}(s_i-\bar s).\] Cauchy--Schwarz bounds it by \[\varepsilon\sqrt{\sum_i p_i(s_i-\bar s)^2}.\] Equality is attained when \(\delta_i\) is a positive multiple of \(p_i(s_i-\bar s)\), unless all scores are equal. The sum-zero constraint is satisfied by that vector.

Thus the Fisher natural gradient identifies the direction of greatest score increase for a specified small statistical budget. The completed-square calculation chooses a step along that direction by trading gain against a quadratic cost. The exact exponential optimization uses the full KL cost instead of its local quadratic approximation.

This interpretation makes the geometry operational. A rare coordinate is expensive to change by a fixed absolute amount because its squared displacement is divided by a small \(p_i\). The maximizing direction adjusts changes to that cost while also preserving total probability. Steepness is meaningful only after the measure of movement has been chosen.

\section{Comparing the finite update with the local direction}\label{ch09-comparing-the-finite-update-with-the-local-direction}

For small \(\eta\), the exponential satisfies \(e^{\eta s_i}\approx1+\eta s_i\). Its normalizing factor is approximately \(1+\eta\bar s\). Dividing and retaining first-order changes gives \[r_i-p_i\approx\eta p_i(s_i-\bar s).\] The same direction emerges from the exact update and from the local quadratic optimization.

For \(p=(1/2,1/3,1/6)\) and \(s=(0,1,2)\), the mean is \(2/3\) and \[v=\left(-\frac13,\frac19,\frac29\right).\] For a small intensity, the change should be close to \(\eta v\). As intensity grows, the exact normalization accounts for finite effects that this first-order approximation omits.

The relation is also the algebra behind the \textbf{replicator direction} in population dynamics: a type gains share when its score exceeds the population average. The present course studies the exact finite recurrence. The companion calculus course develops the limiting differential equation.

\section{Intensity, support, and stable computation}\label{ch09-intensity-support-and-stable-computation}

Adding a common number \(C\) to all scores multiplies every unnormalized weight by \(e^{\eta C}\), which cancels. Multiplying all scores by a positive number rescales the intensity. Only score differences affect the normalized result.

Let \(s_{\min}\) and \(s_{\max}\) be the smallest and largest scores, and let \(R=s_{\max}-s_{\min}\). Since \(Z\) is a weighted average of the exponentials, \[e^{\eta s_{\min}}\leq Z\leq e^{\eta s_{\max}}.\] It follows that \[e^{-\eta R}\leq\frac{r_i}{p_i}\leq e^{\eta R}.\] This gives a finite bound on relative reweighting. A small product \(\eta R\) limits how dramatically one step can change each positive coordinate.

For computation, subtract \(s_{\max}\) before exponentiating. The update becomes \[r_i=\frac{p_i e^{\eta(s_i-s_{\max})}}
{\sum_jp_j e^{\eta(s_j-s_{\max})}}.\] The largest exponent is now zero, reducing overflow risk without changing the answer.

If a current coordinate is zero, multiplicative updating keeps it zero. The optimization explanation agrees: a candidate assigning positive mass there would have infinite \(D(q\|p)\). To reopen missing choices, one can first mix with a positive reference table. That is an additional operation with a clear interpretation.

A finite update should be evaluated at the scale at which its assumptions apply. If a score summarizes one new batch of observations, using that score repeatedly without receiving new data would count the same evidence more than once in a Bayesian interpretation. In a fixed-payoff population model, repeated use may be exactly what the model intends. The formula alone does not settle this distinction.

It is also useful to separate the score used for an update from the score used to evaluate its eventual performance. The one-step improvement inequality is exact for the supplied vector. It does not guarantee that a prediction method will perform better on the next unseen batch. Such a claim would need assumptions relating future observations to current scores.

When comparing update intensities experimentally, keep the score vector and current distribution fixed at first. Compare \(\Delta p\), the new average score, and \(D(r\|p)\). These measurements expose how a larger response trades greater score improvement against greater movement. After that controlled comparison, allow scores to vary across rounds.

\section{Repeated rounds accumulate evidence}\label{ch09-repeated-rounds-accumulate-evidence}

For a specified score sequence \(s_{i,n}\) and constant intensity, \[p_{i,T}=
\frac{p_{i,0}\exp\left(\eta\sum_{n=0}^{T-1}s_{i,n}\right)}
{\sum_jp_{j,0}\exp\left(\eta\sum_{n=0}^{T-1}s_{j,n}\right)}.\] Each round multiplies the previous weights by another exponential. All intermediate normalizing factors cancel in the final ratio.

If scores are fixed at \(c_i\), this reduces to \[p_{i,T}=\frac{p_{i,0}e^{T\eta c_i}}
{\sum_jp_{j,0}e^{T\eta c_j}}.\] A uniquely highest-scoring available choice approaches probability one. Tied leaders retain their initial ratios while lower-scoring choices fade. Finite positive multipliers preserve positive coordinates at every finite time, even when their limiting values are zero.

If score vectors are supplied independently of the evolving state, their order does not change the final accumulated scores. Intermediate states can still differ. When scores depend on the state, reordering rounds may change the scores themselves, so this order-independence conclusion no longer follows.

\section{Bayesian updating is a particular reweighting rule}\label{ch09-bayesian-updating-is-a-particular-reweighting-rule}

Suppose \(p_i\) is a current probability for hypothesis \(i\), and an observation has likelihood \(L_i>0\) under that hypothesis. Bayes' rule gives \[r_i=\frac{p_iL_i}{\sum_jp_jL_j}.\] Setting \(s_i=\ln L_i\) and \(\eta=1\) turns this into the exponential update. Likelihood is the probability of the observation under a hypothesis; it is not the hypothesis's prior probability.

Start with equal probabilities for three hypotheses. One observation has likelihoods \((0.8,0.4,0.4)\), proportional to \((2,1,1)\). A second conditionally independent observation has likelihoods \((0.25,0.75,0.25)\), proportional to \((1,3,1)\). Multiplying gives weights proportional to \((2,3,1)\), hence posterior probabilities \((1/3,1/2,1/6)\).

The example connects relative growth, accumulated log evidence, and learning within the same normalization formula. The interpretation depends on what the weights and scores represent.

Conditional independence is the assumption that allows the likelihoods of the two observations to multiply once a hypothesis is fixed. Without it, use the second observation's conditional likelihood given both the hypothesis and the first observation. Bayes' rule still reweights, but the table of multipliers may change.

A zero likelihood has a direct meaning: the hypothesis says the observed event is impossible. It removes that hypothesis from the posterior support. Our exponential formula with finite scores uses positive likelihoods because \(\ln0\) is not finite. A model that permits exact zeros should handle them explicitly, rather than substitute an arbitrary very negative score and call the result exact.

The optimization interpretation and the evidence interpretation therefore overlap under stated assumptions. Both normalize positive reweighting factors, while each supplies its own reason for choosing those factors. This is why the recurrence appears in several subjects without making their underlying mechanisms identical.

\section{Seeking uncertainty is another objective}\label{ch09-seeking-uncertainty-is-another-objective}

A high-score update need not increase entropy. If the scores favor one option, the table may become more concentrated. To seek uncertainty deliberately, a simple rule is mixing toward uniform: \[p^+=(1-\lambda)p+\lambda u.\] Concavity and maximal uniform entropy give \[H(p^+)\geq(1-\lambda)H(p)+\lambda H(u)\geq H(p).\] Relative entropy to uniform decreases at the same time because it is the entropy gap.

The \href{https://bobovski66.github.io/IntroDiscreteInfoGeom/demos/week09.html}{learning demonstration} lets you compare finite reweighting steps. Begin with the worked three-method example, record the displacement, and then lower the intensity to compare \((r-p)/\eta\) with the natural-gradient direction.

We can now choose an update by an objective and explain its geometry. Chapter 10 asks what remains observable when distributions move through a network or pass through a noisy recording system.

\section{Exercises}\label{ch09-exercises}

\begin{exercises}

\item 

Compute the exponential update of \((1/4,1/4,1/2)\) with scores \((2,1,0)\) and \(\eta=\ln2\).

\item 

Derive \(J(q)=\ln Z-D(q\|r)\) and identify where normalization of \(q\) enters the proof.

\item 

For \(p=(1/2,1/2)\), scores \((0,1)\), and \(\eta=\ln3\), compare the objectives at \(p\), the optimizer, and \((0,1)\).

\item 

Use optimality to prove \(\eta(r-p)\cdot s\geq D(r\|p)\).

\item 

Derive the exact log-ratio increment and use it for two rounds with binary score gaps \(\ln2\) and \(\ln3\), starting at equal probabilities and \(\eta=1\).

\item 

Calculate the natural-gradient direction for \(p=(1/2,1/3,1/6)\) and scores \((0,1,2)\). Verify its coordinates sum to zero.

\item 

Expand the completed-square formula for \(Q(\delta)\) and explain why its maximizer satisfies the mass constraint.

\item 

Starting from equal probabilities, use fixed scores \((0,1,1)\) and \(\eta=\ln2\). Find the table after four rounds.

\item 

Repeat Exercise 8 from \((1/2,1/3,1/6)\) and explain the ratio of the tied leaders.

\item 

Show that shifting scores \((1000,1001,999)\) by their largest entry leaves the normalized answer unchanged.

\item 

Apply the two Bayesian likelihood updates in the worked example. State the independence assumption needed to multiply their likelihoods.

\item 

Compare entropy change under a score-favoring update and under mixing toward uniform. Explain why the update objective matters.

\end{exercises}

\section{Selected solutions}\label{ch09-selected-solutions}

\textbf{1.} The unnormalized weights are \((1,1/2,1/2)\), giving \((1/2,1/4,1/4)\).

\textbf{3.} The three objective values are \(\tfrac12\ln3\), \(\ln2\), and \(\ln(3/2)\), respectively. The optimizer has the largest value.

\textbf{5.} The odds are multiplied first by two and then by three. Final odds are six, giving probabilities \((6/7,1/7)\). The accumulated log-odds change is \(\ln6\).

\textbf{6.} The mean score is \(2/3\), giving \((-1/3,1/9,2/9)\). Their sum is zero.

\textbf{8.} The weights are proportional to \((1,16,16)\), so the table is \((1/33,16/33,16/33)\).

\textbf{9.} Unnormalized weights are \((1/2,16/3,8/3)\), totaling \(17/2\). The table is \((1/17,32/51,16/51)\); the tied leaders retain their original ratio two to one.
\labheading\begin{enumerate}
\def\labelenumi{\arabic{enumi}.}
\tightlist
\item
  Let \(p=(1/2,1/3,1/6)\), scores \(s=(0,1,2)\), and intensity \(\eta=\ln2\). (a) Find the three unnormalized weights. (b) Normalize. (c) Compute each coordinate of \(\Delta p\). (d) Verify that the three differences sum to zero.
\item
  Take \(p=(1/2,1/2)\), score difference \(d=\ln2\), and \(\eta=1\). (a) Calculate the next A probability. (b) Calculate \(\Delta x\). (c) Find the initial and new odds. (d) Verify \(\Delta\log O=\eta d\).
\item
  Using your three-category example, compute \(J(q)=\eta q\cdot s-D(q\|p)\) for the original p and for the exponential optimizer r. (a) Evaluate both. (b) Find \(Z\). (c) Check \(J(r)=\ln Z\).
\item
  Repeat the three-category update for \(\eta=0.1\) and \(\eta=0.5\). (a) Compute the actual coordinate changes. (b) Find \(\eta p_i(s_i-\bar s)\) at each intensity. (c) Compare approximation errors and explain why the smaller step agrees better.
\item
  Design scores for which one category begins rare but becomes more likely. (a) Choose a valid positive starting table. (b) Calculate four updates. (c) List the probability differences and the score-weighted means. (d) Explain what each number reveals about the dynamics.
\end{enumerate}

\FloatBarrier

\chapter{Graph divergence, diffusion, and information flow}
A network can transport something, record something, or do both. A random walker changes location on a graph. A sensor converts a detailed physical state into an imperfect report. An observer may see only a group label rather than the original category. Each operation transforms a probability distribution.

The tools developed throughout the course now meet in one setting. Edge differences produce currents; divergence gives conservation; entropy measures uncertainty; relative entropy and spherical distance measure distinctions between models. We will examine which quantities a network preserves, which it changes, and what becomes impossible to recover after observations are grouped.

\section{From an edge gradient to diffusion}\label{ch10-from-an-edge-gradient-to-diffusion}

Let a finite undirected graph have vertices \(i\) and symmetric nonnegative edge weights \(c_{ij}=c_{ji}\). A positive weight means that the edge permits transport. Let \(p_i\) be the probability stored at vertex \(i\).

The graph gradient on an oriented edge is \[(\nabla_G p)_{ij}=p_j-p_i.\] It is positive when the probability is higher at the destination. Diffusion moves down this difference, from larger concentration to smaller concentration. With step parameter \(\tau\geq0\), define the outward current \[J_{ij}=\tau c_{ij}(p_i-p_j)=-\tau c_{ij}(\nabla_Gp)_{ij}.\] The current reverses sign when the edge orientation is reversed.

The \textbf{graph Laplacian} is the vertex operator \[\boxed{(Lp)_i=\sum_{j\sim i}c_{ij}(p_i-p_j).}\] It adds the weighted differences between the value at a vertex and the values at its neighbors. With our outward-divergence convention, \(\operatorname{div}J=\tau Lp\). The continuity equation becomes \[\boxed{\Delta p=-\tau Lp,\qquad p^+=(I-\tau L)p.}\]

This is a complete discrete diffusion equation. The time difference on the left records one update; the graph operator on the right measures imbalance across adjacent locations.

\section{A step-size condition preserves probabilities}\label{ch10-a-step-size-condition-preserves-probabilities}

Let \(d_i=\sum_{j\sim i}c_{ij}\) be the weighted degree of vertex \(i\). Expanding the update gives \[p_i^+=(1-\tau d_i)p_i+\tau\sum_{j\sim i}c_{ij}p_j.\] A sufficient condition for all coefficients to be nonnegative is \[0\leq\tau d_i\leq1\quad\text{for every vertex }i.\] When the graph has a positive edge, this is ensured by \[\tau\leq\frac1{\max_i d_i}.\] On a graph with no edges, \(L=0\) and nothing changes.

Symmetry of the edge weights makes \(K=I-\tau L\) symmetric. Its rows and columns each sum to one. Under the step-size condition it is therefore \textbf{doubly stochastic}: both a valid transition matrix and a table whose row sums are one.

The algebra explains why excessive steps are problematic. If \(\tau d_i>1\), the retained-mass coefficient at vertex \(i\) is negative. Starting with all mass at that vertex can then produce a negative probability. A time step belongs to the model, just as correction strength did in Chapter 2.

\section{The three-node path gives an exact example}\label{ch10-the-three-node-path-gives-an-exact-example}

For the path \(A\)--\(B\)--\(C\) with unit edge weights, \[L=\begin{pmatrix}1&-1&0\\-1&2&-1\\0&-1&1\end{pmatrix}.\] The largest degree is two, so \(0\leq\tau\leq1/2\) preserves nonnegativity. Taking \(\tau=1/4\) recovers the matrix from Chapter 4: \[K=\begin{pmatrix}3/4&1/4&0\\1/4&1/2&1/4\\0&1/4&3/4\end{pmatrix}.\]

Starting from \((1,0,0)\) gives first \((3/4,1/4,0)\) and then \((5/8,5/16,1/16)\). The uniform vector \(u=(1/3,1/3,1/3)\) is stationary.

We can solve every trajectory by splitting deviations into two patterns: \[a(1,0,-1)+b(1,-2,1).\] Both patterns have coordinate sum zero. Every probability vector can be written \[p_0=u+a(1,0,-1)+b(1,-2,1),\] where \[a=\frac{p_{A,0}-p_{C,0}}2,\qquad
b=\frac{p_{A,0}+p_{C,0}-2p_{B,0}}6.\]

Direct matrix multiplication shows \[K(1,0,-1)=\frac34(1,0,-1),\qquad
K(1,-2,1)=\frac14(1,-2,1).\] The first imbalance shrinks by three quarters each step; the second shrinks by one quarter. Linearity then gives \[\boxed{
p_n=u+a(3/4)^n(1,0,-1)+b(1/4)^n(1,-2,1).
}\] Both powers tend to zero, proving convergence to uniform from every initial distribution.

For \(p_0=(1,0,0)\), the coefficients are \(a=1/2\) and \(b=1/6\). The slower pattern measures imbalance between the two ends of the path. This decomposition shows how a network can have several rates of relaxation without requiring us to solve a new recurrence at each vertex.

\begin{example}[A centered initial state relaxes at the faster rate]

Start the three-node walk at \(p_0=(0,1,0)\). Its end imbalance coefficient is \(a=0\), while \(b=-1/3\). Therefore \[p_n=\left(\frac13-\frac13(1/4)^n,\;
\frac13+\frac23(1/4)^n,\;
\frac13-\frac13(1/4)^n\right).\] The endpoint probabilities remain equal by symmetry. Only the faster pattern is present. Starting at one endpoint excites both patterns and retains the slower factor \((3/4)^n\).

\end{example}

A network's relaxation rate depends both on its transition rule and on which patterns the initial state contains. The exact formula makes that dependence visible. A plot of one initial state may miss a slowly decaying pattern that another initial state excites.

For this example, the error at the middle vertex is \((2/3)(1/4)^n\). Requiring it to be below \(0.01\) reduces to a comparison of powers, just like the scalar stopping questions in Chapter 2. The network problem has been decomposed into scalar difference equations.

These patterns are often called modes. A mode is a vector pattern reproduced up to a multiplier by the update. No new notation is needed to verify one: multiply the vector by the matrix and compare its coordinates with the proposed multiple.

\section{Entropy rises under a doubly stochastic update}\label{ch10-entropy-rises-under-a-doubly-stochastic-update}

Let \(h(x)=-x\ln x\), with \(h(0)=0\). Chapter 5 proved its concavity. If \(K\) is doubly stochastic, each output coordinate is a weighted average of input coordinates because its row sums to one: \[h\left(\sum_jK_{ij}p_j\right)\geq\sum_jK_{ij}h(p_j).\] Summing over output coordinates and using the column sums gives \[\boxed{H(Kp)\geq H(p).}\] Thus the nonnegative entropy increment \[\sigma_n=H(p_{n+1})-H(p_n)\] is guaranteed for this diffusion rule.

For the three-node example, \(H(p_0)=0\), \(H(p_1)\approx0.562335\), and the exact trajectory approaches entropy \(\ln3\). Telescoping yields \[\sum_{n=0}^{T-1}\sigma_n=H(p_T)-H(p_0)\leq\ln3-H(p_0).\] The total possible entropy gain is bounded.

Doubly stochastic does not by itself imply convergence. A permutation matrix is doubly stochastic and preserves entropy exactly. The two-vertex diffusion rule at \(\tau=1\) simply swaps its coordinates. A convergence claim needs information about the dynamics beyond entropy monotonicity.

If the graph is disconnected, mass cannot travel between components. Its stationary behavior is constrained by the mass initially present in each component. Global uniformity is not generally reachable from every initial state.

\section{A general observation channel has another interpretation}\label{ch10-a-general-observation-channel-has-another-interpretation}

A \textbf{channel} is a column-stochastic matrix \(K\) whose entry \(K_{yi}\) is the probability of reporting output \(y\) given input label \(i\). The output and input label sets may have different sizes. An input distribution \(p\) produces output distribution \[P_y=\sum_iK_{yi}p_i.\] This is the same matrix arithmetic as transport, now describing observation.

A deterministic grouping is a special channel. If labels one and two are reported as group \(A\), while labels three and four are reported as group \(B\), then \[K=\begin{pmatrix}1&1&0&0\\0&0&1&1\end{pmatrix}.\] The output retains group totals and discards the original label within each group.

An arbitrary channel need not increase entropy. The two-input channel that always reports the same single symbol sends every input to a certain output of entropy zero. Even a square channel can do this, for example \[K=\begin{pmatrix}1&1\\0&0\end{pmatrix}.\] It is column-stochastic but not doubly stochastic.

There is, however, a general statement about distinguishing two input models: passing both through the same channel cannot increase their KL divergence or their Fisher--Rao distance. The reference to the \textbf{same} channel is essential.

The observational interpretation differs from the physical transport interpretation in what is claimed to happen to an individual state. A grouping report does not require the original categories to merge physically. It records only a coarser label. The fine state can remain as detailed as before while the observer loses access to some of its distinctions.

This distinction matters when reading data-processing inequalities. They compare the probability tables available before and after the chosen recording operation. They do not claim that every internal feature of the system has ceased to exist. A more informative measurement scheme may retain distinctions that another scheme erases.

Observation rules can be composed. If \(K\) converts a fine label to an intermediate report and \(R\) converts that report to a final report, the combined channel is \(RK\). Its entries are nonnegative, and its column sums are one because \[\sum_z(RK)_{zi}
=\sum_yK_{yi}\sum_zR_{zy}
=\sum_yK_{yi}=1.\] Two stages of processing therefore obey the same framework as one. Data processing can be applied successively to compare the distinctions visible at each stage.

\section{The log-sum inequality organizes the proof}\label{ch10-the-log-sum-inequality-organizes-the-proof}

For nonnegative lists \(a_i,b_i\), let \(A=\sum_i a_i\) and \(B=\sum_i b_i\). When the relevant sums are positive and finite, define normalized tables \(\alpha_i=a_i/A\) and \(\beta_i=b_i/B\). Expanding gives \[\sum_i a_i\ln\frac{a_i}{b_i}
=A\ln\frac AB+A D(\alpha\|\beta).\] Gibbs' inequality therefore implies \[\boxed{A\ln(A/B)\leq\sum_i a_i\ln(a_i/b_i).}\] This is the \textbf{log-sum inequality}. It says grouping terms can reduce a directed logarithmic discrepancy. The established zero and infinity conventions extend it to boundary cases.

Apply the inequality separately to each output label, with \[a_i=K_{yi}p_i,\qquad b_i=K_{yi}q_i.\] Their sums are \((Kp)_y\) and \((Kq)_y\). For nonzero channel entries, the factor \(K_{yi}\) cancels inside the ratio. Summing over outputs gives \[\begin{aligned}
D(Kp\|Kq)
&\leq\sum_{y,i}K_{yi}p_i\ln\frac{p_i}{q_i}\\
&=\sum_i p_i\ln\frac{p_i}{q_i}\sum_yK_{yi}\\
&=D(p\|q).
\end{aligned}\] This is the \textbf{data-processing inequality}.

The proof uses only finite sums and Gibbs' inequality. Its meaning is observational: a common reporting mechanism cannot make two models more distinguishable by this measure. It can preserve a distinction or erase some of it.

\section{Relative entropy to stationarity cannot increase}\label{ch10-relative-entropy-to-stationarity-cannot-increase}

If \(\pi\) is a stationary distribution, then \(K\pi=\pi\). Data processing gives \[D(Kp\|\pi)=D(Kp\|K\pi)\leq D(p\|\pi).\] Thus \(D(p_n\|\pi)\) is a nonincreasing potential whenever its initial value is finite.

This generalizes the Lyapunov viewpoint from Chapter 6 to Markov transport. It does not automatically provide the geometric contraction factor proved for the mixture rule. Some channels preserve the divergence. A separate estimate is needed to obtain a numerical convergence rate.

When \(\pi\) is uniform, the relation \(D(p\|\pi)=\ln k-H(p)\) turns divergence decrease into entropy increase. With a nonuniform stationary distribution, relative entropy can decrease without ordinary entropy increasing. The correct target is determined by the transition mechanism.

An asymmetric two-location channel shows why stationarity, rather than uniformity, is the appropriate reference in general. Take \[K=\begin{pmatrix}3/4&1/2\\1/4&1/2\end{pmatrix}.\] Its stationary distribution is \(\pi=(2/3,1/3)\), from the two-location calculation in Chapter 4. Starting at uniform gives the next state \((5/8,3/8)\). Entropy decreases because the new table is less balanced. Nevertheless, \[D((5/8,3/8)\|\pi)\leq D((1/2,1/2)\|\pi).\] The process is approaching its own stationary state rather than maximizing uncertainty over all possible states.

A model can also have several stationary distributions. On a graph with disconnected components, any compatible allocation of total mass among component equilibria remains stationary. A target must then account for those conserved component masses. Measuring divergence to an unreachable reference could misrepresent what the dynamics can accomplish.

In the connected symmetric three-node example, the exact solution identifies a unique stationary distribution and proves convergence. In a new graph, begin by checking connectivity, retained mass, and possible periodicity before generalizing that conclusion.

\section{Square-root distance also contracts}\label{ch10-square-root-distance-also-contracts}

For each output label, Cauchy--Schwarz gives \[\sqrt{\left(\sum_iK_{yi}p_i\right)
\left(\sum_iK_{yi}q_i\right)}
\geq\sum_iK_{yi}\sqrt{p_iq_i}.\] Add over outputs and use the channel column sums: \[B(Kp,Kq)\geq B(p,q).\] The output distributions have at least as much overlap. Since inverse cosine decreases as its argument increases, \[\boxed{d_{\mathrm{FR}}(Kp,Kq)\leq d_{\mathrm{FR}}(p,q).}\] The Hellinger distance also cannot increase because its square is \(1-B\).

This result links the sphere geometry to observation. A channel is not merely a linear map between numerical lists: it has a definite effect on the statistical distances we constructed.

\section{Grouping shows exactly where a distinction is lost}\label{ch10-grouping-shows-exactly-where-a-distinction-is-lost}

Take \[p=(0.1,0.4,0.1,0.4),\qquad
q=(0.4,0.1,0.4,0.1).\] The input models differ within both pairs of labels. Their divergence is \[D(p\|q)=0.6\ln4>0.\] After grouping \(\{1,2\}\) and \(\{3,4\}\), both output tables are \((1/2,1/2)\). The output divergence and spherical distance are zero.

\begin{figure}[htbp]\centering
\resizebox{.98\linewidth}{!}{\begin{tikzpicture}[every node/.style={font=\small},>=Stealth]
\node[anchor=west] at (0,4.5){Input label};\node[anchor=west] at (2,4.5){$q_i$};\node[anchor=west] at (3.1,4.5){$p_i$};
\node[anchor=west] at (7,4.5){Recorded group};
\foreach \i/\y/\q/\p in {1/3.5/.4/.1,2/2.5/.1/.4,3/1.5/.4/.1,4/.5/.1/.4}{
\node at (.5,\y){\i};\node at (2.3,\y){\q};\node at (3.4,\y){\p};}
\draw[->,teal,thick] (4,3.5)--(7,3);\draw[->,teal,thick] (4,2.5)--(7,3);
\draw[->,copper,thick] (4,1.5)--(7,1);\draw[->,copper,thick] (4,.5)--(7,1);
\node[anchor=west,align=left] at (7.2,3){Group $\{1,2\}$\\Both models: $1/2$};
\node[anchor=west,align=left] at (7.2,1){Group $\{3,4\}$\\Both models: $1/2$};
\node[anchor=west] at (0,-.5){$D(q\Vert p)=0.6\ln4\quad\longrightarrow\quad D(Kq\Vert Kp)=0$};
\end{tikzpicture}}
\caption{Grouping can erase a distinction carried entirely by the probabilities inside each group. Both recorded output tables are identical.}\end{figure}

We can account for the lost divergence exactly. For each group \(g\), let \[P_g=\sum_{i\in g}p_i,\qquad Q_g=\sum_{i\in g}q_i.\] When the group masses are positive, define conditional distributions within that group: \[p(i\mid g)=p_i/P_g,\qquad q(i\mid g)=q_i/Q_g.\] The conditional probability describes the fine label after the group is known.

Factor the ratio as \[\frac{p_i}{q_i}=\frac{P_g}{Q_g}
\frac{p(i\mid g)}{q(i\mid g)}.\] Taking logarithms, multiplying by \(p_i\), and summing first within groups gives \[\boxed{
D(p\|q)=D(P\|Q)+\sum_g P_gD(p(\cdot\mid g)\|q(\cdot\mid g)).
}\] For positive tables this is an exact identity. Terms of zero group mass are omitted under the usual conventions.

The first term is the distinction still visible in the reported groups. The remaining terms describe distinctions inside groups, which that report no longer resolves. In the four-label example, all discrepancy lies in these conditional tables.

A whole family can collapse under grouping. For \[p_\theta=(\theta/2,(1-\theta)/2,\theta/2,(1-\theta)/2),
\qquad 0\leq\theta\leq1,\] the same grouping always outputs \((1/2,1/2)\). The parameter changes the fine distribution, but this observation scheme cannot identify it. This is a concrete information-geometric question: which directions through a family of probability models remain observable?

The conditional decomposition provides an equality condition for deterministic grouping. If the two models have identical conditional distributions inside every group carrying positive first-model mass, then grouping loses no KL divergence. Their entire distinction was already visible in the group totals.

For example, \(p=(0.3,0.3,0.2,0.2)\) and \(q=(0.1,0.1,0.4,0.4)\) differ in the totals of the two pairs, but both are uniform within each pair. Grouping the pairs retains their full KL divergence. This differs from the earlier example, where the group totals agreed and every distinction was internal.

The comparison gives a concrete way to design observations. If a question concerns relative mass between groups, recording group labels may retain the relevant distinction. If it concerns how probability is divided inside a group, that same recording scheme is insufficient. A suitable observation depends on the family of models and the question, rather than on the number of recorded symbols alone.

For the parameter family \(p_\theta\), grouping \(\{1,3\}\) and \(\{2,4\}\) returns \((\theta,1-\theta)\) and retains the parameter exactly. The alternative grouping returns a constant table. Both reports have two possible outputs; the decisive difference is which fine distinctions their outputs encode.

\section{Building a complete investigation}\label{ch10-building-a-complete-investigation}

A useful final investigation begins with a small transport graph and a precise update. Choose an initial probability vector and a step size satisfying the nonnegativity condition. Record several states, their differences, and their edge currents. Check conservation directly.

Next track entropy and divergence to a stationary reference. Separate a numerical pattern from a proof: the continuity equation proves conservation, double stochasticity proves entropy monotonicity, and a particular contraction estimate or exact solution can prove convergence.

Finally pass two initial models through the same observation channel. Compare their fine and reported divergences, overlaps, and spherical distances. If labels are grouped, use the conditional decomposition to identify the distinctions that disappear.

The \href{https://bobovski66.github.io/IntroDiscreteInfoGeom/demos/week10.html}{network demonstration} explores channels and model comparisons. The transport view of the \href{https://bobovski66.github.io/IntroDiscreteInfoGeom/demos/operator-explorer.html}{operator explorer} exposes currents and divergences directly. The \href{https://bobovski66.github.io/IntroDiscreteInfoGeom/guide.html\#capstone}{course guide} suggests longer investigations.

These calculations bring the course's original question into a concrete setting: a probability distribution changes through a sequence of finite operations, and the available observations determine which changes can still be distinguished.

\section{Exercises}\label{ch10-exercises}

\begin{exercises}

\item 

Derive the three-node Laplacian from its edge weights. Verify that \(I-\tfrac14L\) is the matrix used in Chapter 4.

\item 

Find a sufficient step-size bound for nonnegative diffusion on a four-vertex star with unit edge weights. Explain the role of its central degree.

\item 

For the three-node path and \(p=(1,0,0)\), calculate currents, divergence, and one probability update at \(\tau=1/4\).

\item 

Verify that \((1,0,-1)\) and \((1,-2,1)\) are multiplied by \(3/4\) and \(1/4\) under the path matrix.

\item 

Use the exact trajectory formula to calculate \(p_2\) from \(p_0=(1,0,0)\).

\item 

Prove entropy monotonicity for a doubly stochastic matrix using concavity of \(-x\ln x\).

\item 

Give a column-stochastic channel that lowers entropy. Explain why it does not contradict Exercise 6.

\item 

Derive the log-sum inequality from Gibbs' inequality by normalizing the two lists.

\item 

Prove data processing for \(D(Kp\|Kq)\), identifying where the channel column sums enter.

\item 

For the four-label grouping example, calculate input and output overlaps and compare Fisher--Rao distances.

\item 

Derive the conditional divergence decomposition for grouping. Use it to locate the lost distinction in the worked example.

\item 

Choose two distinct parameter values in the family \(p_\theta\). Show that the fine distributions differ but the grouped observations agree. Propose a different grouping that would retain information about \(\theta\).

\end{exercises}

\section{Selected solutions}\label{ch10-selected-solutions}

\textbf{2.} The central weighted degree is three and each outer degree is one. The bound is \(0\leq\tau\leq1/3\).

\textbf{3.} Currents are \(J_{AB}=1/4\) and \(J_{BC}=0\). Divergence is \((1/4,-1/4,0)\), giving next state \((3/4,1/4,0)\).

\textbf{5.} With \(a=1/2\), \(b=1/6\), the formula gives \((5/8,5/16,1/16)\), agreeing with direct multiplication.

\textbf{7.} The channel \(\left(\begin{smallmatrix}1&1\\0&0\end{smallmatrix}\right)\) sends a fair input to \((1,0)\), lowering entropy from \(\ln2\) to zero. Its first row sums to two, so it is not doubly stochastic.

\textbf{10.} Each input overlap term is \(\sqrt{0.04}=0.2\), so \(B(p,q)=0.8\) and \(d_{\mathrm{FR}}=2\arccos0.8\). Output overlap is one and output distance is zero.

\textbf{12.} Grouping \(\{1,3\}\) and \(\{2,4\}\) yields \((\theta,1-\theta)\), which retains the parameter. The original grouping always yields \((1/2,1/2)\).
\labheading\begin{enumerate}
\def\labelenumi{\arabic{enumi}.}
\tightlist
\item
  On A--B--C with \(p=(1,0,0)\) and \(\tau=1/4\), (a) compute both signed edge currents \(J_{AB},J_{BC}\), (b) calculate outward divergence at A, B, C, (c) use \(\Delta p=-\operatorname{div}J\) to determine the next state, and (d) check mass conservation.
\item
  Using the same path, calculate \(p_1\) through \(p_4\). (a) Write each state as a probability vector. (b) Compute all coordinate differences. (c) Verify every difference vector sums to zero. (d) Explain the approach toward uniform.
\item
  For \(p_0=(1,0,0)\) and the first three transport updates, (a) compute entropy at each step, (b) calculate its first differences, (c) identify the sign of every difference, and (d) explain which property of the diffusion matrix supports the pattern.
\item
  Compare two channels: the symmetric path diffusion and a channel sending every input to A. (a) Construct each stochastic matrix. (b) Calculate both output distributions from uniform input. (c) Compare input/output entropy. (d) Explain why a general stochastic channel does not have to raise entropy.
\item
  Capstone: choose a four- or five-node undirected graph and an initial probability distribution. (a) Define a legal diffusion step using \(\tau\leq1/d_{\max}\). (b) Calculate at least four iterates with a table or short program. (c) Show currents, divergence, mass conservation, and entropy changes. (d) Compare two states by Fisher--Rao distance and explain what the mathematical results say about the flow of information.
\end{enumerate}

\FloatBarrier

\backmatter

\chapter*{Further reading}\addcontentsline{toc}{chapter}{Further reading}
The following sources provide broader context and more advanced developments. Their prerequisites vary; the present chapters supply the definitions needed for the accompanying course.
\begin{description}[style=nextline,leftmargin=1.5em]
\item[Arora, S., Hazan, E., and Kale, S. (2012).] The multiplicative weights update method: a meta-algorithm and applications. \emph{Theory of Computing} 8, 121--164. \url{https://theoryofcomputing.org/articles/v008a006/}. An extensive account of finite learning updates and their applications.
\item[Levin, D. A., and Peres, Y. (2017).] \emph{Markov Chains and Mixing Times}, second edition, with contributions by E. L. Wilmer. American Mathematical Society. Author resources: \url{https://pages.uoregon.edu/dlevin/MARKOV/}. Develops finite Markov chains, stationarity, and convergence.
\item[Nielsen, F. (2020).] An elementary introduction to information geometry. \emph{Entropy} 22(10), 1100. \url{https://doi.org/10.3390/e22101100}. A survey connecting statistical divergences with geometric structures.
\item[Shannon, C. E. (1948).] A mathematical theory of communication. \emph{Bell System Technical Journal} 27, 379--423 and 623--656. The foundational paper on entropy and communication.
\end{description}

\end{document}
